\documentclass[11pt, a4paper]{article}

\usepackage{estilo-mate}
\tema{Álgebra Lineal}

\title{\textbf{Notas de Álgebra Lineal}}
\author{Lucas Allara}
\date{}

\begin{document}

\maketitle
\newpage

\tableofcontents
\newpage

\section{Espacios Vectoriales y Subespacios}

\subsection{Espacios Vectoriales}

\begin{definicion}{$\mathbb{K}$-Espacio Vectorial}{espacio_vectorial}
Sea $\mathbb{K}$ un cuerpo. Se dice que un conjunto no vacío $V$ es un $\mathbb{K}$-espacio vectorial (o espacio vectorial sobre $\mathbb{K}$) si está provisto de dos operaciones:
\begin{enumerate}
    \item[i)] Suma: $+: V \times V \to V$
    \item[ii)] Producto por escalar: $\cdot: \mathbb{K} \times V \to V$
\end{enumerate}
tales que, para todo $u, v, w \in V$ y $\alpha, \beta \in \mathbb{K}$, se cumplen las siguientes condiciones:

Propiedades de la suma (Estructura de Grupo Abeliano):
\begin{itemize}
    \item Conmutatividad: $u+v = v+u$.
    \item Asociatividad: $(u+v)+w = u+(v+w)$.
    \item Existencia de neutro: existe $0_V \in V$ tal que $v + 0_V = v$.
    \item Existencia de opuesto: para cada $v$ existe $-v$ tal que $v + (-v) = 0_V$.
\end{itemize}

Propiedades del producto por escalar:
\begin{itemize}
    \item Asociatividad: $\alpha(\beta v) = (\alpha\beta)v$.
    \item Neutro del cuerpo: $1 \cdot v = v$ (donde $1$ es la unidad de $\mathbb{K}$).
    \item Distributividad respecto a la suma de vectores: $\alpha(u+v) = \alpha u + \alpha v$.
    \item Distributividad respecto a la suma de escalares: $(\alpha+\beta)v = \alpha v + \beta v$.
\end{itemize}
\end{definicion}

\obs Por un abuso de notación habitual, en lo sucesivo se dice que $V$ es un $\mathbb{K}$-espacio vectorial, en referencia a la terna $(V, +, \cdot)$ de la Definición \ref{def:espacio_vectorial}. Se identifica el espacio con su conjunto subyacente y se sobreentienden las operaciones de suma y producto por escalar, salvo que el contexto requiera especificarlas.

\subsection{Vectores y Operaciones}

\begin{definicion}{Vector en Sentido General}{vector_general}
Sea $V$ un $\mathbb{K}$-espacio vectorial. Se llama \textbf{vector} a cualquier elemento $v \in V$.
\end{definicion}

\obs Lo esencial de un espacio vectorial no es la naturaleza de sus elementos (polinomios, matrices, funciones), sino su comportamiento frente a la suma y el producto por escalar definidos en la Definición \ref{def:espacio_vectorial}.

Desde la perspectiva del álgebra lineal, si dos espacios comparten la misma estructura operativa, son esencialmente el mismo objeto. Distinguir una matriz de $2\times 2$ de un vector de $\mathbb{R}^4$ basándose solo en cómo se escriben sería tan artificial como sostener que una misma palabra cambia de significado según se la escriba en cursiva o en imprenta. La representación cambia, pero la entidad es la misma.

\begin{ejemplo}{Matrices}{ej_matrices}
Sea $M_{2\times 2}$ el conjunto de las matrices de $2 \times 2$. Un vector en este espacio es una matriz:
\[
A = \begin{pmatrix} a & b \\ c & d \end{pmatrix}.
\]
Aquí la suma de vectores es la suma de matrices y el producto por escalar se realiza entrada a entrada.
\end{ejemplo}

\begin{ejemplo}{El Espacio $\mathbb{R}^n$}{ej_rn_general}
Sea $\mathbb{R}^n$ el conjunto de todas las listas ordenadas de $n$ números reales.
Un vector en este espacio tiene la forma:
\[
\vec{v} = (v_1, v_2, \dots, v_n), \quad \text{donde } v_i \in \mathbb{R}.
\]
Aquí las operaciones se definen coordenada a coordenada:
\begin{itemize}
    \item Suma: si $\vec{u} = (u_1, \dots, u_n)$ y $\vec{v} = (v_1, \dots, v_n)$,
    \[
    \vec{u} + \vec{v} = (u_1 + v_1, \ u_2 + v_2, \ \dots, \ u_n + v_n).
    \]
    \item Producto por escalar: si $\alpha \in \mathbb{R}$,
    \[
    \alpha \vec{v} = (\alpha v_1, \ \alpha v_2, \ \dots, \ \alpha v_n).
    \]
\end{itemize}
El vector cero en este espacio es la lista de ceros: $\vec{0} = (0, 0, \dots, 0)$.
\end{ejemplo}

\begin{ejemplo}{Polinomios}{ej_polinomios}
Sea $P_n$ el conjunto de los polinomios de grado menor o igual a $n$ con coeficientes reales.
Un vector en este espacio es un polinomio de la forma:
\[
p(x) = a_0 + a_1 x + a_2 x^2 + \dots + a_n x^n.
\]
\end{ejemplo}

\subsection{Subespacios Vectoriales}

\begin{definicion}{Subespacio Vectorial}{subespacio}
Sea $V$ un $\mathbb{K}$-espacio vectorial. Un \textbf{subespacio vectorial} es un subconjunto $W \subseteq V$ tal que:
\begin{enumerate}
    \item $0 \in W$,
    \item Si $u,v \in W$, entonces $u+v \in W$,
    \item Si $u \in W$ y $\alpha \in \mathbb{K}$, entonces $\alpha u \in W$.
\end{enumerate}
\end{definicion}

\subsection{Combinaciones Lineales y Subespacio Generado}

\begin{definicion}{Combinación Lineal}{comb_lineal}
Sean $v_1,v_2,\dots,v_k \in V$. Una \textbf{combinación lineal} de estos vectores es un vector de la forma
\[
\alpha_1 v_1 + \alpha_2 v_2 + \cdots + \alpha_k v_k,
\quad \alpha_i \in \mathbb{K}.
\]
\end{definicion}

\begin{definicion}{Subespacio Generado}{span}
El conjunto de todas las combinaciones lineales de $v_1,v_2,\dots,v_k$ se llama \textbf{span} o \textbf{subespacio generado} y se denota
\[
\operatorname{span}\{v_1,v_2,\dots,v_k\}
= \{ \alpha_1 v_1 + \cdots + \alpha_k v_k \mid \alpha_i \in \mathbb{K} \}.
\]
Los vectores $v_1,\dots,v_k$ se llaman \textbf{generadores} de este subespacio.
\end{definicion}

\begin{proposicion}{El Span es Subespacio}{span_es_subespacio}
Sea $V$ un $\mathbb{K}$-espacio vectorial y $v_1,\dots,v_k \in V$. Entonces $\operatorname{span}\{v_1,\dots,v_k\}$ es un subespacio de $V$.
\end{proposicion}

\begin{proof}
Se verifican las tres condiciones de subespacio (Definición \ref{def:subespacio}):
\begin{enumerate}
    \item el vector nulo pertenece al span, pues con $\alpha_i = 0$ se obtiene $0 = 0v_1 + \cdots + 0v_k$;
    \item dados $u = \sum_{i=1}^k \alpha_i v_i$ y $w = \sum_{i=1}^k \beta_i v_i$ del span,
    \[
    u+w = \sum_{i=1}^k (\alpha_i+\beta_i) v_i
    \]
    es otra combinación lineal, de modo que $u+w \in \operatorname{span}\{v_1,\dots,v_k\}$;
    \item dado $u = \sum_{i=1}^k \alpha_i v_i$ del span y $\lambda \in \mathbb{K}$,
    \[
    \lambda u = \sum_{i=1}^k (\lambda \alpha_i) v_i
    \]
    también es combinación lineal, luego $\lambda u \in \operatorname{span}\{v_1,\dots,v_k\}$.
\end{enumerate}
Por lo tanto el span es un subespacio de $V$.
\end{proof}

\subsection{Representación de Subespacios}

\begin{enumerate}
    \item \textit{Forma paramétrica.} Se describe $S$ por cómo se construye (mediante el span,
    Definición \ref{def:span}): se da un conjunto de vectores y $S$ es el conjunto de todas sus
    combinaciones lineales.
    \[ S = \operatorname{span}\{u_1, \dots, u_k\} = \{ \alpha_1 u_1 + \dots + \alpha_k u_k \mid \alpha_i \in \mathbb{K} \}. \]
    Con esta forma es directo exhibir puntos del subespacio (dando valores a los $\alpha_i$), pero es difícil decidir si un punto dado pertenece a él.

    \item \textit{Forma implícita.} Se describe $S$ por las condiciones que debe cumplir: se da un
    conjunto de restricciones (ecuaciones lineales homogéneas) y $S$ es el conjunto donde todas se
    satisfacen.
    \[ S = \{ x \in V \mid Ax = 0 \}. \]
    Con esta forma es directo decidir si un punto pertenece al subespacio (verificando las restricciones), pero es difícil exhibir sus puntos.
\end{enumerate}

\subsection{Intersección y Suma de Subespacios}

\begin{definicion}{Intersección de Subespacios}{interseccion}
Sean $S$ y $T$ subespacios de un $\mathbb{K}$-espacio vectorial $V$. La intersección de $S$ y $T$ se define como:
\[
S \cap T = \{ v \in V \mid v \in S \text{ y } v \in T \}.
\]
\end{definicion}

\begin{proposicion}{Intersección de Subespacios}{prop_interseccion}
La intersección $S \cap T$ es un subespacio de $V$, está contenido en ambos y es el mayor subespacio con esta propiedad.
\end{proposicion}

\begin{proof}
\textit{$S \cap T$ es subespacio} (Definición \ref{def:subespacio}):
\begin{itemize}
    \item como $S$ y $T$ son subespacios, $0 \in S$ y $0 \in T$, luego $0 \in S \cap T$;
    \item si $u, v \in S \cap T$, entonces $u, v \in S$ y $u, v \in T$; por ser ambos subespacios, $u+v \in S$ y $u+v \in T$, de modo que $u+v \in S \cap T$;
    \item si $u \in S \cap T$ y $\alpha \in \mathbb{K}$, entonces $u \in S$ y $u \in T$, y por cerradura $\alpha u \in S$ y $\alpha u \in T$, luego $\alpha u \in S \cap T$.
\end{itemize}

\textit{Está contenido en $S$ y en $T$:} por definición de intersección, si $x \in S \cap T$ entonces $x \in S$ y $x \in T$, es decir, $S \cap T \subseteq S$ y $S \cap T \subseteq T$.

\textit{Es el mayor subespacio contenido en $S$ y $T$:} sea $W$ un subespacio con $W \subseteq S$ y $W \subseteq T$; hay que ver que $W \subseteq S \cap T$. Dado $w \in W$, como $W \subseteq S$ se tiene $w \in S$, y como $W \subseteq T$, $w \in T$; al estar en ambos, $w \in S \cap T$. Luego $W \subseteq S \cap T$.
\end{proof}

\begin{definicion}{Suma de Subespacios}{suma_subespacios}
Sean $S$ y $T$ subespacios de un $\mathbb{K}$-espacio vectorial $V$. La suma de $S$ y $T$ se define como el menor subespacio de $V$ que contiene a $S$ y $T$:
\[
S + T = \bigcap \{ W \subseteq V \mid W \text{ es subespacio de $V$ y } S \subseteq W, \ T \subseteq W \}.
\]
\end{definicion}

\begin{proposicion}{Equivalencia de Suma}{suma_equiv}
La siguiente definición:
\[
S+T = \{ s+t \mid s \in S, t \in T \}
\]
es equivalente a la definición anterior.
\end{proposicion}

\begin{proof}
Sea $W = \{ s+t \mid s \in S, t \in T \}$.

\textit{$W$ es subespacio} (Definición \ref{def:subespacio}):
\begin{itemize}
    \item con $s=0 \in S$ y $t=0 \in T$, $0 = 0+0 \in W$;
    \item si $w_1 = s_1+t_1$ y $w_2 = s_2+t_2$ están en $W$ (con $s_i \in S$, $t_i \in T$), entonces
    \[
    w_1 + w_2 = (s_1+s_2) + (t_1+t_2) \in W,
    \]
    pues $s_1+s_2 \in S$ y $t_1+t_2 \in T$;
    \item para $\alpha \in \mathbb{K}$, $\alpha w = (\alpha s) + (\alpha t) \in W$, porque $\alpha s \in S$ y $\alpha t \in T$.
\end{itemize}

\textit{Contiene a $S$ y $T$:} para $s \in S$, $s+0 \in W$, de modo que $S \subseteq W$; para $t \in T$, $0+t \in W$, de modo que $T \subseteq W$.

\textit{Es el menor que contiene a $S$ y $T$:} sea $W'$ un subespacio con $S \subseteq W'$ y $T \subseteq W'$. Para $w = s+t \in W$ con $s \in S$, $t \in T$, se tiene $s, t \in W'$, y como $W'$ es subespacio, $s+t \in W'$. Luego $W \subseteq W'$.
\end{proof}

\subsection{Suma Directa y Escritura Única}

\begin{definicion}{Unicidad de la Escritura}{unicidad_general}
Sea $v$ un vector. La escritura de $v$ es \textbf{única} si dos representaciones cualesquiera de $v$ coinciden componente a componente.
\end{definicion}

\begin{ejemplo}{Aplicación a Combinaciones Lineales}{}
Que $v$ se escriba de forma única como combinación lineal de $\{u_1, \dots, u_n\}$ significa que si
\[
v = \sum \alpha_i u_i \quad \text{y} \quad v = \sum \beta_i u_i,
\]
entonces obligatoriamente $\alpha_i = \beta_i$ para todo $i$.
\end{ejemplo}

\begin{ejemplo}{Aplicación a Suma de Subespacios}{unicidad_suma}
Que un vector $v \in S+T$ se escriba de forma única como suma de un elemento de $S$ y otro de $T$ (Definición \ref{def:unicidad_general}) significa que si se plantean dos descomposiciones:
\[
v = s + t \quad \text{y} \quad v = s' + t',
\]
con $s, s' \in S$ y $t, t' \in T$, la condición de unicidad implica necesariamente que:
\[
s = s' \quad \text{y} \quad t = t'.
\]
\end{ejemplo}

\begin{definicion}{Suma Directa}{suma_directa}
Si todo vector $v \in S+T$ tiene escritura única, $S+T$ es la \textbf{suma directa} de $S$ y $T$, y se escribe
\[
S \oplus T = S+T.
\]
\end{definicion}

\begin{proposicion}{Caracterización de Suma Directa}{suma_directa_carac}
Sean $S$ y $T$ subespacios de $V$. La suma $S+T$ es directa si y solo si
\[
S \cap T = \{0\}.
\]
\end{proposicion}

\begin{proof}
Se prueban las dos implicaciones.

\textit{($\Rightarrow$).} Supóngase $S \oplus T$ y sea $v \in S \cap T$. Entonces $v \in S$ y $v \in T$, y
\[
v = v + 0 = 0 + v,
\]
con $v, 0 \in S$ y $v, 0 \in T$. Como la escritura es única (Definición \ref{def:suma_directa}), $v = 0$. Luego $S \cap T = \{0\}$.

\textit{($\Leftarrow$).} Supóngase $S \cap T = \{0\}$ y sea $v \in S+T$ con dos escrituras
\[
v = s_1 + t_1 = s_2 + t_2, \quad s_i \in S, \ t_i \in T.
\]
Restando, $0 = (s_1 - s_2) + (t_1 - t_2)$, con $s_1 - s_2 \in S$ y $t_1 - t_2 \in T$. Como $S \cap T = \{0\}$,
\[
s_1 - s_2 = 0, \quad t_1 - t_2 = 0,
\]
es decir, $s_1 = s_2$ y $t_1 = t_2$. Por lo tanto la escritura de $v$ es única y $S+T$ es suma directa.
\end{proof}

\subsection{Producto Interno}

\begin{definicion}{Producto Interno}{producto_interno_general}
Sea $V$ un $\mathbb{K}$-espacio vectorial (donde $\mathbb{K}$ suele ser $\mathbb{R}$ o $\mathbb{C}$). Un \textbf{producto interno} en $V$ es una función
\[
\langle \cdot , \cdot \rangle : V \times V \to \mathbb{K}
\]
que satisface las siguientes propiedades para todo $x, y, z \in V$ y $\alpha \in \mathbb{K}$:

\begin{enumerate}
    \item Simetría Conjugada (o Hermítica):
    \[
    \langle x, y \rangle = \overline{\langle y, x \rangle}.
    \]
    \item Linealidad en la primera variable:
    \[
    \langle x + y, z \rangle = \langle x, z \rangle + \langle y, z \rangle,
    \quad
    \langle \alpha x, y \rangle = \alpha \langle x, y \rangle.
    \]
    \item Definido positivo:
    \[
    \langle x, x \rangle \in \mathbb{R}_{\geq 0} \quad (\text{es un número real no negativo}),
    \]
    y además $\langle x, x \rangle = 0 \iff x = 0$.
\end{enumerate}
\end{definicion}

\obs Si $\mathbb{K} = \mathbb{C}$, la propiedad 2 implica que en la segunda variable salen los escalares conjugados:
\[
\langle x, \alpha y \rangle = \overline{\langle \alpha y, x \rangle} = \overline{\alpha \langle y, x \rangle} = \overline{\alpha} \overline{\langle y, x \rangle} = \overline{\alpha} \langle x, y \rangle.
\]
Si $\mathbb{K} = \mathbb{R}$, el conjugado es la identidad ($\bar{a} = a$), y se recupera la simetría y bilinealidad.

\obs El producto interno usual en $\mathbb{R}^n$ puede verse como una multiplicación de matrices: representando $\vec{x}$ como vector fila e $\vec{y}$ como vector columna,
\[
\langle \vec{x}, \vec{y} \rangle = 
\begin{bmatrix} x_1 & x_2 & \cdots & x_n \end{bmatrix}
\begin{bmatrix} y_1 \\ y_2 \\ \vdots \\ y_n \end{bmatrix}.
\]

\subsection{Norma y Cauchy-Schwarz}

\begin{definicion}{Norma Inducida}{norma}
Sea $(V,\langle\cdot,\cdot\rangle)$ un espacio con producto interno. La \textbf{norma inducida} de un vector $x \in V$ es:
\[
\|x\| := \sqrt{\langle x,x\rangle}.
\]
\end{definicion}

\begin{proposicion}{Desigualdad de Cauchy-Schwarz}{cauchy_schwarz}
En todo espacio con producto interno $(V,\langle\cdot,\cdot\rangle)$ se cumple
\[
\;|\langle x,y\rangle| \le \|x\|\,\|y\|\; \qquad\forall x,y\in V.
\]
\end{proposicion}

\begin{proof}
Se prueba el caso $\mathbb{K} = \mathbb{R}$; el caso complejo se reduce a él en la observación
que sigue. La idea es usar que la cuadrática $\varphi(t) = \langle tx+y, tx+y\rangle$ es no negativa. Sean
$x, y \in V$; para todo $t \in \mathbb{R}$ se tiene $\varphi(t) \ge 0$, y por bilinealidad
\[
\varphi(t) = t^2\langle x,x\rangle + 2t\langle x,y\rangle + \langle y,y\rangle,
\]
una cuadrática $a t^2 + b t + c$ con $a = \|x\|^2$, $b = 2\langle x,y\rangle$, $c = \|y\|^2$. Como
$\varphi(t)\ge 0$, su discriminante satisface
\[
\Delta = b^2 - 4ac = 4\Big(\langle x,y\rangle^2 - \|x\|^2\|y\|^2\Big) \le 0,
\]
de donde $\langle x,y\rangle^2 \le \|x\|^2\|y\|^2$ y, tomando raíces, $|\langle x,y\rangle| \le \|x\|\,\|y\|$.
\end{proof}

\obs
Si el cuerpo es $\mathbb{C}$, no puede formarse el polinomio directamente con $\|x+ty\|^2$ porque el término central sería complejo y $\varphi$ no sería una cuadrática real en $t$. El problema se transforma al caso real como sigue.

Supóngase $\langle x, y \rangle \neq 0$ (si es $0$, la desigualdad vale directamente). Este número complejo se escribe en forma polar:
\[
\langle x, y \rangle = |\langle x, y \rangle| e^{i\theta}.
\]
Se define un vector auxiliar $\tilde{y}$ rotando $y$ para cancelar esa fase:
\[
\tilde{y} = e^{i\theta} y.
\]
Sobre $\tilde{y}$ se observan dos cosas:
\begin{enumerate}
    \item Tiene la misma norma que $y$: $\|\tilde{y}\| = |e^{i\theta}| \|y\| = 1 \cdot \|y\| = \|y\|$.
    \item Su producto interno con $x$ es ahora un número real:
    \[
    \langle x, \tilde{y} \rangle = \langle x, e^{i\theta} y \rangle = \overline{e^{i\theta}} \langle x, y \rangle = e^{-i\theta} \left( |\langle x, y \rangle| e^{i\theta} \right) = |\langle x, y \rangle|.
    \]
\end{enumerate}

Ahora se plantea el polinomio cuadrático con variable real $t \in \mathbb{R}$:
\[
\varphi(t) = \| x + t\tilde{y} \|^2 \ge 0.
\]
Desarrollando, usando que $\langle x, \tilde{y} \rangle$ es real,
\begin{align*}
\| x + t\tilde{y} \|^2 &= \langle x + t\tilde{y}, x + t\tilde{y} \rangle \\
&= \langle x, x \rangle + t \langle x, \tilde{y} \rangle + t \langle \tilde{y}, x \rangle + t^2 \langle \tilde{y}, \tilde{y} \rangle \\
&= \|x\|^2 + t |\langle x, y \rangle| + t \overline{|\langle x, y \rangle|} + t^2 \|\tilde{y}\|^2 \\
&= \|x\|^2 + 2t |\langle x, y \rangle| + t^2 \|y\|^2.
\end{align*}
Desde aquí la demostración continúa como en el caso real, y como $\langle x, \tilde{y} \rangle = |\langle x, y \rangle|$ y $\|\tilde{y}\| = \|y\|$, la conclusión es la desigualdad buscada.

\begin{proposicion}{Propiedades de la Norma Inducida}{prop_norma}
La norma inducida cumple las siguientes propiedades para todo $x, y \in V$ y $\alpha \in \mathbb{K}$:
\begin{enumerate}
    \item Positividad: $\|x\| \ge 0$, y $\|x\|=0 \iff x=0$.
    \item Homogeneidad: $\|\alpha x\| = |\alpha|\,\|x\|$.
    \item Desigualdad Triangular: $\|x+y\| \le \|x\| + \|y\|$.
\end{enumerate}
Es decir, la norma inducida es una norma.
\end{proposicion}

\begin{proof}
Las dos primeras se verifican a partir de la definición de producto interno (Definición \ref{def:producto_interno_general}); la última se deduce de la desigualdad de Cauchy-Schwarz (Proposición \ref{prop:cauchy_schwarz}).
\end{proof}

\begin{proposicion}{Condición de Igualdad en Cauchy-Schwarz}{igualdad_cs}
Sean $x, y$ vectores en un espacio con producto interno. La igualdad
\[
|\langle x, y \rangle| = \|x\| \|y\|
\]
se cumple si y solo si uno es múltiplo escalar del otro.
\end{proposicion}

\begin{proof}
Se prueban las dos implicaciones.

\textit{($\Rightarrow$).} Supóngase que $|\langle x, y \rangle| = \|x\| \|y\|$. Si $x = 0$, entonces
$x$ es múltiplo de $y$ y no hay nada que probar; supóngase $x \neq 0$ y considérese el vector $z$:
\[
z = y - \frac{\langle y, x \rangle}{\|x\|^2} x.
\]
Su norma al cuadrado, expandiendo $\langle z, z \rangle$, es:
\[
\|z\|^2 = \left\langle y - \frac{\langle y, x \rangle}{\|x\|^2} x, \ y - \frac{\langle y, x \rangle}{\|x\|^2} x \right\rangle.
\]
Al desarrollar (en el caso complejo, $\langle \alpha u, v \rangle = \alpha \langle u, v \rangle$ y $\langle u, \alpha v \rangle = \bar{\alpha} \langle u, v \rangle$):
\[
\|z\|^2 = \|y\|^2 - \frac{\overline{\langle y, x \rangle}}{\|x\|^2}\langle y, x \rangle - \frac{\langle y, x \rangle}{\|x\|^2}\langle x, y \rangle + \left| \frac{\langle y, x \rangle}{\|x\|^2} \right|^2 \|x\|^2.
\]
Como $\overline{\langle y, x \rangle} = \langle x, y \rangle$, los términos centrales se agrupan:
\[
\|z\|^2 = \|y\|^2 - \frac{|\langle x, y \rangle|^2}{\|x\|^2} - \frac{|\langle x, y \rangle|^2}{\|x\|^2} + \frac{|\langle x, y \rangle|^2}{\|x\|^4} \|x\|^2.
\]
Simplificando:
\[
\|z\|^2 = \|y\|^2 - 2\frac{|\langle x, y \rangle|^2}{\|x\|^2} + \frac{|\langle x, y \rangle|^2}{\|x\|^2} = \|y\|^2 - \frac{|\langle x, y \rangle|^2}{\|x\|^2}.
\]
Por hipótesis $|\langle x, y \rangle|^2 = \|x\|^2 \|y\|^2$, de donde:
\[
\|z\|^2 = \|y\|^2 - \frac{\|x\|^2 \|y\|^2}{\|x\|^2} = \|y\|^2 - \|y\|^2 = 0.
\]
Como $\|z\|^2 = 0$ implica $z = 0$ (positividad de la norma, Proposición \ref{prop:prop_norma}), se obtiene:
\[
y - \frac{\langle y, x \rangle}{\|x\|^2} x = 0 \implies y = \left( \frac{\langle y, x \rangle}{\|x\|^2} \right) x.
\]
Por lo tanto $y$ es un múltiplo escalar de $x$.

\textit{($\Leftarrow$).} Existe $\lambda \in \mathbb{K}$ con $y = \lambda x$. Sustituyendo en ambos lados:
\begin{itemize}
    \item Lado izquierdo:
    \[
    |\langle x, y \rangle| = |\langle x, \lambda x \rangle| = |\bar{\lambda} \langle x, x \rangle| = |\lambda| \|x\|^2.
    \]
    \item Lado derecho:
    \[
    \|x\| \|y\| = \|x\| \|\lambda x\| = \|x\| (|\lambda| \|x\|) = |\lambda| \|x\|^2.
    \]
\end{itemize}
Ambos lados coinciden.
\end{proof}

\subsection{Distancia y Ángulos}

La existencia de un producto interno dota al espacio de una estructura geométrica. A partir de él se inducen conceptos métricos como la longitud de los vectores (norma) y la distancia entre ellos.

Hay una jerarquía: un producto interno induce una norma, y una norma induce una distancia; el camino inverso no siempre vale (existen distancias que no provienen de normas, y normas que no provienen de productos internos).

\begin{definicion}{Distancia Inducida}{distancia}
Sea $(V,\|.\|)$ un espacio normado. La \textbf{distancia inducida} entre dos vectores $x, y \in V$ es:
\[
d(x,y) := \|x-y\|.
\]
\end{definicion}

\begin{proposicion}{Propiedades de la Distancia Inducida}{prop_distancia}
La función $d$ define una métrica en $V$, es decir:
\begin{enumerate}
    \item $d(x,y) \ge 0$ y $d(x,y)=0 \iff x=y$.
    \item Simetría: $d(x,y) = d(y,x)$.
    \item Desigualdad Triangular: $d(x,z) \le d(x,y) + d(y,z)$.
\end{enumerate}
\end{proposicion}

\begin{proof}
Las tres propiedades se verifican a partir de las de la norma (Proposición \ref{prop:prop_norma}).
\end{proof}

\begin{definicion}{Ángulo entre Vectores}{angulo}
Sean $x, y \in V$ vectores no nulos. El \textbf{ángulo} $\theta$ entre ellos es el único número real tal que:
\[
\cos\theta = \frac{\langle x,y\rangle}{\|x\|\cdot\|y\|}, \qquad 0 \le \theta \le \pi.
\]
\end{definicion}

\obs Esta definición tiene sentido gracias a la desigualdad de Cauchy-Schwarz (Proposición \ref{prop:cauchy_schwarz}).
Se tiene $|\langle x, y \rangle| \le \|x\| \|y\|$; dividiendo por el producto de las normas,
\[
\frac{|\langle x, y \rangle|}{\|x\| \|y\|} \le 1 \implies -1 \le \frac{\langle x, y \rangle}{\|x\| \|y\|} \le 1.
\]
Como el valor siempre cae en el intervalo $[-1, 1]$, está dentro del dominio de la función arco coseno, lo que garantiza que el ángulo $\theta$ esté bien definido.

\newpage

\section{Espacios Ortogonales, Independencia Lineal, Base y Dimensión}

\subsection{Ortogonalidad y Complemento Ortogonal}

\begin{definicion}{Ortogonalidad}{ortogonalidad}
Un vector $x$ es \textbf{ortogonal} a $y$ ($x \perp y$) si $\langle x, y \rangle = 0$.
\end{definicion}

\begin{definicion}{Subespacios Ortogonales}{sub_ortogonales}
Sea $V$ un $\mathbb{K}$-espacio vectorial con producto interno y sean $S$ y $T$ subespacios de $V$.
$S$ y $T$ son \textbf{ortogonales} ($S \perp T$) si:
\[
\langle x, y \rangle = 0, \quad \forall x \in S, \ \forall y \in T.
\]
\end{definicion}

\begin{proposicion}{Intersección Trivial}{int_trivial}
Si $S \perp T$, entonces $S \cap T = \{0\}$.
\end{proposicion}

\begin{proof}
Sea $z \in S \cap T$. Como $z \in S$, $z \in T$ y $S \perp T$, se tiene $\langle z, z \rangle = 0$, luego $z = 0$.
\end{proof}

\begin{definicion}{Complemento Ortogonal}{complemento_ortogonal}
Sea $V$ un $\mathbb{K}$-espacio vectorial con producto interno y sea $S$ un subespacio de $V$. El \textbf{complemento ortogonal} de $S$, denotado $S^\perp$, es:
\[
S^\perp = \{ x \in V \mid \langle x, s \rangle = 0, \ \forall s \in S \}.
\]
\end{definicion}

\begin{proposicion}{Propiedades del Complemento Ortogonal}{prop_complemento}
Sea $V$ un $\mathbb{K}$-espacio vectorial con producto interno y sea $S$ un subespacio de $V$. El complemento ortogonal $S^\perp$ es un subespacio de $V$, es ortogonal a $S$ y es el mayor subespacio que cumple esta propiedad.
\end{proposicion}

\begin{proof}
\textit{$S^\perp$ es subespacio:}
\begin{itemize}
    \item el vector nulo pertenece a $S^\perp$, pues $\langle 0, s \rangle = 0$ para todo $s \in S$;
    \item si $x, y \in S^\perp$, para todo $s \in S$ se tiene $\langle x+y, s \rangle = \langle x, s \rangle + \langle y, s \rangle = 0$, luego $x+y \in S^\perp$;
    \item si $x \in S^\perp$ y $\alpha \in \mathbb{K}$, $\langle \alpha x, s \rangle = \alpha \langle x, s \rangle = 0$, luego $\alpha x \in S^\perp$.
\end{itemize}

\textit{$S^\perp \perp S$:} por definición de complemento ortogonal (Definición \ref{def:complemento_ortogonal}), $\langle x, s \rangle = 0$ para todo $x \in S^\perp$ y todo $s \in S$.

\textit{Es el mayor con esa propiedad:} sea $W$ un subespacio con $W \perp S$. Dado $w \in W$, como $W \perp S$ se tiene $\langle w, s \rangle = 0$ para todo $s \in S$, de modo que $w \in S^\perp$. Luego $W \subseteq S^\perp$.
\end{proof}

\subsection{Proyección sobre Subespacios y Complementariedad}

\begin{definicion}{Dimensión Finita}{dim_finita}
Un $\mathbb{K}$-espacio vectorial $V$ se dice \textbf{de dimensión finita} (o finitamente generado) si posee un conjunto generador finito, es decir, si existe un subconjunto finito $G \subset V$ tal que:
\[ \operatorname{span}(G) = V \]
En caso contrario, se dice \textbf{de dimensión infinita}.
\end{definicion}

El siguiente teorema es un resultado clásico del análisis y se enuncia sin demostración.

\begin{teorema}{Mejor Aproximación en Subespacios}{mejor_aprox}
Sea $V$ un $\mathbb{K}$-espacio vectorial de dimensión finita con producto interno y sea $S$ un subespacio de $V$.
Para todo vector $v \in V$, existe un único elemento $p \in S$ tal que:
\[ \|v - p\| \leq \|v - s\| \quad \text{para todo } s \in S. \]
Es decir, $p$ es el vector de $S$ más cercano a $v$.
\end{teorema}

\begin{definicion}{Proyección Ortogonal}{proyeccion_ortogonal}
Al vector $p \in S$ del Teorema \ref{teo:mejor_aprox} se lo llama \textbf{proyección ortogonal} de $v$ sobre $S$ y se denota $P_S(v)$.
\end{definicion}

\begin{teorema}{Descomposición Ortogonal}{descomposicion_ortogonal}
Sea $V$ un $\mathbb{K}$-espacio vectorial de dimensión finita con producto interno y $S$ un subespacio de $V$. Entonces:
\[
V = S \oplus S^\perp
\]
Esto significa que todo vector $v \in V$ se puede escribir de manera única como la suma de dos componentes ortogonales:
\[ v = p + q, \quad \text{donde } p \in S \text{ y } q \in S^\perp. \]
Específicamente, $p = P_S(v)$ y $q = v - P_S(v)$.
\end{teorema}

\begin{proof}
Sea $p \in S$ el elemento que minimiza la distancia a $v$ (Teorema \ref{teo:mejor_aprox}) y $q = v - p$. Hay que probar que $q \in S^\perp$, es decir, $\langle q, y \rangle = 0$ para todo $y \in S$.

Fíjese $y \in S$. Para todo $\alpha \in \mathbb{R}$, el vector $p + \alpha y$ pertenece a $S$, y la función
\[
f(\alpha) = \|v - (p + \alpha y)\|^2 = \|q - \alpha y\|^2
\]
tiene, por la minimalidad de $p$ (Teorema \ref{teo:mejor_aprox}), un mínimo absoluto en $\alpha = 0$. Desarrollando,
\[
f(\alpha) = \langle q - \alpha y, q - \alpha y \rangle = \|q\|^2 - 2\alpha \operatorname{Re}\langle q, y \rangle + \alpha^2 \|y\|^2.
\]
Como $f$ es derivable y tiene un mínimo en $0$, $f'(0) = -2 \operatorname{Re}\langle q, y \rangle = 0$, de modo que $\operatorname{Re}\langle q, y \rangle = 0$. En el caso real esto da $\langle q, y \rangle = 0$; en el complejo, el mismo argumento aplicado a $iy$ (Observación siguiente) da $\langle q, y \rangle = 0$. Luego $q \perp S$ y $v = p + q$.

Finalmente, como $S$ y $S^\perp$ son ortogonales, $S \cap S^\perp = \{0\}$ (Proposición \ref{prop:int_trivial}), de modo que la suma es directa y la escritura es única.
\end{proof}

\obs En la demostración anterior, con $\alpha \in \mathbb{R}$ se prueba que la parte real del producto interno es nula. Para la parte imaginaria, en un espacio complejo:
\begin{itemize}
    \item como $S$ es subespacio, $y \in S \implies iy \in S$;
    \item aplicando el resultado anterior a $iy$, $\operatorname{Re}\langle q, iy \rangle = 0$;
    \item como los escalares en la segunda variable salen conjugados, $\langle q, iy \rangle = -i \langle q, y \rangle$;
    \item entonces $0 = \operatorname{Re}(-i\langle q, y \rangle) = \operatorname{Im}(\langle q, y \rangle)$.
\end{itemize}
Como la parte real y la imaginaria son cero, el producto interno es cero.

\begin{corolario}{Doble Ortogonal}{doble_orto}
$(S^\perp)^\perp = S$.
\end{corolario}

\begin{proof}
Se prueba la doble inclusión.

\textit{($\supseteq$).} Sea $s \in S$. Para todo $z \in S^\perp$, $\langle s, z \rangle = 0$, luego $s \in (S^\perp)^\perp$.

\textit{($\subseteq$).} Sea $x \in (S^\perp)^\perp$ y escríbase $x = s + z$ con $s \in S$, $z \in S^\perp$. Entonces $\langle x, z \rangle = \langle s, z \rangle + \langle z, z \rangle$; como $x \in (S^\perp)^\perp$ y $z \in S^\perp$, $\langle x, z \rangle = 0$, y como $s \in S$, $\langle s, z \rangle = 0$. Queda $0 = \|z\|^2$, de modo que $z = 0$ y $x = s \in S$.
\end{proof}

\begin{proposicion}{Relaciones de Inclusión y Operaciones}{relaciones_orto}
Sea $V$ un $\mathbb{K}$-espacio vectorial de dimensión finita con producto interno y sean $S$ y $T$ subespacios de $V$. Entonces:
\begin{enumerate}
    \item $S \subset T \iff T^\perp \subset S^\perp$.
    \item $(S+T)^\perp = S^\perp \cap T^\perp$.
    \item $(S \cap T)^\perp = S^\perp + T^\perp$.
\end{enumerate}
\end{proposicion}

\begin{proof}
\textit{(1)} Se prueban las dos implicaciones.

\textit{($\Rightarrow$).} Sea $x \in T^\perp$, es decir, $\langle x, t \rangle = 0$ para todo $t \in T$. Como $S \subset T$, en particular $\langle x, s \rangle = 0$ para todo $s \in S$, de modo que $x \in S^\perp$. Luego $T^\perp \subset S^\perp$.

\textit{($\Leftarrow$).} Aplicando la ida a $T^\perp \subset S^\perp$ se obtiene $(S^\perp)^\perp \subseteq (T^\perp)^\perp$, y por el doble complemento ortogonal (Corolario \ref{cor:doble_orto}), $S \subset T$.

\textit{(2)} Doble inclusión. \textit{($\subseteq$)} si $x \in (S+T)^\perp$, entonces $\langle x, s+t \rangle = 0$; con $t=0$, $\langle x, s \rangle = 0$, de modo que $x \in S^\perp$, y con $s=0$, $x \in T^\perp$, luego $x \in S^\perp \cap T^\perp$. \textit{($\supseteq$)} si $x \in S^\perp \cap T^\perp$, entonces $\langle x, s \rangle = 0$ y $\langle x, t \rangle = 0$, de modo que $\langle x, s+t \rangle = 0$ y $x \in (S+T)^\perp$.

\textit{(3)} Por (2), $(S^\perp + T^\perp)^\perp = S^{\perp\perp} \cap T^{\perp\perp} = S \cap T$; tomando ortogonal a ambos lados, $S^\perp + T^\perp = (S \cap T)^\perp$.
\end{proof}

\subsection{Independencia Lineal}

\begin{proposicion}{Equivalencias de Independencia Lineal}{equiv_li}
Sea $V$ un $\mathbb{K}$-espacio vectorial y sean $v_1, v_2, \dots, v_n$ vectores en $V$. Las siguientes afirmaciones son equivalentes:
\begin{enumerate}[label=(\roman*)] % etiquetas romanas para que coincidan con las referencias (i)/(ii)/(iii) de la demostración
    \item Cada vector $x \in \operatorname{span}\{v_1, \dots, v_n\}$ se escribe de manera única como combinación lineal de los $v_i$.
    \item Ningún $v_j$ es combinación lineal de los demás vectores del conjunto.
    \item La única combinación lineal nula es la trivial:
    \[
    \sum_{i=1}^{n} \alpha_i v_i = 0 \implies \alpha_1 = \alpha_2 = \dots = \alpha_n = 0.
    \]
\end{enumerate}
\end{proposicion}

\begin{proof}
Se prueba en cadena (i) $\implies$ (ii) $\implies$ (iii) $\implies$ (i).

\textit{(i) $\implies$ (ii).} Por el absurdo: si algún $v_j$ fuera combinación lineal de los demás, $v_j = \sum_{i \neq j} \beta_i v_i$, habría dos escrituras distintas de $v_j$ (esa y $1 \cdot v_j$), contra la unicidad de (i).

\textit{(ii) $\implies$ (iii).} Por contrarrecíproco: si hubiera una combinación lineal nula $\sum_{i} \alpha_i v_i = 0$ con algún $\alpha_j \neq 0$, despejando,
\[
v_j = \sum_{i \neq j} \left( -\frac{\alpha_i}{\alpha_j} \right) v_i,
\]
de modo que $v_j$ sería combinación de los demás, contra (ii).

\textit{(iii) $\implies$ (i).} Sea $x$ en el generado con dos escrituras $x = \sum_{i} \alpha_i v_i = \sum_{i} \beta_i v_i$. Restando,
\[
0 = \sum_{i=1}^{n} (\alpha_i - \beta_i) v_i,
\]
y por (iii) todos los coeficientes son nulos, $\alpha_i = \beta_i$ para todo $i$. Luego la escritura es única.
\end{proof}

\begin{definicion}{Conjunto Linealmente Independiente}{def_li}
Sea $V$ un $\mathbb{K}$-espacio vectorial y sean $v_1, v_2, \dots, v_n$ vectores en $V$.
\begin{itemize}
    \item El conjunto es \textbf{linealmente independiente} si cumple cualquiera de las condiciones equivalentes de la Proposición \ref{prop:equiv_li}.
    \item En caso contrario, es \textbf{linealmente dependiente}.
\end{itemize}
\end{definicion}

\subsection{Base y Dimensión}

\begin{definicion}{Base de un Espacio (Hamel)}{def_base}
Sea $V$ un $\mathbb{K}$-espacio vectorial. Un conjunto de vectores $\mathcal{B} \subset V$ es una \textbf{base} de $V$ si cumple dos condiciones:
\begin{enumerate}
    \item Todo vector $v \in V$ se puede expresar como una combinación lineal finita de elementos de $\mathcal{B}$.
    Es decir, existen $b_1, \dots, b_k \in \mathcal{B}$ y escalares $\alpha_1, \dots, \alpha_k$ tales que:
    \[ v = \sum_{i=1}^k \alpha_i b_i. \]
    \item Es linealmente independiente.
\end{enumerate}
\end{definicion}

\obs La condición de finitud en la suma es esencial: no se define la suma infinita de vectores, de modo que, aunque una base pueda tener infinitos elementos, cada vector se construye con una cantidad finita de ellos.

\begin{proposicion}{Reducción a una Base}{reduccion_base}
Sea $A = \{v_1, \dots, v_k\}$ un conjunto de vectores que genera un $\mathbb{K}$-espacio vectorial $V$. Si $A$ es linealmente dependiente, entonces existe un subconjunto propio de $A$ que todavía genera a $V$.
\end{proposicion}

\begin{proof}
Como $A$ es linealmente dependiente, por la Proposición \ref{prop:equiv_li} algún $v_j \in A$ es combinación lineal de los demás:
\[
v_j = \sum_{i \neq j} \alpha_i v_i.
\]
Hay que ver que $A' = A \setminus \{v_j\}$ sigue generando $V$. La inclusión $\operatorname{span}(A') \subseteq \operatorname{span}(A)$ vale porque $A' \subset A$. Para la otra, sea $x \in \operatorname{span}(A) = V$,
\[
x = \beta_j v_j + \sum_{i \neq j} \beta_i v_i;
\]
sustituyendo $v_j$,
\[
x = \beta_j \left( \sum_{i \neq j} \alpha_i v_i \right) + \sum_{i \neq j} \beta_i v_i,
\]
combinación lineal de los $v_i$ con $i \neq j$, de modo que $x \in \operatorname{span}(A')$. Luego se puede quitar $v_j$ sin alterar el generado.
\end{proof}

\begin{proposicion}{Extensión a una Base}{extension_base}
Sea $A = \{v_1, \dots, v_k\}$ un conjunto linealmente independiente contenido en un $\mathbb{K}$-espacio vectorial $V$. Si $A$ no genera a $V$, entonces se puede agregar un vector $v_{k+1} \in V$ tal que el nuevo conjunto $\{v_1, \dots, v_k, v_{k+1}\}$ sigue siendo linealmente independiente.
\end{proposicion}

\begin{proof}
Como $\operatorname{span}(A) \neq V$, existe $v_{k+1} \in V$ con $v_{k+1} \notin \operatorname{span}\{v_1, \dots, v_k\}$; hay que ver que $\{v_1, \dots, v_k, v_{k+1}\}$ es linealmente independiente.

Por el absurdo, supóngase que es linealmente dependiente: existen escalares $\alpha_1, \dots, \alpha_{k+1}$ no todos nulos con
\[
\alpha_1 v_1 + \dots + \alpha_k v_k + \alpha_{k+1} v_{k+1} = 0.
\]
\begin{itemize}
    \item Si $\alpha_{k+1} = 0$, la ecuación se reduce a $\alpha_1 v_1 + \dots + \alpha_k v_k = 0$; como $\{v_1, \dots, v_k\}$ es linealmente independiente, todos los $\alpha_i$ son $0$, contra que no todos sean nulos.
    \item Si $\alpha_{k+1} \neq 0$, despejando,
    \[
    v_{k+1} = \sum_{i=1}^k \left( -\frac{\alpha_i}{\alpha_{k+1}} \right) v_i \in \operatorname{span}\{v_1, \dots, v_k\},
    \]
    contra la elección de $v_{k+1}$.
\end{itemize}
En ambos casos hay contradicción, de modo que el conjunto es linealmente independiente.
\end{proof}

\begin{teorema}{Existencia de Base (Caso Finito)}{existencia_base}
Todo $\mathbb{K}$-espacio vectorial $V \neq \{0\}$ de dimensión finita tiene una base.
\end{teorema}

\begin{proof}
Por la definición de dimensión finita (Definición \ref{def:dim_finita}), $V$ admite un generador finito $G = \{v_1, \dots, v_k\}$. Se reduce $G$ hasta una base:
\begin{itemize}
    \item si $G$ es linealmente independiente, como además genera $V$, ya es una base;
    \item si $G$ es linealmente dependiente, por la Proposición \ref{prop:reduccion_base} se puede quitar un vector obteniendo $G'$ con $\operatorname{span}(G') = V$ y un elemento menos.
\end{itemize}
Como $G$ es finito, el proceso de eliminación se detiene en a lo sumo $k-1$ pasos. El conjunto final $\mathcal{B}$ genera $V$ y es linealmente independiente, luego es una base de $V$.
\end{proof}

\obs El subespacio nulo $S = \{0\}$ es el único subespacio que no tiene una base.

\begin{teorema}{Relación entre Generadores e Independencia Lineal}{teorema_intercambio}
Sea $V$ un $\mathbb{K}$-espacio vectorial generado por un conjunto de vectores $\{u_1, \dots, u_r\}$ (es decir, $V = \operatorname{span}\{u_1, \dots, u_r\}$).
Si $\{v_1, \dots, v_k\}$ es un conjunto de vectores linealmente independientes contenidos en $V$, entonces:
\[
k \le r.
\]
\end{teorema}

\begin{proof}
Por el absurdo, supóngase $k > r$. Como $\{u_1, \dots, u_r\}$ genera $V$ y cada $v_j \in V$, existen escalares $a_{ij}$ con
\[
v_j = \sum_{i=1}^{r} a_{ij} u_i.
\]
Una combinación lineal nula $\sum_{j=1}^{k} x_j v_j = 0$, sustituyendo, es
\[
\sum_{i=1}^{r} \left( \sum_{j=1}^{k} a_{ij} x_j \right) u_i = 0,
\]
y basta anular cada coeficiente, es decir, resolver el sistema homogéneo
\[
\begin{cases}
a_{11}x_1 + a_{12}x_2 + \dots + a_{1k}x_k = 0 \\
a_{21}x_1 + a_{22}x_2 + \dots + a_{2k}x_k = 0 \\
\vdots \\
a_{r1}x_1 + a_{r2}x_2 + \dots + a_{rk}x_k = 0
\end{cases}
\]
de $r$ ecuaciones y $k$ incógnitas. Como $k > r$, tiene una solución no trivial, con algún $x_j \neq 0$, de modo que $\sum_{j} x_j v_j = 0$ con coeficientes no todos nulos y $\{v_1, \dots, v_k\}$ sería linealmente dependiente, contra la hipótesis. Luego $k \le r$.
\end{proof}

\begin{corolario}{Invariancia del Número de Elementos de una Base}{invariancia_base}
Todas las bases de un mismo $\mathbb{K}$-espacio vectorial de dimensión finita $V$ tienen la misma cantidad de elementos.
\end{corolario}

\begin{proof}
Sean $B_1 = \{v_1, \dots, v_n\}$ y $B_2 = \{w_1, \dots, w_m\}$ dos bases de $V$. Considerando $B_1$ como generador y $B_2$ como conjunto linealmente independiente, el Teorema \ref{teo:teorema_intercambio} da $m \le n$. Invirtiendo los roles ($B_2$ generador, $B_1$ independiente), el mismo teorema da $n \le m$. Luego $n = m$, es decir, ambas bases tienen la misma cantidad de elementos.
\end{proof}

\begin{definicion}{Dimensión de un Espacio Vectorial}{definicion_dimension}
Sea $V$ un $\mathbb{K}$-espacio vectorial de dimensión finita.
La \textbf{dimensión} de $V$, denotada $\dim(V)$, es la cantidad de vectores de cualquier base de $V$.

Este número está bien definido (es único para cada espacio) gracias al Corolario \ref{cor:invariancia_base}, el cual garantiza que todas las bases de un mismo espacio vectorial poseen la misma cantidad de elementos.
\end{definicion}

\obs Para $V = \{0\}$ se conviene que su base es el conjunto vacío $\emptyset$; como tiene $0$ elementos, $\dim(\{0\}) = 0$.

\begin{proposicion}{Relación entre Subespacios y Dimensión}{relacion_subespacios}
Sea $V$ un $\mathbb{K}$-espacio vectorial. Sean $S$ y $T$ subespacios de dimensión finita de $V$ tales que $T \subseteq S$. Entonces se cumple:
\begin{enumerate}
    \item $\dim(T) \le \dim(S)$.
    \item Si $\dim(T) = \dim(S)$, entonces $T = S$.
\end{enumerate}
\end{proposicion}

\begin{proof}
\textit{(1)} Sea $\mathcal{B}_T = \{v_1, \dots, v_k\}$ una base de $T$, con $k = \dim(T)$. Como $T \subseteq S$, sus vectores están en $S$ y son linealmente independientes. Sea $\mathcal{B}_S = \{u_1, \dots, u_r\}$ una base de $S$, con $r = \dim(S)$. Por el Teorema \ref{teo:teorema_intercambio} (los $v_i$ independientes, los $u_i$ generadores), $k \le r$, es decir, $\dim(T) \le \dim(S)$.

\textit{(2)} Sea $\dim(T) = \dim(S) = n$ y, por el absurdo, supóngase $S \neq T$. Como $T \subseteq S$, existe $x \in S$ con $x \notin T$. Si $\mathcal{B}_T = \{v_1, \dots, v_n\}$ es base de $T$, como $x \notin T = \operatorname{span}(\mathcal{B}_T)$, el conjunto $\{v_1, \dots, v_n, x\}$ es linealmente independiente en $S$, con $n+1$ vectores. Esto es absurdo, pues $\dim(S) = n$ y, por el Corolario \ref{cor:invariancia_base}, no hay más de $n$ vectores independientes en $S$. Luego $S = T$.
\end{proof}

\begin{teorema}{Teorema de la Extensión de la Base}{extension_base}
Sea $V$ un $\mathbb{K}$-espacio vectorial de dimensión finita y sea $A$ un conjunto de vectores linealmente independientes en $V$.
Si $A$ no es una base (es decir, no genera a $V$), entonces es posible agregar vectores a $A$ hasta convertirlo en una base de $V$.
\end{teorema}

\begin{proof}
Sea $A_0 = A$ y considérese el proceso iterativo: si $\operatorname{span}(A_k) = V$, entonces $A_k$ es base y termina; si $\operatorname{span}(A_k) \subsetneq V$, existe $v_{k+1} \in V \setminus \operatorname{span}(A_k)$, y agregar a un conjunto linealmente independiente un vector fuera de su generado conserva la independencia, de modo que $A_{k+1} = A_k \cup \{v_{k+1}\}$ sigue siendo linealmente independiente.

Como $V$ es de dimensión finita, admite un generador con $n$ elementos, y por el Corolario \ref{cor:invariancia_base} no hay conjuntos linealmente independientes con más de $n$ vectores; el proceso se detiene en algún paso $m$ (con $|A_m| \le n$). Al detenerse, $\operatorname{span}(A_m) = V$, y como $A_m$ es linealmente independiente, es la base buscada que extiende a $A$.
\end{proof}

\begin{teorema}{Teorema de la Dimensión de la Suma}{dim_suma}
Sean $S$ y $T$ subespacios de dimensión finita de un $\mathbb{K}$-espacio vectorial $V$. Entonces:
\[
\dim(S+T) = \dim(S) + \dim(T) - \dim(S \cap T).
\]
\end{teorema}

\begin{proof}
Sea $\{v_1, \dots, v_k\}$ una base de $S \cap T$, de modo que $\dim(S \cap T) = k$. Como es linealmente independiente y está contenida en $S$, se la extiende a una base de $S$ agregando $s_1, \dots, s_p$ (Teorema \ref{teo:extension_base}):
\[
B_S = \{v_1, \dots, v_k, s_1, \dots, s_p\}.
\]
Entonces $\dim(S) = k + p$. De manera análoga, se extiende a una base de $T$ agregando $t_1, \dots, t_q$:
\[
B_T = \{v_1, \dots, v_k, t_1, \dots, t_q\}.
\]
Entonces $\dim(T) = k + q$. Se demuestra que
\[
B = \{v_1, \dots, v_k, s_1, \dots, s_p, t_1, \dots, t_q\}
\]
es una base de $S+T$.

\textit{Genera $S+T$:} sea $x \in S+T$, esto es, $x = s + t$ con $s \in S$, $t \in T$. Como $B_S$ es base de $S$ y $B_T$ de $T$,
\[
s = \sum_{i=1}^{k} \alpha_i v_i + \sum_{j=1}^{p} \beta_j s_j, \qquad
t = \sum_{i=1}^{k} \delta_i v_i + \sum_{l=1}^{q} \gamma_l t_l,
\]
y sumando,
\[
x = \sum_{i=1}^{k} (\alpha_i + \delta_i) v_i + \sum_{j=1}^{p} \beta_j s_j + \sum_{l=1}^{q} \gamma_l t_l,
\]
combinación lineal de los elementos de $B$. Luego $B$ genera $S+T$.

\textit{Independencia lineal:} se plantea una combinación lineal nula
\[
\sum_{i=1}^k \alpha_i v_i + \sum_{j=1}^p \beta_j s_j + \sum_{l=1}^q \gamma_l t_l = 0.
\]
Despejando la parte de los vectores exclusivos de $T$,
\[
\sum_{l=1}^q \gamma_l t_l = - \sum_{i=1}^k \alpha_i v_i - \sum_{j=1}^p \beta_j s_j =: y.
\]
El lado izquierdo está en $T$ y el derecho en $S$, de modo que $y \in S \cap T$, y se escribe como combinación de su base:
\[
y = \sum_{i=1}^k \delta_i v_i.
\]
Igualando con el lado izquierdo,
\[
\sum_{l=1}^q \gamma_l t_l - \sum_{i=1}^k \delta_i v_i = 0,
\]
combinación lineal nula de los vectores de $B_T$, linealmente independientes; en particular $\gamma_l = 0$ para todo $l$. Con esto, la ecuación original se reduce a
\[
\sum_{i=1}^k \alpha_i v_i + \sum_{j=1}^p \beta_j s_j = 0,
\]
combinación nula de los vectores de $B_S$, linealmente independientes, de modo que $\alpha_i = 0$ y $\beta_j = 0$. Como todos los escalares son nulos, $B$ es linealmente independiente.

En consecuencia, $B$ es una base de $S+T$ con $k + p + q$ elementos, y
\[
\dim(S) + \dim(T) - \dim(S \cap T) = (k+p) + (k+q) - k = k + p + q = \dim(S+T). \qedhere
\]
\end{proof}

\subsection{Bases Ortogonales y Ortonormales}

\begin{definicion}{Conjuntos y Bases Ortogonales/Ortonormales}{base_orto}
Sea $(V, \langle \cdot, \cdot \rangle)$ un $\mathbb{K}$-espacio vectorial con producto interno y sea $\mathcal{B} \subset V$ un subconjunto de vectores.

\begin{enumerate}
    \item \textbf{Conjunto ortogonal:} $\mathcal{B}$ es un conjunto ortogonal si todos sus vectores son no nulos y ortogonales dos a dos:
    \[
    \forall u, v \in \mathcal{B}, \quad u \neq v \implies \langle u, v \rangle = 0.
    \]

    \item \textbf{Conjunto ortonormal:} $\mathcal{B}$ es un conjunto ortonormal si es ortogonal y, además, cada vector tiene norma unitaria.

    \item \textbf{Base ortogonal/ortonormal:} si $\mathcal{B}$ es una base de $V$ y un conjunto ortogonal/ortonormal, se llama base ortogonal/ortonormal.
\end{enumerate}
\end{definicion}

\begin{proposicion}{Independencia de Conjuntos Ortogonales}{indep_ortogonal}
Si $\{v_1, \dots, v_k\}$ es un conjunto ortogonal de vectores no nulos, entonces el conjunto es linealmente independiente.
\end{proposicion}

\begin{proof}
Se plantea una combinación lineal nula $\sum_{i=1}^{k} \alpha_i v_i = 0$ y se prueba que todos los $\alpha_i$ son cero. Dado un índice $j$, tomando producto interno contra $v_j$ y usando linealidad,
\[
\sum_{i=1}^{k} \alpha_i \langle v_i, v_j \rangle = \langle 0, v_j \rangle = 0.
\]
Por ortogonalidad, $\langle v_i, v_j \rangle = 0$ para $i \neq j$, de modo que sobrevive solo
\[
\alpha_j \langle v_j, v_j \rangle = \alpha_j \|v_j\|^2 = 0.
\]
Como $v_j \neq 0$, $\|v_j\|^2 \neq 0$, luego $\alpha_j = 0$. Como vale para todo $j$, el conjunto es linealmente independiente.
\end{proof}

\begin{teorema}{Fórmula de la Proyección}{proyeccion_general}
Sea $(V, \langle \cdot, \cdot \rangle)$ un $\mathbb{K}$-espacio vectorial con producto interno y sea $S$ un subespacio de dimensión finita.
Sea $\mathcal{B} = \{v_1, \dots, v_k\}$ una base ortogonal de $S$.

Entonces, para todo vector $v \in V$, la proyección ortogonal de $v$ sobre $S$ está dada por:
\[
P_S(v) = \sum_{i=1}^{k} \frac{\langle v, v_i \rangle}{\|v_i\|^2} v_i.
\]
\end{teorema}

\begin{proof}
Como $S$ es de dimensión finita, $V = S \oplus S^\perp$ (Teorema \ref{teo:descomposicion_ortogonal}), de modo que todo $v \in V$ se escribe de manera única como $v = s + t$ con $s \in S$ y $t \in S^\perp$, donde $s = P_S(v)$. Como $\mathcal{B} = \{v_1, \dots, v_k\}$ es base de $S$,
\[
s = \sum_{j=1}^{k} c_j v_j,
\]
y resta determinar los $c_j$. Tomando producto interno de $v$ con $v_i$ y usando linealidad,
\[
\langle v, v_i \rangle = \langle s, v_i \rangle + \langle t, v_i \rangle = \langle s, v_i \rangle,
\]
pues $t \in S^\perp$ y $v_i \in S$ dan $\langle t, v_i \rangle = 0$. Sustituyendo $s$ y usando la ortogonalidad de la base ($\langle v_j, v_i \rangle = 0$ para $j \neq i$),
\[
\langle v, v_i \rangle = \sum_{j=1}^{k} c_j \langle v_j, v_i \rangle = c_i \|v_i\|^2,
\]
de donde $c_i = \langle v, v_i \rangle / \|v_i\|^2$ (los $v_i$ son no nulos). Por lo tanto,
\[
P_S(v) = \sum_{i=1}^{k} \frac{\langle v, v_i \rangle}{\|v_i\|^2} v_i. \qedhere
\]
\end{proof}

\begin{corolario}{Caso Ortonormal}{caso_ortonormal}
Si la base $\mathcal{B} = \{u_1, \dots, u_k\}$ es ortonormal, la fórmula se simplifica:
\[
P_S(v) = \sum_{i=1}^{k} \langle v, u_i \rangle u_i.
\]
\end{corolario}

\begin{teorema}{Proceso de Ortogonalización de Gram-Schmidt}{gram_schmidt}
Sea $(V, \langle \cdot, \cdot \rangle)$ un $\mathbb{K}$-espacio vectorial con producto interno y sea $S$ un subespacio de dimensión finita.
Dada una base arbitraria $\mathcal{B} = \{u_1, \dots, u_n\}$ de $S$, existe una base ortogonal $\mathcal{O} = \{v_1, \dots, v_n\}$ de $S$ que satisface la propiedad de conservación de los subespacios generados:
\[
\operatorname{span}\{u_1, \dots, u_k\} = \operatorname{span}\{v_1, \dots, v_k\}, \quad \forall k \in \{1, \dots, n\}.
\]
\end{teorema}

\begin{proof}
Se definen los vectores recursivamente:
\begin{align*}
v_1 &= u_1 \\
v_2 &= u_2 - \frac{\langle u_2, v_1 \rangle}{\|v_1\|^2} v_1 \\
&\vdots \\
v_k &= u_k - \sum_{j=1}^{k-1} \frac{\langle u_k, v_j \rangle}{\|v_j\|^2} v_j.
\end{align*}

\textit{Los $v_k$ son no nulos.} Si fuera $v_k = 0$, la definición de $v_k$ daría
$u_k \in \operatorname{span}\{v_1, \dots, v_{k-1}\}$, que por la conservación de los generados
(probada más abajo, con índice $k-1$) es $\operatorname{span}\{u_1, \dots, u_{k-1}\}$; esto
contradice la independencia lineal de $\mathcal{B}$. En particular, las divisiones por
$\|v_j\|^2$ están bien definidas.

\textit{$\{v_1, \dots, v_n\}$ es ortogonal.} Por inducción sobre $k$. El caso $\{v_1\}$ vale por vacuidad. Supóngase $\{v_1, \dots, v_{k-1}\}$ ortogonal; hay que ver que $\langle v_k, v_i \rangle = 0$ para $i < k$. Por la definición de $v_k$ y linealidad,
\[
\langle v_k, v_i \rangle = \langle u_k, v_i \rangle - \sum_{j=1}^{k-1} \frac{\langle u_k, v_j \rangle}{\|v_j\|^2} \langle v_j, v_i \rangle.
\]
Por la hipótesis inductiva, $\langle v_j, v_i \rangle = 0$ salvo $j = i$, de modo que
\[
\langle v_k, v_i \rangle = \langle u_k, v_i \rangle - \frac{\langle u_k, v_i \rangle}{\|v_i\|^2} \|v_i\|^2 = 0.
\]
Luego $\{v_1, \dots, v_k\}$ es ortogonal.

\textit{$\operatorname{span}\{u_1, \dots, u_p\} = \operatorname{span}\{v_1, \dots, v_p\}$ para todo $p$.} Por inducción sobre $p$. El caso $p = 1$ vale porque $v_1 = u_1$. Supóngase la igualdad para $p-1$. De la construcción,
\[
v_p = u_p - \sum_{j=1}^{p-1} c_j v_j \implies u_p = v_p + \sum_{j=1}^{p-1} c_j v_j.
\]
\textit{($\subseteq$).} $u_p \in \operatorname{span}\{v_1, \dots, v_p\}$ por la igualdad anterior, y $u_1, \dots, u_{p-1} \in \operatorname{span}\{v_1, \dots, v_{p-1}\} \subseteq \operatorname{span}\{v_1, \dots, v_p\}$ por hipótesis inductiva. \textit{($\supseteq$).} $v_p = u_p - w$ con $w \in \operatorname{span}\{v_1, \dots, v_{p-1}\} = \operatorname{span}\{u_1, \dots, u_{p-1}\}$, de modo que $v_p \in \operatorname{span}\{u_1, \dots, u_p\}$. Por doble inclusión, los subespacios coinciden.

Con $p = n$, $\mathcal{O}$ genera $S$, y por ser un conjunto ortogonal de vectores no nulos es linealmente independiente (Proposición \ref{prop:indep_ortogonal}); luego $\mathcal{O}$ es una base ortogonal de $S$.
\end{proof}

\obs La construcción de los vectores $v_k$ se llama \textbf{proceso de ortogonalización de Gram-Schmidt}.

\subsection{Coordenadas y Representación Matricial}

\begin{definicion}{Vector de Coordenadas}{vector_coordenadas}
Sea $V$ un espacio vectorial de dimensión finita $n$ y sea $\mathcal{B} = \{v_1, v_2, \dots, v_n\}$ una base ordenada de $V$.
Para cualquier vector $v \in V$, existen únicos escalares $\alpha_1, \dots, \alpha_n$ tales que:
\[ v = \alpha_1 v_1 + \alpha_2 v_2 + \dots + \alpha_n v_n \]
El \textbf{vector de coordenadas} de $v$ en la base $\mathcal{B}$, denotado $[v]_\mathcal{B}$, es el vector columna de estos escalares:
\[ [v]_\mathcal{B} = \begin{pmatrix} \alpha_1 \\ \alpha_2 \\ \vdots \\ \alpha_n \end{pmatrix} \in \mathbb{K}^n \]
\end{definicion}

\obs Hay que distinguir el vector $v$ (objeto abstracto) de $[v]_\mathcal{B}$: al cambiar la base $\mathcal{B}$, el vector $v$ sigue siendo el mismo, pero sus coordenadas $[v]_\mathcal{B}$ cambian.

\subsection{Existencia de Bases en Espacios Generales y Curiosidades}

La lectura de esta sección es opcional. Se sugiere familiarizarse previamente con el Lema de Zorn y el Teorema de Categoría de Baire (en alguna de sus versiones).

\begin{teorema}{Existencia de Base (Caso General)}{existencia_base_zorn}
Todo espacio vectorial $V$ posee una base.
\end{teorema}

\begin{proof}
Sea $\mathcal{X}$ el conjunto de los subconjuntos linealmente independientes de $V$, ordenado por inclusión: $A \le B \iff A \subseteq B$. Es no vacío, pues $\emptyset \in \mathcal{X}$.

Sea $\mathcal{Y}$ una cadena de $\mathcal{X}$; hay que ver que tiene cota superior en $\mathcal{X}$. Sea $C = \bigcup_{B \in \mathcal{Y}} B$; basta ver que $C$ es linealmente independiente. Por el absurdo, supóngase que es dependiente: existen $v_1, \dots, v_n \in C$ y escalares $\lambda_i$ no todos nulos con $\sum_{i=1}^n \lambda_i v_i = 0_V$. Cada $v_i$ está en algún $B_i \in \mathcal{Y}$; como $\{B_1, \dots, B_n\}$ es finito y $\mathcal{Y}$ es cadena, tiene un máximo $B_{\max}$, y todos los $v_i \in B_{\max}$. Pero entonces $B_{\max}$ contendría un conjunto dependiente, contra $B_{\max} \in \mathcal{X}$. Luego $C$ es linealmente independiente y, como $B \subseteq C$ para todo $B \in \mathcal{Y}$, es cota superior de la cadena.

Por el lema de Zorn, $\mathcal{X}$ tiene un elemento maximal $G$, que es una base de $V$:
\begin{itemize}
    \item \textit{Independencia:} $G \in \mathcal{X}$ es linealmente independiente.
    \item \textit{Generación:} si $G$ no generara $V$, habría $v \notin \operatorname{span}(G)$, y $G' = G \cup \{v\}$ sería linealmente independiente (si fuera dependiente, $v$ sería combinación lineal de $G$), con $G \subsetneq G'$ en $\mathcal{X}$, contra la maximalidad de $G$.
\end{itemize}
Luego $G$ genera $V$ y es una base.
\end{proof}

\begin{teorema}{Cardinalidad de la Base}{base_frechet}
Sea $V$ un espacio vectorial normado completo (espacio de Banach) de dimensión infinita.
Entonces, cualquier base de Hamel de $V$ es no numerable.
\end{teorema}

\begin{proof}
Por el absurdo, supóngase que $V$ admite una base de Hamel numerable $\mathcal{B} = \{e_1, e_2, \dots\}$, y sea
\[
F_n = \operatorname{span}\{e_1, \dots, e_n\}.
\]
Estos subespacios cumplen:
\begin{enumerate}
    \item \textit{Son cerrados:} todo subespacio de dimensión finita es cerrado.
    \item \textit{Tienen interior vacío:} como $V$ es de dimensión infinita, $F_n \subsetneq V$ es propio; un subespacio propio con interior no vacío contendría un entorno del origen y generaría todo el espacio, contra $F_n \neq V$. Luego $\operatorname{int}(F_n) = \emptyset$.
\end{enumerate}
Como $\mathcal{B}$ es base, todo $x \in V$ es combinación lineal finita de sus elementos, de modo que
\[
V = \bigcup_{n=1}^{\infty} F_n,
\]
una unión numerable de cerrados con interior vacío (nunca densos). Esto contradice el Teorema de Categoría de Baire, según el cual un espacio métrico completo no es de primera categoría. Luego la base de Hamel no puede ser numerable.
\end{proof}

\obs No debe confundirse este resultado con las bases de espacios de Hilbert o series de Fourier.
\begin{itemize}
    \item Una base de Hamel (la algebraica) requiere combinaciones lineales finitas. En un Banach de dimensión infinita, estas bases son necesariamente no numerables.
    \item Una base de Schauder (la topológica) permite combinaciones lineales infinitas.
\end{itemize}
El teorema anterior dice que no se puede generar un Banach infinito usando solo sumas finitas de una cantidad numerable de vectores.

\obs Sobre la relación entre separabilidad y bases de Schauder:
\begin{enumerate}
    \item Si un espacio de Banach tiene una base de Schauder, entonces es necesariamente separable.
    \item El recíproco es falso: no todo espacio de Banach separable admite una base de Schauder.
\end{enumerate}
Este fue un problema abierto famoso durante 40 años, resuelto negativamente por Per Enflo en 1973, quien construyó un espacio de Banach separable que no admite ninguna base de Schauder.

\newpage

\section{Transformaciones Lineales}

\subsection{Definición y Primeras Propiedades}

\begin{definicion}{Transformación Lineal}{def_trans_lineal}
Sean $V$ y $W$ dos $\mathbb{K}$-espacios vectoriales.
Una función $T: V \to W$ se dice que es una transformación lineal si satisface las siguientes dos condiciones:

\begin{enumerate}
    \item Aditividad:
    Para todo par de vectores $u, v \in V$:
    \[ T(u + v) = T(u) + T(v) \]
    \item Homogeneidad:
    Para todo escalar $\lambda \in \mathbb{K}$ y todo vector $v \in V$:
    \[ T(\lambda v) = \lambda T(v) \]
\end{enumerate}
\end{definicion}

\begin{proposicion}{Propiedades}{prop_estructurales_tl}
Sea $T: V \to W$ una transformación lineal. Se cumplen las siguientes propiedades:

\begin{enumerate}
    \item El origen es un punto fijo: $T(0_V) = 0_W$.
    \item Si se conocen los valores que toma $T$ sobre una base $\{v_1, \dots, v_n\}$ de $V$, entonces $T(v)$ queda determinado para cualquier vector $v \in V$.
    \item Si $S$ es un subespacio vectorial de $V$, entonces su imagen $T(S)$ es un subespacio vectorial de $W$.
    \item Dadas dos transformaciones lineales $T: U \to V$ y $L: V \to W$, su composición $L \circ T: U \to W$ es también una transformación lineal.
    \item Si $W'$ es un subespacio de $W$, entonces la preimagen de $W'$, definida como:
    \[ T^{-1}(W') = \{ v \in V \mid T(v) \in W' \} \]
    es un subespacio vectorial de $V$.
\end{enumerate}
\end{proposicion}

\begin{proof}
\textit{(1)} De $0_V + 0_V = 0_V$ y la aditividad, $T(0_V) = T(0_V) + T(0_V)$; restando $T(0_V)$, queda $0_W = T(0_V)$.

\textit{(2)} Sea $\mathcal{B} = \{v_1, \dots, v_n\}$ una base de $V$. Cada $v \in V$ se escribe de manera única $v = \alpha_1 v_1 + \dots + \alpha_n v_n$, y por linealidad
\[
T(v) = \alpha_1 T(v_1) + \dots + \alpha_n T(v_n),
\]
de modo que $T(v)$ queda determinado por los $\alpha_i$ y las imágenes $T(v_i)$.

\textit{(3)} Sea $S$ un subespacio de $V$; se verifica que $T(S)$ es subespacio de $W$. Como $0_V \in S$ y $T(0_V) = 0_W$, $0_W \in T(S)$. Si $y_1 = T(s_1)$ e $y_2 = T(s_2)$ con $s_1, s_2 \in S$, entonces $y_1 + y_2 = T(s_1 + s_2)$ con $s_1 + s_2 \in S$; y para $\lambda \in \mathbb{K}$, $\lambda y_1 = T(\lambda s_1)$ con $\lambda s_1 \in S$. Luego $T(S)$ es subespacio.

\textit{(4)} Para $u, v \in U$ y $\lambda \in \mathbb{K}$, $(L \circ T)(u + v) = L(T(u) + T(v)) = (L \circ T)(u) + (L \circ T)(v)$ y $(L \circ T)(\lambda u) = \lambda (L \circ T)(u)$, de modo que $L \circ T$ es lineal.

\textit{(5)} Sea $Z = T^{-1}(W')$. Como $T(0_V) = 0_W \in W'$, $0_V \in Z$. Si $u, v \in Z$, entonces $T(u), T(v) \in W'$ y $T(u+v) = T(u) + T(v) \in W'$, de modo que $u + v \in Z$. Si $v \in Z$ y $\lambda \in \mathbb{K}$, $T(\lambda v) = \lambda T(v) \in W'$, de modo que $\lambda v \in Z$. Luego $T^{-1}(W')$ es subespacio de $V$.
\end{proof}

\subsection{El Teorema de Representación Matricial}

Como complemento visual de la relación entre una transformación lineal y el producto matricial, puede consultarse \href{https://youtu.be/3gxjK9zkKqc?si=_vvIvbiUcNyruPvx}{este video explicativo}.

\begin{definicion}{Base Canónica de $\mathbb{K}^n$}{base_canonica_kn}
Se llama \textbf{base canónica} de $\mathbb{K}^n$ al conjunto $\{e_1, \dots, e_n\}$ donde $e_j$ tiene un $1$ en la posición $j$ y $0$ en las demás. Todo $x \in \mathbb{K}^n$ se escribe de forma única como $x = \sum_{j=1}^n x_j e_j$.
\end{definicion}

\begin{teorema}{Representación Matricial}{teo_matriz_asociada}
Sean $V$ y $W$ espacios de dimensiones $n$ y $m$, con bases $\mathcal{B}_V$ y $\mathcal{B}_W$. Sea $T: V \to W$ lineal.
Existe una única matriz $A \in \mathbb{K}^{m \times n}$, denotada $[T]_{\mathcal{B}_V}^{\mathcal{B}_W}$, tal que para todo vector $v \in V$:
\[
[T(v)]_{\mathcal{B}_W} = [T]_{\mathcal{B}_V}^{\mathcal{B}_W} \cdot [v]_{\mathcal{B}_V}
\]
Las columnas de dicha matriz se construyen calculando las coordenadas de las imágenes de la base del dominio:
\[ \text{Col}_j(A) = [T(v_j)]_{\mathcal{B}_W} \]
\end{teorema}

\begin{proof}
Para la existencia, se construye $A$ con la fórmula de las columnas del enunciado. Sea $v \in V$ con coordenadas $x = [v]_{\mathcal{B}_V}$, esto es, $v = x_1 v_1 + \dots + x_n v_n$. En $\mathbb{K}^n$, $x = x_1 e_1 + \dots + x_n e_n$, y por la distributividad del producto matricial,
\[ A x = x_1 (A e_1) + \dots + x_n (A e_n) = x_1 [T(v_1)]_{\mathcal{B}_W} + \dots + x_n [T(v_n)]_{\mathcal{B}_W}, \]
pues $A e_j$ es la columna $j$ de $A$, igual a $[T(v_j)]_{\mathcal{B}_W}$ por construcción. Por la linealidad de tomar coordenadas y la de $T$,
\[ A x = \left[ x_1 T(v_1) + \dots + x_n T(v_n) \right]_{\mathcal{B}_W} = \left[ T( x_1 v_1 + \dots + x_n v_n ) \right]_{\mathcal{B}_W} = [T(v)]_{\mathcal{B}_W}. \]
Luego la matriz construida cumple la propiedad requerida.

Para la unicidad, sean $A$ y $A'$ dos matrices que cumplen la propiedad. Tomando $v = v_j$ (cuyas coordenadas son $[v_j]_{\mathcal{B}_V} = e_j$), se obtiene $A e_j = [T(v_j)]_{\mathcal{B}_W} = A' e_j$ para todo $j$: las columnas coinciden y $A = A'$.
\end{proof}

\subsection{Composición y Producto de Matrices}

\obs \textit{Orden de la multiplicación.} Dadas $T: U \to V$ y $L: V \to W$, la matriz de la composición $L \circ T$ es el producto de las matrices asociadas:
\[ [L \circ T]_{\mathcal{B}_U}^{\mathcal{B}_W} = [L]_{\mathcal{B}_V}^{\mathcal{B}_W} \cdot [T]_{\mathcal{B}_U}^{\mathcal{B}_V} \]
Nótese que el orden se invierte respecto a la escritura habitual de izquierda a derecha, pero coincide con el orden de evaluación de funciones: $(L \circ T)(x) = L(T(x))$.

\subsection{El Producto Matricial como Composición}

La regla de multiplicación de matrices (fila por columna) se obtiene al exigir que el producto de matrices represente la composición de transformaciones lineales.

Sean dos transformaciones lineales consecutivas
\[ T: U \to V \quad \text{y} \quad S: V \to W, \]
con matrices asociadas $A$ (para $T$) y $B$ (para $S$). Dado $x \in U$, al aplicar $T$ se obtiene $y = Ax$, y al aplicar $S$ se obtiene $z = By$. Se busca la matriz $C$ que represente la composición $S \circ T$, esto es, tal que $z = Cx$. Sustituyendo,
\[ z = B(Ax). \]

A nivel de componentes, la componente $i$-ésima de $z$ es $z_i = \sum_{k} B_{ik} y_k$, y la componente $k$-ésima de $y$ es $y_k = \sum_{j} A_{kj} x_j$. Sustituyendo,
\[ z_i = \sum_{k} B_{ik} \left( \sum_{j} A_{kj} x_j \right) = \sum_{j} \left( \sum_{k} B_{ik} A_{kj} \right) x_j. \]
Para que esto adopte la forma $z_i = \sum_{j} C_{ij} x_j$, la entrada $(i, j)$ de $C$ debe ser el término entre paréntesis.

\obs
La fórmula del producto matricial
\[ (BA)_{ij} = \sum_{k} B_{ik} A_{kj} \]
es la única definición que garantiza que multiplicar matrices equivalga a componer las transformaciones lineales correspondientes.

\subsubsection{La Asociatividad como Consecuencia}

La prueba aritmética de que el producto de matrices es asociativo ($(AB)C = A(BC)$) es laboriosa, pues involucra sumatorias triples. La prueba conceptual es directa al usar que multiplicar matrices equivale a componer transformaciones lineales.

\begin{proposicion}{Asociatividad del Producto Matricial}{asoc_matrices}
El producto de matrices es asociativo. Es decir, para matrices $A, B, C$ de tamaños compatibles:
\[ (AB)C = A(BC) \]
\end{proposicion}

\begin{proof}
% Corrección: con T, S, R representadas por A, B, C, el producto AB representa T∘S
% (el orden se invierte), no S∘T como decía antes.
Sean $T, S, R$ las transformaciones lineales representadas por $A, B, C$. Como el producto de matrices representa la composición en el mismo orden ($AB$ representa $T \circ S$), la matriz $(AB)C$ representa la composición $(T \circ S) \circ R$ y la matriz $A(BC)$ representa $T \circ (S \circ R)$. Como la composición de funciones es asociativa, ambas coinciden, y por lo tanto $(AB)C = A(BC)$.
\end{proof}

\subsection{Núcleo e Imagen}

\begin{definicion}{Núcleo e Imagen}{def_nucleo_imagen}
Sea $T: V \to W$ una transformación lineal.
\begin{enumerate}
    \item Núcleo: $\ker(T) = \{v \in V \mid T(v) = 0_W\}$.
    \item Imagen: $\operatorname{Im}(T) = \{w \in W \mid \exists v \in V, T(v) = w\}$.
\end{enumerate}
\end{definicion}

\begin{proposicion}{Estructura de Subespacio}{prop_subespacios_T}
El núcleo $\ker(T)$ es un subespacio de $V$ y la imagen $\operatorname{Im}(T)$ es un subespacio de $W$.
\end{proposicion}

\begin{proof}
El núcleo es la preimagen del subespacio $\{0_W\}$, subespacio por la Proposición \ref{prop:prop_estructurales_tl} (5). La imagen es $T(V)$, subespacio por la Proposición \ref{prop:prop_estructurales_tl} (3).
\end{proof}

\begin{proposicion}{Monomorfismos y Núcleo}{prop_monomorfismo}
$T$ es inyectiva si y solo si $\ker(T) = \{0_V\}$.
\end{proposicion}

\begin{proof}
Se prueban las dos implicaciones.

\textit{($\Rightarrow$).} Sea $T$ inyectiva y $v \in \ker(T)$, esto es, $T(v) = 0_W = T(0_V)$; por inyectividad, $v = 0_V$. Luego $\ker(T) = \{0_V\}$.

\textit{($\Leftarrow$).} Sea $\ker(T) = \{0_V\}$ y $u, v$ con $T(u) = T(v)$. Por linealidad, $T(u - v) = 0_W$, de modo que $u - v \in \ker(T) = \{0_V\}$, es decir, $u = v$. Luego $T$ es inyectiva.
\end{proof}

\begin{proposicion}{Caracterización de Sobreyectividad}{prop_epi_img}
Sea $T: V \to W$ una transformación lineal.
$T$ es sobreyectiva si y solo si su imagen coincide con todo el codominio $W$.
\end{proposicion}

\begin{proof}
Se prueban las dos implicaciones.

\textit{($\Rightarrow$).} Si $T$ es sobreyectiva, para todo $w \in W$ existe $v$ con $T(v) = w$, de modo que $W \subseteq \operatorname{Im}(T)$; como siempre $\operatorname{Im}(T) \subseteq W$, se tiene $\operatorname{Im}(T) = W$.

\textit{($\Leftarrow$).} Si $\operatorname{Im}(T) = W$, dado $w \in W = \operatorname{Im}(T)$ existe $v$ con $T(v) = w$, de modo que $T$ es sobreyectiva.
\end{proof}

\subsection{Interpretación Matricial: Filas y Columnas}

\begin{proposicion}{Imagen y Columnas}{prop_img_cols}
Sea $A$ la matriz asociada a $T$. Entonces la imagen de $T$ es el espacio generado por las columnas de $A$:
\[ \operatorname{Im}(T) = \operatorname{Col}(A) \]
\end{proposicion}

\begin{proof}
Sea $x \in \mathbb{K}^n$ y $C_1, \dots, C_n$ las columnas de $A$. Entonces
\[
Ax = \begin{bmatrix} | & & | \\ C_1 & \dots & C_n \\ | & & | \end{bmatrix} \begin{bmatrix} x_1 \\ \vdots \\ x_n \end{bmatrix} = x_1 C_1 + x_2 C_2 + \dots + x_n C_n,
\]
combinación lineal de las columnas. Como todo vector de la imagen es de la forma $Ax$, la imagen es exactamente el conjunto de combinaciones lineales de las columnas, es decir, $\operatorname{Col}(A)$.
\end{proof}

\begin{proposicion}{Núcleo y Filas}{prop_ker_rows}
% TODO: para K = C el argumento del complemento ortogonal requiere el producto hermítico
% (conjugar); el enunciado y la prueba están escritos para el caso real.
El núcleo de $A$ es el complemento ortogonal del espacio fila:
\[ \ker(A) = (\operatorname{Fil}(A))^\perp \]
\end{proposicion}

\begin{proof}
Sean $R_1, \dots, R_m$ las filas de $A$. Un vector $x$ pertenece al núcleo si y solo si $Ax = \mathbf{0}$, lo que, desglosando el producto fila por fila mediante el producto interno estándar, equivale a
\[
Ax = \begin{bmatrix} \text{---} & R_1 & \text{---} \\ & \vdots & \\ \text{---} & R_m & \text{---} \end{bmatrix} \begin{bmatrix} x \end{bmatrix} = \begin{bmatrix} R_1 \cdot x \\ \vdots \\ R_m \cdot x \end{bmatrix} = \begin{bmatrix} 0 \\ \vdots \\ 0 \end{bmatrix},
\]
es decir, a $R_i \cdot x = 0$ para todo $i$. Esto significa que $x$ es ortogonal a todas las filas, y un vector ortogonal a un conjunto de generadores es ortogonal a todo el subespacio que generan. Luego $\ker(A) = (\operatorname{Fil}(A))^\perp$.
\end{proof}

\subsection{Teorema de la Dimensión y Rango}

\begin{teorema}{Teorema de la Dimensión (Rango-Nulidad)}{teo_dimension}
Sea $T: V \to W$ lineal con $\dim(V) < \infty$. Entonces:
\[ \dim(V) = \dim(\ker(T)) + \dim(\operatorname{Im}(T)) \]
\end{teorema}

\begin{proof}
La estrategia es construir una base de $\operatorname{Im}(T)$ a partir de una base del núcleo extendida a una base de $V$. Sea $n = \dim(V)$ y $k = \dim(\ker(T))$, y tómese una base del núcleo $\mathcal{B}_K = \{u_1, \dots, u_k\}$. Por ser un conjunto linealmente independiente en $V$, por el Teorema de la Extensión de la Base \ref{teo:extension_base} se completa a una base
\[ \mathcal{B}_V = \{u_1, \dots, u_k, v_1, \dots, v_r\}, \]
con $n = k + r$. Se prueba que $B' = \{T(v_1), \dots, T(v_r)\}$ es base de $\operatorname{Im}(T)$, de donde $\dim(\operatorname{Im}(T)) = r$.

\textit{$B'$ genera $\operatorname{Im}(T)$.} Sea $w \in \operatorname{Im}(T)$, esto es, $w = T(x)$ para algún $x \in V$. Escribiendo $x = \sum_{i=1}^k \alpha_i u_i + \sum_{j=1}^r \beta_j v_j$ y aplicando $T$, como $T(u_i) = 0_W$,
\[ w = \sum_{j=1}^r \beta_j T(v_j), \]
combinación lineal de $B'$.

\textit{$B'$ es linealmente independiente.} Supóngase $\sum_{j=1}^r c_j T(v_j) = 0_W$. Por linealidad, $T(\sum_j c_j v_j) = 0_W$, de modo que $\sum_j c_j v_j \in \ker(T)$ y se escribe $\sum_{j=1}^r c_j v_j = \sum_{i=1}^k d_i u_i$. Entonces
\[ \sum_{j=1}^r c_j v_j - \sum_{i=1}^k d_i u_i = 0_V \]
es una combinación lineal nula de los vectores de la base $\mathcal{B}_V$, lo que fuerza todos los escalares a ser $0$; en particular, $c_j = 0$ para todo $j$.

Luego $\dim(\operatorname{Im}(T)) = r$ y $\dim(V) = k + r = \dim(\ker(T)) + \dim(\operatorname{Im}(T))$.
\end{proof}

\begin{corolario}{Rango de una Matriz}{cor_rango_matriz}
Sea $A \in \mathbb{K}^{m \times n}$. El número de filas linealmente independientes es igual al número de columnas linealmente independientes.
\[ \dim(\operatorname{Col}(A)) = \dim(\operatorname{Fil}(A)) \]
\end{corolario}

\begin{proof}
Aplicando el Teorema de la Dimensión \ref{teo:teo_dimension} a la transformación $T(x) = Ax$ y usando que $\dim(\mathbb{K}^n) = n$,
\[ n = \dim(\ker(A)) + \dim(\operatorname{Col}(A)). \quad (*) \]
Por la Proposición \ref{prop:prop_ker_rows}, $\ker(A) = (\operatorname{Fil}(A))^\perp$ dentro de $\mathbb{K}^n$. Como la suma de las dimensiones de un subespacio y su complemento ortogonal es la dimensión del espacio total,
\[ \dim(\operatorname{Fil}(A)) + \dim(\ker(A)) = n. \quad (**) \]
Igualando $(*)$ y $(**)$ y cancelando $\dim(\ker(A))$, se obtiene $\dim(\operatorname{Col}(A)) = \dim(\operatorname{Fil}(A))$.
\end{proof}

\subsection{Tipos de Morfismos}

\begin{definicion}{Homomorfismo y Endomorfismo}{homo_endo}
Sean $V$ y $W$ dos $\mathbb{K}$-espacios vectoriales. A toda transformación lineal $f: V \to W$ se la denomina homomorfismo.
Si el espacio de salida y de llegada son el mismo ($V = W$), es decir $f: V \to V$, se dice que $f$ es un endomorfismo.

\end{definicion}

\obs El concepto de autovalor y autovector tiene sentido únicamente para endomorfismos, ya que se compara un vector $v$ con su imagen $f(v)$ en el mismo espacio.

\begin{definicion}{Monomorfismo, Epimorfismo e Isomorfismo}{mono_epi_iso}
Sea $f: V \to W$ una transformación lineal.
\begin{enumerate}
    \item Monomorfismo: si $f$ es inyectiva, es decir,
    \[ v_1 \neq v_2 \implies f(v_1) \neq f(v_2), \]
    lo que equivale a $\ker(f) = \{0\}$ (Proposición \ref{prop:prop_monomorfismo}).
    \item Epimorfismo: si $f$ es sobreyectiva, es decir, $\operatorname{Im}(f) = W$.
    \item Isomorfismo: si $f$ es biyectiva (inyectiva y sobreyectiva a la vez). En este caso, $f$ es invertible y su inversa $f^{-1}: W \to V$ es también lineal.
\end{enumerate}
\end{definicion}

\begin{definicion}{Automorfismo}{automorfismo}
Si $f: V \to V$ es un endomorfismo que además es biyectivo (es decir, es un isomorfismo de un espacio en sí mismo), se denomina automorfismo.
\end{definicion}

\obs El conjunto de todos los automorfismos de $V$, junto con la operación de composición de funciones, forma un grupo conocido como el Grupo General Lineal de $V$, denotado por $GL(V)$. En términos matriciales, corresponde a las matrices cuadradas invertibles.

\subsection{Equivalencias en Endomorfismos}

\begin{corolario}{Equivalencias}{cor_endomorfismos}
Sea $T: V \to V$ una transformación lineal ($\dim(V)=n$). Son equivalentes:
\begin{enumerate}
    \item[(i)] $T$ es biyectiva.
    \item[(ii)] $T$ es inyectiva.
    \item[(iii)] $T$ es sobreyectiva.
\end{enumerate}
\end{corolario}

\begin{proof}
Se prueba (ii) $\Leftrightarrow$ (iii); junto con las definiciones, esto da la equivalencia con (i).

\textit{(ii $\Rightarrow$ iii).} Si $T$ es inyectiva, $\dim(\ker(T)) = 0$, y por el Teorema de la Dimensión \ref{teo:teo_dimension}, $\dim(\operatorname{Im}(T)) = n$. Como $\operatorname{Im}(T)$ es subespacio de $V$ con su misma dimensión, $\operatorname{Im}(T) = V$, de modo que $T$ es sobreyectiva.

\textit{(iii $\Rightarrow$ ii).} Si $T$ es sobreyectiva, $\operatorname{Im}(T) = V$ y $\dim(\operatorname{Im}(T)) = n$. Por el Teorema de la Dimensión \ref{teo:teo_dimension}, $\dim(\ker(T)) = 0$, esto es, $\ker(T) = \{0\}$, de modo que $T$ es inyectiva.

Al ser equivalentes inyectividad y sobreyectividad, cualquiera de ellas implica la biyectividad (i).
\end{proof}

\begin{proposicion}{Versión Matricial}{prop_matrices_invertibles}
Sea $A \in \mathbb{K}^{n \times n}$. Son equivalentes:
\begin{enumerate}
    \item $A$ es invertible.
    \item Las columnas de $A$ son linealmente independientes.
    \item Las columnas de $A$ generan $\mathbb{K}^n$.
\end{enumerate}
\end{proposicion}

\begin{proof}
Sea $T: \mathbb{K}^n \to \mathbb{K}^n$ la transformación $T(x) = Ax$. La matriz $A$ es invertible si y solo si $T$ es biyectiva. Como $Ax = x_1 C_1 + \dots + x_n C_n$ (combinación lineal de las columnas), $\ker(T) = \{0\}$ equivale a que las columnas sean linealmente independientes, e $\operatorname{Im}(T) = \operatorname{Col}(A) = \mathbb{K}^n$ (Proposición \ref{prop:prop_img_cols}) equivale a que las columnas generen $\mathbb{K}^n$. Por el Corolario \ref{cor:cor_endomorfismos}, inyectividad, sobreyectividad y biyectividad de $T$ son equivalentes, de donde las tres condiciones del enunciado también lo son.
\end{proof}

\subsection{Ejemplo: Proyección Ortogonal}

\begin{proposicion}{Proyección Ortogonal}{prop_proyeccion}
Sea $V$ un espacio con producto interno y $S$ un subespacio tal que $V = S \oplus S^\perp$.
Se define $P_S(v) = s$, donde $v = s + s^\perp$.
Entonces $P_S: V \to V$ es lineal y su matriz asociada $M$ cumple $M^2 = M$ (idempotencia).
\end{proposicion}

\begin{proof}
\textit{Linealidad.} Sean $u, v \in V$ con descomposiciones $u = u_S + u_\perp$ y $v = v_S + v_\perp$ respecto de $V = S \oplus S^\perp$, de modo que $P_S(u) = u_S$ y $P_S(v) = v_S$. Entonces
\[ u + v = (u_S + v_S) + (u_\perp + v_\perp), \]
con $u_S + v_S \in S$ y $u_\perp + v_\perp \in S^\perp$; por definición de la proyección, $P_S(u+v) = u_S + v_S = P_S(u) + P_S(v)$. Análogamente, para $\lambda \in \mathbb{K}$, $\lambda u = \lambda u_S + \lambda u_\perp$ con $\lambda u_S \in S$, luego $P_S(\lambda u) = \lambda u_S = \lambda P_S(u)$.

\textit{Idempotencia del operador.} Sea $s = P_S(v) \in S$. Su descomposición es $s = s + 0$, de modo que $P_S(s) = s$, y entonces $P_S(P_S(v)) = s = P_S(v)$, es decir, $P_S^2 = P_S$.

\textit{Idempotencia de la matriz.} Sea $M$ la matriz asociada. La matriz de $P_S \circ P_S$ es $M^2$, y como $P_S \circ P_S = P_S$, se tiene $M^2 = M$.
\end{proof}

\newpage

\section{El Espacio Dual}

\subsection{Definición y Motivación}

\begin{definicion}{Espacio Dual}{def_espacio_dual}
Si $V$ es un espacio vectorial sobre un cuerpo $\mathbb{K}$, el \textbf{espacio dual}, denotado $V^*$, es el conjunto de todas las transformaciones lineales que van de $V$ al cuerpo de escalares.
\[
V^* = \{ f: V \to \mathbb{K} \mid f \text{ es lineal} \}.
\]
A los elementos de $V^*$ se los llama \textbf{funcionales lineales}.
\end{definicion}

El espacio dual interesa por dos motivos:

\begin{itemize}
    \item Distinción objeto/medición: en el producto matricial $x^T y$, los dos vectores no cumplen el mismo papel:
    \begin{itemize}
        \item el vector columna $y \in V$ es el objeto geométrico;
        \item el vector fila $x^T \in V^*$ actúa como operador o regla de medición.
    \end{itemize}
    El espacio dual formaliza esta diferencia: el objeto medido no coincide con la regla con que se lo mide.

    \item Correspondencia (Riesz): en espacios con producto interno de dimensión finita (como $\mathbb{R}^n$) existe una correspondencia biunívoca entre $V$ y $V^*$; para cada funcional lineal $f \in V^*$ existe un único vector $x \in V$ tal que $f(y) = \langle x, y \rangle$.
\end{itemize}

\subsection{Dimensión y Base Dual}

\begin{definicion}{Base Dual}{def_base_dual}
Sea $V$ un espacio vectorial de dimensión finita $n$, y sea $\mathcal{B} = \{v_1, v_2, \dots, v_n\}$ una base de $V$. 
La \textbf{base dual}, denotada $\mathcal{B}^* = \{f_1, f_2, \dots, f_n\}$, es el conjunto de funcionales lineales en $V^*$ definidos por la condición:
\[
f_i(v_j) = \delta_{ij} = \begin{cases} 
1 & \text{si } i = j \\
0 & \text{si } i \neq j
\end{cases}
\]
\end{definicion}

\begin{teorema}{Existencia y Unicidad de la Base Dual}{teo_base_dual}
Si $V$ es un espacio vectorial de dimensión finita $n$ con base $\mathcal{B}$, entonces el conjunto $\mathcal{B}^*$ definido anteriormente es efectivamente una base de $V^*$ y, además, es la única base que satisface la condición de los deltas de Kronecker (los $\delta_{ij}$).
\end{teorema}

\begin{proof}
La estrategia es probar que $\mathcal{B}^*$ es linealmente independiente, que genera $V^*$ y, por último, su unicidad.

\textit{Independencia.} Supóngase $\sum_{i=1}^n c_i f_i = 0$. Evaluando en $v_j$,
\[
\left( \sum_{i=1}^n c_i f_i \right)(v_j) = \sum_{i=1}^n c_i \delta_{ij} = c_j = 0,
\]
de modo que $c_j = 0$ para todo $j$.

\textit{Generación.} Dado $\phi \in V^*$, sea $\psi = \sum_{i=1}^n \phi(v_i) f_i$. Para cada $v_j$,
\[
\psi(v_j) = \sum_{i=1}^n \phi(v_i) \delta_{ij} = \phi(v_j),
\]
y como $\psi$ y $\phi$ coinciden en la base, $\phi = \psi$. Luego $\mathcal{B}^*$ es base de $V^*$, y $\dim(V^*) = n$.

\textit{Unicidad.} Si $\mathcal{G} = \{g_1, \dots, g_n\}$ satisface también $g_i(v_j) = \delta_{ij}$, entonces $g_i$ y $f_i$ coinciden sobre toda la base $\mathcal{B}$, y como una transformación lineal queda determinada por sus valores en una base, $g_i = f_i$ para todo $i$.
\end{proof}

\subsection{Importancia de la Base Dual}

\begin{proposicion}{Extracción de Coordenadas}{prop_coords_dual}
Sea $v \in V$. Si $v$ se expresa en términos de la base $\mathcal{B}$ como $v = \sum_{i=1}^n x_i v_i$, entonces los coeficientes $x_i$ están dados exactamente por la aplicación de la base dual:
\[
x_i = f_i(v)
\]
\end{proposicion}

\begin{proof}
Aplicando $f_i$ a $v = \sum_{j=1}^n x_j v_j$ y usando la linealidad y la condición $f_i(v_j) = \delta_{ij}$,
\[
f_i(v) = \sum_{j=1}^n x_j f_i(v_j) = \sum_{j=1}^n x_j \delta_{ij} = x_i. \qedhere
\]
\end{proof}

Esta propiedad permite escribir cualquier vector $v$ y cualquier funcional $\phi$ de forma explícita:

\begin{itemize}
    \item Para vectores: $v = \sum_{i=1}^n f_i(v) v_i$
    \item Para funcionales: $\phi = \sum_{i=1}^n \phi(v_i) f_i$
\end{itemize}

\obs
En física, esto subyace a la diferencia entre componentes contravariantes (las coordenadas del vector $x_i$) y componentes covariantes (las componentes del funcional).

\begin{ejemplo}{Caso en $\mathbb{R}^3$}{ej_R3_dual}
Sea la base estándar $e_1, e_2, e_3$ en $\mathbb{R}^3$ (vectores columna). La base dual correspondiente $e^*_1, e^*_2, e^*_3$ son los vectores fila transpuestos:
\[
e^*_1(v) = \begin{bmatrix} 1 & 0 & 0 \end{bmatrix} v, \quad 
e^*_2(v) = \begin{bmatrix} 0 & 1 & 0 \end{bmatrix} v, \quad 
e^*_3(v) = \begin{bmatrix} 0 & 0 & 1 \end{bmatrix} v
\]
Su función es extraer la componente correspondiente del vector sobre el que actúan.
\end{ejemplo}

La relación entre los funcionales del espacio dual y los vectores del espacio (vía producto interno) admite una visualización dinámica en \href{https://www.youtube.com/watch?v=LyGKycYT2v0}{este video explicativo}.

\subsection{Recuperación del Espacio Primal}

\begin{proposicion}{Existencia de Base Pre-dual}{prop_base_pre_dual}
Sea $V$ un $\mathbb{K}$-espacio vectorial de dimensión $n$, y sea $V^*$ su espacio dual. 
Sea $\mathcal{B}^* = \{\phi_1, \dots, \phi_n\}$ una base de $V^*$. 
Entonces existe una única base $\mathcal{B} = \{v_1, \dots, v_n\}$ de $V$ tal que su dual es exactamente $\mathcal{B}^*$ (es decir, $\phi_i(v_j) = \delta_{ij}$).
\end{proposicion}

\begin{proof}
\textit{Existencia.} Sea $\mathcal{W} = \{w_1, \dots, w_n\}$ una base auxiliar de $V$ y $M \in \mathbb{K}^{n \times n}$ la matriz $M_{ij} = \phi_i(w_j)$. La fila $i$ de $M$ es el vector de coordenadas de $\phi_i$ en la base dual $\mathcal{W}^*$ (pues $\phi_i = \sum_j \phi_i(w_j)\, w_j^*$); como $\{\phi_1, \dots, \phi_n\}$ es linealmente independiente en $V^*$ y tomar coordenadas es un isomorfismo, las filas de $M$ son linealmente independientes y $M$ es inversible. Sea $A = (a_{ij}) = M^{-1}$, esto es, $M \cdot A = I_n$, y defínase
\[
v_j = \sum_{k=1}^n a_{kj} w_k, \qquad 1 \le j \le n.
\]
Entonces
\[
\phi_i(v_j) = \sum_{k=1}^n a_{kj} \phi_i(w_k) = \sum_{k=1}^n a_{kj} M_{ik} = (M \cdot A)_{ij} = \delta_{ij}.
\]
Además, $\{v_1, \dots, v_n\}$ es base de $V$: como $\dim(V) = n$, basta la independencia lineal, y si $\sum_j \alpha_j v_j = 0$, aplicando $\phi_i$ se obtiene $\alpha_i = 0$ para todo $i$.

\textit{Unicidad.} Si $\mathcal{B}' = \{u_1, \dots, u_n\}$ es otra base con dual $\mathcal{B}^*$, sus coordenadas en la base $\mathcal{B} = \{v_1, \dots, v_n\}$ son
\[
(u_i)_{\mathcal{B}} = (\phi_1(u_i), \dots, \phi_n(u_i)) = e_i = (v_i)_{\mathcal{B}},
\]
de modo que $u_i = v_i$ para todo $i$.
\end{proof}

\begin{ejemplo}{Interpolación de Lagrange}{ej_lagrange}
Sea $V = \mathbb{R}_2[X]$ (polinomios de grado $\le 2$). Se definen los funcionales $\epsilon_0, \epsilon_1, \epsilon_2 \in V^*$ como las evaluaciones en puntos específicos:
\[
\epsilon_0(P) = P(0), \quad \epsilon_1(P) = P(1), \quad \epsilon_2(P) = P(2).
\]
Puede verificarse que $\{\epsilon_0, \epsilon_1, \epsilon_2\}$ es una base de $V^*$. Por la proposición anterior, debe existir una base de polinomios $\{P_0, P_1, P_2\}$ tal que $\epsilon_i(P_j) = \delta_{ij}$.

Esto implica, por ejemplo, que $P_0$ debe cumplir:
\[
P_0(0)=1, \quad P_0(1)=0, \quad P_0(2)=0.
\]
El polinomio que cumple esto es $P_0(X) = \frac{1}{2}(X-1)(X-2)$.

Para hallar un polinomio $P$ tal que $P(0)=\alpha$, $P(1)=\beta$, $P(2)=\gamma$, se usan las coordenadas en esta base:
\[
P = \alpha P_0 + \beta P_1 + \gamma P_2,
\]
que es la fórmula de interpolación de Lagrange.
\end{ejemplo}

\subsection{Anulador de un Subespacio}

\begin{definicion}{Anulador}{def_anulador}
Sea $V$ un espacio vectorial y $S$ un subespacio de $V$. El \textbf{anulador} de $S$, denotado $S^\circ$, es el conjunto de funcionales que se anulan en todo vector de $S$:
\[
S^\circ = \{ f \in V^* \mid f(s) = 0, \forall s \in S \} = \{ f \in V^* \mid S \subseteq \ker(f) \}.
\]
\end{definicion}

\begin{teorema}{Dimensión del Anulador}{teo_dim_anulador}
Sea $V$ de dimensión finita $n$ y $S$ un subespacio. Entonces:
\[
\dim(S^\circ) = n - \dim(S).
\]
\end{teorema}

\begin{proof}
Sea $\{v_1, \dots, v_r\}$ una base de $S$, extendida a una base $\mathcal{B} = \{v_1, \dots, v_n\}$ de $V$, y sea $\mathcal{B}^* = \{\phi_1, \dots, \phi_n\}$ la base dual. Para $i > r$, $\phi_i$ se anula en $v_1, \dots, v_r$, de modo que $\phi_i \in S^\circ$, y $\{\phi_{r+1}, \dots, \phi_n\}$ es linealmente independiente. Recíprocamente, todo $g \in S^\circ$ se escribe $g = \sum_i \alpha_i \phi_i$, y como $g(v_k) = 0$ para $k \le r$, los coeficientes $\alpha_1, \dots, \alpha_r$ son nulos; luego $g$ depende solo de $\phi_{r+1}, \dots, \phi_n$, que por lo tanto generan $S^\circ$. En consecuencia, $\dim(S^\circ) = n - r = n - \dim(S)$.
\end{proof}

\begin{proposicion}{Recuperación del Subespacio}{prop_recuperacion_sub}
Se puede recuperar el subespacio original a partir de su anulador:
\[
S = \{ x \in V \mid f(x) = 0, \forall f \in S^\circ \}.
\]
\end{proposicion}
\begin{proof}
Sea $T = \{ x \in V \mid f(x) = 0, \ \forall f \in S^\circ \}$ el conjunto del enunciado. Por la definición del anulador, $S \subseteq T$. Para la inclusión recíproca, razónese por el absurdo: si existiera $x \in T$ con $x \notin S$, al completar una base de $S$ con $x$ y tomar la base dual habría un funcional $\phi$ que se anula en $S$ (luego $\phi \in S^\circ$) pero con $\phi(x) = 1$, lo que contradice que $x \in T$. Por lo tanto $T = S$.
\end{proof}

\subsection{Aplicación: El Gradiente y la Dualidad}

\begin{definicion}{El Diferencial (El Covector)}{def_diferencial}
Sea $f: V \to \mathbb{R}$ una función diferenciable en un punto $x \in V$. El \textbf{diferencial} de $f$ en $x$, denotado $df_x$, es el funcional lineal en $V^*$ que mejor aproxima el cambio de la función:
\[
f(x + v) \approx f(x) + df_x(v)
\]
Es decir, $df_x \in V^*$. Toma un vector de desplazamiento $v$ y devuelve un número (el cambio infinitesimal).
\end{definicion}

Ahora bien, $df_x$ es un funcional, no una dirección de movimiento. Para convertirlo en un vector se requiere una estructura adicional: un producto interno.

\begin{teorema}{El Gradiente (Representante de Riesz)}{teo_gradiente_riesz}
Dado un espacio vectorial $V$ con producto interno $\langle \cdot, \cdot \rangle$, por el Teorema de Representación de Riesz, para el funcional $df_x \in V^*$ existe un único vector $g \in V$ tal que:
\[
df_x(v) = \langle g, v \rangle \quad \forall v \in V
\]
A este vector $g$ se lo llama \textbf{gradiente} y se lo denota $\nabla f(x)$.
\end{teorema}

\begin{proof}
Se prueba el caso real de dimensión finita. Sea $\{e_1, \dots, e_n\}$ una base ortonormal de $V$ y defínase $g = \sum_{i=1}^n df_x(e_i)\, e_i$. Para $v = \sum_i v_i e_i$, por la linealidad de $df_x$ y la ortonormalidad de la base,
\[
\langle g, v \rangle = \sum_{i=1}^n df_x(e_i)\, v_i = df_x\!\left( \sum_{i=1}^n v_i e_i \right) = df_x(v).
\]
Para la unicidad, si $\langle g, v \rangle = \langle g', v \rangle$ para todo $v$, entonces $\langle g - g', v \rangle = 0$; tomando $v = g - g'$ queda $\|g - g'\|^2 = 0$, es decir, $g = g'$.
\end{proof}

\obs
El diferencial $df_x$ existe independientemente de la geometría. El gradiente $\nabla f(x)$ depende del producto interno elegido. Si se cambia la forma de medir distancias, el gradiente cambia, aunque la función sea la misma.

\subsection{El Descenso de Gradiente}

\begin{proposicion}{Dirección de Máximo Descenso}{prop_max_descenso}
La dirección $v$ que minimiza $df_x(v)$ sujeto a $\|v\|=1$ es opuesta al gradiente:
\[
v = -\frac{\nabla f(x)}{\|\nabla f(x)\|}
\]
\end{proposicion}

\begin{proof}
Por la definición del gradiente, $df_x(v) = \langle \nabla f(x), v \rangle$. Por la desigualdad de Cauchy-Schwarz, para $\|v\| = 1$,
\[
|\langle \nabla f(x), v \rangle| \le \|\nabla f(x)\| \|v\| = \|\nabla f(x)\|,
\]
y el valor mínimo (más negativo) se alcanza cuando el coseno del ángulo entre $v$ y $\nabla f(x)$ es $-1$, esto es, cuando $v = -\nabla f(x)/\|\nabla f(x)\|$.
\end{proof}

\newpage

\section{Determinantes}

\subsection{Notación}

Sea $\mathcal{M}_{n}(\mathbb{K})$ el espacio de las matrices cuadradas de tamaño $n \times n$ con coeficientes en un cuerpo $\mathbb{K}$.
Dada una matriz $A \in \mathcal{M}_{n}(\mathbb{K})$, sus columnas se denotan como vectores $A^1, A^2, \dots, A^n$ en $\mathbb{K}^n$.

La matriz $A$ se escribe explícitamente como la lista ordenada de estas columnas:
\[
A = (A^1 \mid A^2 \mid \dots \mid A^n)
\]
De esta manera, una función $D: \mathcal{M}_{n}(\mathbb{K}) \to \mathbb{K}$ se puede escribir como:
\[
D(A) = D(A^1 \mid A^2 \mid \dots \mid A^n)
\]

\subsection{Funciones Multilineales Alternadas}

\begin{definicion}{Función Multilineal Alternada}{multilineal_matriz}
Una función $D: \mathcal{M}_{n}(\mathbb{K}) \to \mathbb{K}$ se dice \textbf{multilineal alternada} si cumple las siguientes condiciones:

1. Multilinealidad: la función es lineal respecto a cada columna por separado. Al fijar todas las columnas excepto la $i$-ésima, la función se comporta como una transformación lineal en esa variable.
Para todo $i \in \{1, \dots, n\}$:
\begin{itemize}
    \item Aditividad:
    \[ D(\dots \mid A^i + {A^i}' \mid \dots) = D(\dots \mid A^i \mid \dots) + D(\dots \mid {A^i}' \mid \dots) \]
    \item Homogeneidad:
    \[ D(\dots \mid \lambda A^i \mid \dots) = \lambda D(\dots \mid A^i \mid \dots) \]
\end{itemize}

2. Alternancia: la función se anula si la matriz tiene dos columnas iguales, es decir,
\[
A^i = A^j \text{ para algunos } i \neq j \implies D(A) = 0.
\]
\end{definicion}

\begin{proposicion}{Propiedades Elementales}{determinantes}
Sea $D: \mathcal{M}_{n}(\mathbb{K}) \to \mathbb{K}$ una función multilineal alternada. Entonces:

\begin{enumerate}
    \item
    Si una matriz tiene una columna de ceros ($A^i = 0_{\mathbb{K}^n}$), la función vale 0.
    \[ D(\dots \mid 0 \mid \dots) = 0. \]

    \item
    Si se intercambian dos columnas de lugar, el valor de la función cambia de signo.
    \[ D(\dots \mid A^i \mid \dots \mid A^j \mid \dots) = - D(\dots \mid A^j \mid \dots \mid A^i \mid \dots). \]

    \item
    El valor de la función no cambia si a una columna se le suma un múltiplo de otra columna distinta.
    \[ D(\dots \mid A^i + \alpha A^j \mid \dots \mid A^j \mid \dots) = D(\dots \mid A^i \mid \dots \mid A^j \mid \dots). \]

    \item
    Si una columna es combinación lineal de las demás, la función vale 0.
    \[ A^i = \sum_{j \neq i} \alpha_j A^j \implies D(A) = 0. \]
\end{enumerate}
\end{proposicion}

\begin{proof}
\textit{(1)} Si $A^i = 0_{\mathbb{K}^n}$, escribiendo $0_{\mathbb{K}^n} = 0 \cdot 0_{\mathbb{K}^n}$ y usando la homogeneidad en la columna $i$,
\[
D(\dots \mid 0_{\mathbb{K}^n} \mid \dots) = 0 \cdot D(\dots \mid 0_{\mathbb{K}^n} \mid \dots) = 0.
\]

\textit{(2)} Sean $i \neq j$. Por la alternancia (Definición \ref{def:multilineal_matriz}), una matriz con el vector $A^i + A^j$ en las posiciones $i$ y $j$ tiene dos columnas iguales, luego
\[
0 = D(\dots \mid \underbrace{A^i+A^j}_{\text{pos } i} \mid \dots \mid \underbrace{A^i+A^j}_{\text{pos } j} \mid \dots).
\]
Al expandir por aditividad en las posiciones $i$ y $j$ se obtienen cuatro términos; los dos con columna repetida ($A^i$ con $A^i$, $A^j$ con $A^j$) se anulan, y queda
\[
0 = D(\dots \mid A^i \mid \dots \mid A^j \mid \dots) + D(\dots \mid A^j \mid \dots \mid A^i \mid \dots),
\]
de donde $D(\dots \mid A^i \mid \dots \mid A^j \mid \dots) = - D(\dots \mid A^j \mid \dots \mid A^i \mid \dots)$.

\textit{(3)} Por aditividad en la columna $i$ y luego homogeneidad,
\[
D(\dots \mid A^i + \alpha A^j \mid \dots \mid A^j \mid \dots) = D(\dots \mid A^i \mid \dots \mid A^j \mid \dots) + \alpha\, D(\dots \mid A^j \mid \dots \mid A^j \mid \dots).
\]
El segundo término tiene $A^j$ en las posiciones $i$ y $j$, de modo que por la alternancia vale $0$, y el valor no cambia.

\textit{(4)} Si $A^i = \sum_{k \neq i} c_k A^k$, por multilinealidad
\[
D(A) = \sum_{\substack{k=1 \\ k \neq i}}^{n} c_k \cdot D(\dots \mid \underbrace{A^k}_{\text{pos } i} \mid \dots \mid \underbrace{A^k}_{\text{pos } k} \mid \dots).
\]
Cada término tiene $A^k$ repetido en las posiciones $i$ y $k$, de modo que por la alternancia es nulo y $D(A) = 0$.
\end{proof}

\subsection{Existencia y Unicidad del Determinante}

El siguiente lema, que relaciona una función multilineal alternada definida en $\mathbb{K}^{(n+1) \times (n+1)}$ con otras definidas en $\mathbb{K}^{n \times n}$, es la base del argumento recursivo que permite probar la unicidad de la función determinante.

\begin{lema}\label{lem:insercion}
% Corrección: el enunciado decía f_i(I_n) = α, pero para i ≥ 2 la matriz construida
% sobre I_n es una permutación cíclica de I_{n+1} y el valor correcto es (−1)^{i+1}·α
% (la versión anterior contradecía la fórmula final de la unicidad).
Sea $f: \mathbb{K}^{(n+1) \times (n+1)} \to \mathbb{K}$ una función multilineal alternada tal que $f(I_{n+1}) = \alpha$, y sea $i$ con $1 \le i \le n+1$. Dada $A \in \mathbb{K}^{n \times n}$, se construye $\tilde{A} \in \mathbb{K}^{(n+1) \times (n+1)}$ así:
\begin{itemize}
    \item la columna $1$ de $\tilde{A}$ es el vector canónico $e_i \in \mathbb{K}^{n+1}$;
    \item la fila $i$ de $\tilde{A}$ tiene ceros en todas las posiciones excepto la $(i,1)$;
    \item el resto de los coeficientes corresponden a la matriz $A$ manteniendo su orden relativo (las filas de $A$ ocupan, en orden, las filas de $\tilde{A}$ distintas de la $i$, y sus columnas, las columnas $2, \dots, n+1$).
\end{itemize}
Se define $f_i: \mathbb{K}^{n \times n} \to \mathbb{K}$ por $f_i(A) = f(\tilde{A})$.
Entonces $f_i$ es una función multilineal alternada y $f_i(I_n) = (-1)^{i+1} \alpha$.
\end{lema}

\begin{proof}
Las columnas de $A$ corresponden a las columnas $2, \dots, n+1$ de $\tilde{A}$: sumar dos columnas de $A$ o multiplicar una por un escalar produce exactamente la misma operación sobre las columnas correspondientes de $\tilde{A}$ (en la fila $i$ esas columnas valen $0$, de modo que no se ven afectadas), y dos columnas iguales de $A$ dan dos columnas iguales de $\tilde{A}$. Así, la multilinealidad y la alternancia de $f$ se transfieren a $f_i$.

Para el valor en la identidad: si $A = I_n$, las columnas de $\tilde{I_n}$ son, en orden, $e_i$, $e_1$, $e_2$, \dots, $e_{i-1}$, $e_{i+1}$, \dots, $e_{n+1}$. Llevando la primera columna a la posición $i$ mediante $i-1$ intercambios consecutivos se obtiene $I_{n+1}$, y cada intercambio cambia el signo (Proposición \ref{prop:determinantes}), de modo que
\[
f_i(I_n) = f(\tilde{I_n}) = (-1)^{i-1} f(I_{n+1}) = (-1)^{i+1} \alpha. \qedhere
\]
\end{proof}

\begin{teorema}{Existencia y Unicidad}{existencia_unicidad}
Sea $\alpha \in \mathbb{K}$. Para cada $n \in \mathbb{N}$, existe una única función multilineal alternada $f: \mathbb{K}^{n \times n} \to \mathbb{K}$ tal que $f(I_n) = \alpha$.
\end{teorema}

\begin{proof}
Dada $A \in \mathbb{K}^{(n+1) \times (n+1)}$ se denota $A(i|j) \in \mathbb{K}^{n \times n}$ a la matriz que resulta de suprimir la fila $i$ y la columna $j$ de $A$. La existencia y la unicidad se prueban por inducción en $n$.

\textit{Existencia.} Para $n=1$, $f(x) = \alpha \cdot x$ es multilineal alternada y cumple $f(1) = \alpha$. Supóngase que existe $g: \mathbb{K}^{n \times n} \to \mathbb{K}$ multilineal alternada con $g(I_n) = \alpha$, y defínase $f: \mathbb{K}^{(n+1) \times (n+1)} \to \mathbb{K}$ por
\[
f(A) = \sum_{i=1}^{n+1} (-1)^{i+1} a_{1i} \cdot g(A(1|i)), \qquad A=(a_{jl})_{1 \le j,l \le n+1}.
\]
Se verifica que $f$ es multilineal alternada y que $f(I_{n+1}) = \alpha$.

\textit{(i) Multilinealidad.} Sean $A = (A_1|\dots|A_k|\dots|A_{n+1})$, $A' = (A_1|\dots|A'_k|\dots|A_{n+1})$ y $\tilde{A} = (A_1|\dots|A_k+A'_k|\dots|A_{n+1})$. Entonces
\[
f(\tilde{A}) = \sum_{i \neq k}^{n+1} (-1)^{i+1} a_{1i} g(\tilde{A}(1|i)) + (-1)^{k+1} (a_{1k} + a'_{1k}) g(\tilde{A}(1|k)).
\]
Si $i \neq k$, la columna $k$ de $\tilde{A}(1|i)$ es la suma de las de $A(1|i)$ y $A'(1|i)$, y por la multilinealidad de $g$, $g(\tilde{A}(1|i)) = g(A(1|i)) + g(A'(1|i))$. Si $i = k$, la columna $k$ se elimina y $g(\tilde{A}(1|k)) = g(A(1|k))$. Reagrupando,
\[
f(\tilde{A}) = \left( \sum_{i=1}^{n+1} (-1)^{i+1} a_{1i} g(A(1|i)) \right) + \left( \sum_{i=1}^{n+1} (-1)^{i+1} a'_{1i} g(A'(1|i)) \right) = f(A) + f(A').
\]
El caso del producto por escalar $\lambda A_k$ es análogo.

\textit{(ii) Alternancia.} Sea $A = (A_1|\dots|A_j|\dots|A_k|\dots|A_{n+1})$ con la columna $k$ igual a la $j$ ($k > j$). Entonces
\[
f(A) = \sum_{i \neq k,j}^{n+1} (-1)^{i+1} a_{1i} g(A(1|i)) + (-1)^{j+1} a_{1j} g(A(1|j)) + (-1)^{k+1} a_{1k} g(A(1|k)).
\]
Para cada $i \neq j,k$, la matriz $A(1|i)$ tiene dos columnas iguales, de modo que $g(A(1|i)) = 0$. Además, $A(1|j)$ se obtiene de $A(1|k)$ mediante $k-1-j$ intercambios de columnas, de modo que $g(A(1|k)) = (-1)^{k-1-j} g(A(1|j))$. Luego
\begin{align*}
f(A) &= (-1)^{j+1} a_{1j} g(A(1|j)) + (-1)^{k+1} a_{1j} (-1)^{k-1-j} g(A(1|j)) \\
&= a_{1j} g(A(1|j)) \big( (-1)^{j+1} + (-1)^{2k-j} \big),
\end{align*}
y como $(-1)^{2k-j} = (-1)^{j} = -(-1)^{j+1}$, la suma es $0$.

\textit{(iii) Valor en la identidad.}
\[
f(I_{n+1}) = \sum_{i=1}^{n+1} (-1)^{i+1} (I_{n+1})_{1i} \cdot g(I_{n+1}(1|i)) = (-1)^{2} \cdot 1 \cdot g(I_n) = \alpha.
\]

\textit{Unicidad.} Para $n=1$, una función $f: \mathbb{K}^{1 \times 1} \to \mathbb{K}$ es multilineal alternada si y solo si es lineal, y el valor $f(1)$ la determina: $f(x) = f(1) \, x$. Supóngase la unicidad en $\mathbb{K}^{n \times n}$ para cada valor prescrito en la identidad (para cada $\beta \in \mathbb{K}$ hay una única función multilineal alternada que vale $\beta$ en $I_n$), sea $g$ la que cumple $g(I_n) = \alpha$, y sea $f: \mathbb{K}^{(n+1) \times (n+1)} \to \mathbb{K}$ multilineal alternada con $f(I_{n+1}) = \alpha$. Por linealidad en la primera columna, para $A = (a_{ij})$,
\[
f(A) = \sum_{i=1}^{n+1} a_{i1} f(e_i \mid A^2 \mid \dots \mid A^{n+1}),
\]
con $e_i$ el $i$-ésimo vector canónico. En el término $f(e_i \mid A^2 \mid \dots \mid A^{n+1})$, restando a cada columna $j$ ($2 \le j \le n+1$) la primera multiplicada por $a_{ij}$ —operación que no altera el valor por la Proposición \ref{prop:determinantes}— la matriz resultante tiene el vector $e_i$ en la columna $1$, ceros en la fila $i$ salvo en la posición $(i,1)$, y la submatriz $A(i|1)$ en las filas y columnas restantes: es exactamente la matriz $\tilde{A}$ del Lema \ref{lem:insercion} construida sobre $A(i|1)$. Por lo tanto,
\[
f(e_i \mid A^2 \mid \dots \mid A^{n+1}) = f_i(A(i|1)).
\]
Por el Lema \ref{lem:insercion}, $f_i$ es multilineal alternada con $f_i(I_n) = (-1)^{i+1} \alpha$; la función $(-1)^{i+1} g$ también es multilineal alternada y toma ese mismo valor en $I_n$, de modo que por la hipótesis inductiva $f_i = (-1)^{i+1} g$. Sustituyendo,
\[
f(A) = \sum_{i=1}^{n+1} (-1)^{i+1} a_{i1} g(A(i|1)),
\]
con lo que $f$ queda determinada de forma única.
\end{proof}

\begin{definicion}{Determinante de Orden $n$}{determinante}
La única función multilineal alternada $f: \mathbb{K}^{n \times n} \to \mathbb{K}$ tal que $f(I_n) = 1$ se llama el \textbf{determinante} de orden $n$.
Se la denota usualmente como $\det: \mathbb{K}^{n \times n} \to \mathbb{K}$, y se escribe $\det(A)$ o $|A|$.
\end{definicion}

\begin{corolario}{Desarrollo del Determinante}{calculo_det}
Sea $A \in \mathbb{K}^{n \times n}$. Si para $r \in \mathbb{N}$, $\det: \mathbb{K}^{r \times r} \to \mathbb{K}$ denota la función determinante de orden $r$, entonces valen las siguientes fórmulas recursivas:

\[ \det(A) = \sum_{i=1}^{n} (-1)^{i+1} a_{i1} \det(A(i|1)) = \sum_{j=1}^{n} (-1)^{1+j} a_{1j} \det(A(1|j)) \]
\end{corolario}

\obs Como complemento visual del significado del determinante (el factor de escala de un área o volumen), puede consultarse \href{https://www.youtube.com/watch?v=Ip3X9LOh2dk}{este video explicativo}.

\obs Cualquier otra función multilineal alternada $g$ en $\mathbb{K}^{n \times n}$ es un múltiplo escalar del determinante.
Específicamente, si $g$ es multilineal alternada y $g(I_n) = \alpha$, entonces por la unicidad del Teorema \ref{teo:existencia_unicidad}:
\[ g(A) = \alpha \cdot \det(A) \]

\subsection{Propiedades}

\begin{proposicion}{Determinante de la Transpuesta}{transpuesta}
Sea $A \in \mathbb{K}^{n \times n}$. Entonces:
\[ \det(A) = \det(A^t) \]
\end{proposicion}

\begin{proof}
Por inducción en $n$.

\textit{Caso base ($n=1$).} Si $A \in \mathbb{K}^{1 \times 1}$, entonces $A = A^t$ y la igualdad vale.

\textit{Paso inductivo.} Supóngase la igualdad para $n$ y sea $A \in \mathbb{K}^{(n+1) \times (n+1)}$. Sea $B = A^t$, con $b_{ij} = a_{ji}$. Desarrollando $\det(B)$ por su primera fila,
\[
\det(B) = \sum_{i=1}^{n+1} (-1)^{1+i} b_{1i} \det(B(1|i)).
\]
Como $b_{1i} = a_{i1}$ y $B(1|i) = (A(i|1))^t$, por la hipótesis inductiva $\det((A(i|1))^t) = \det(A(i|1))$, y entonces
\[
\det(A^t) = \sum_{i=1}^{n+1} (-1)^{1+i} a_{i1} \det(A(i|1)),
\]
que es el desarrollo de $\det(A)$ por la primera columna (Corolario \ref{cor:calculo_det}). Luego $\det(A^t) = \det(A)$.
\end{proof}

\obs Como consecuencia, el determinante es una función multilineal alternada por filas. Todo lo demostrado para columnas vale para filas.

\begin{proposicion}{Determinante de Triangular Superior}{triangular}
Sea $A = (a_{ij}) \in \mathbb{K}^{n \times n}$ una matriz triangular superior ($a_{ij} = 0$ si $i > j$). Entonces:
\[ \det(A) = \prod_{i=1}^n a_{ii} \]
\end{proposicion}

\begin{proof}
Por inducción en $n$. Para $n=1$, $\det(A) = a_{11}$. Para el paso inductivo, sea $A \in \mathbb{K}^{(n+1) \times (n+1)}$ triangular superior. Desarrollando por la primera columna,
\[
\det(A) = \sum_{i=1}^{n+1} (-1)^{i+1} a_{i1} \det(A(i|1)),
\]
y como $a_{i1} = 0$ para $i \ge 2$, solo subsiste el término $i=1$:
\[
\det(A) = a_{11} \det(A(1|1)).
\]
La submatriz $A(1|1)$ es triangular superior de tamaño $n \times n$ con diagonal $a_{22}, \dots, a_{n+1,n+1}$, de modo que por la hipótesis inductiva $\det(A(1|1)) = \prod_{j=2}^{n+1} a_{jj}$, y
\[
\det(A) = a_{11} \prod_{j=2}^{n+1} a_{jj} = \prod_{j=1}^{n+1} a_{jj}.
\]
\end{proof}

\begin{proposicion}{Desarrollo General de Laplace}{laplace_general}
Sea $A \in \mathbb{K}^{n \times n}$.
\begin{enumerate}
    \item Por columna $j$: $\det(A) = \sum_{i=1}^n (-1)^{i+j} a_{ij} \det(A(i|j))$
    \item Por fila $i$: $\det(A) = \sum_{j=1}^n (-1)^{i+j} a_{ij} \det(A(i|j))$
\end{enumerate}
\end{proposicion}

\begin{proof}
Se prueba la fórmula por columnas; el caso por filas es análogo vía la transpuesta. Para desarrollar por la columna $j$, se lleva esa columna a la posición $1$ mediante $j-1$ intercambios sucesivos con las columnas anteriores; si $A'$ es la matriz reordenada,
\[ \det(A) = (-1)^{j-1} \det(A'). \]
Desarrollando $A'$ por su primera columna (la columna $j$ original de $A$),
\[ \det(A') = \sum_{i=1}^n (-1)^{i+1} a'_{i1} \det(A'(i|1)), \]
y como $a'_{i1} = a_{ij}$ y $\det(A'(i|1)) = \det(A(i|j))$,
\[
\det(A) = (-1)^{j-1} \sum_{i=1}^n (-1)^{i+1} a_{ij} \det(A(i|j)) = \sum_{i=1}^n (-1)^{i+j} a_{ij} \det(A(i|j)).
\]
\end{proof}

\begin{proposicion}{Determinante del Producto}{producto}
Sean $A, B \in \mathbb{K}^{n \times n}$. Entonces:
\[ \det(A \cdot B) = \det(A) \cdot \det(B) \]
\end{proposicion}

\begin{proof}
La estrategia es fijar $A$ y verificar que $f(X) = \det(A \cdot X)$ es multilineal alternada en las columnas de $X$. Sean $X_1, \dots, X_n$ las columnas de $X$; las de $AX$ son $AX_1, \dots, AX_n$.

\textit{Multilinealidad.} Para la columna $k$, por la distributividad del producto matriz-vector y la linealidad de $\det$ en cada columna,
\begin{align*}
f(\dots \mid X_k + \lambda X'_k \mid \dots) &= \det(\dots \mid AX_k + \lambda AX'_k \mid \dots) \\
&= \det(\dots \mid AX_k \mid \dots) + \lambda \det(\dots \mid AX'_k \mid \dots).
\end{align*}

\textit{Alternancia.} Si $X_i = X_j$ con $i \neq j$, entonces $AX_i = AX_j$, de modo que $f(X) = \det(AX) = 0$.

Por la unicidad del Teorema \ref{teo:existencia_unicidad}, $f(X) = f(I_n) \cdot \det(X)$, y como $f(I_n) = \det(A \cdot I_n) = \det(A)$, tomando $X = B$ se obtiene $\det(A \cdot B) = \det(A) \cdot \det(B)$.
\end{proof}

\begin{proposicion}{Criterio de Inversibilidad}{inversibilidad}
Sea $A \in \mathbb{K}^{n \times n}$. La matriz $A$ es inversible si y solo si su determinante es no nulo:
\[
A \text{ es inversible} \iff \det(A) \neq 0
\]
\end{proposicion}

\begin{proof}
Se prueban las dos implicaciones.

\textit{($\Rightarrow$).} Si $A$ es inversible, existe $B$ con $A \cdot B = I_n$. Aplicando el determinante y la propiedad multiplicativa (Proposición \ref{prop:producto}), $\det(A) \cdot \det(B) = \det(I_n) = 1$, de modo que $\det(A) \neq 0$.

\textit{($\Leftarrow$).} Si $\det(A) \neq 0$, sus columnas son linealmente independientes (Proposición \ref{prop:determinantes}), luego $A$ es inversible.
\end{proof}

\subsection{Adjunta de una Matriz}

\begin{definicion}{Adjunta}{adjunta}
Sea $A=(a_{ij}) \in \mathbb{K}^{n \times n}$. Se llama \textbf{adjunta} de $A$, y se nota $adj(A)$, a la matriz $adj(A) \in \mathbb{K}^{n \times n}$ definida por:
\[
(adj(A))_{ij} = (-1)^{i+j} \det(A(j|i))
\]
\end{definicion}

\begin{proposicion}{Propiedad de la Adjunta}{adjunta}
Sea $A \in \mathbb{K}^{n \times n}$. Entonces:
\[ A \cdot adj(A) = \det(A) \cdot I_n \]
Luego, si $\det(A) \neq 0$, se tiene que:
\[ A^{-1} = \frac{1}{\det(A)} \cdot adj(A) \]
\end{proposicion}

\begin{proof}
Sea $A = (a_{ij})$. El coeficiente $(k,l)$ del producto es
\[
(A \cdot adj(A))_{kl} = \sum_{i=1}^{n} a_{ki} \cdot (adj(A))_{il} = \sum_{i=1}^{n} a_{ki} (-1)^{i+l} \det(A(l|i)).
\]

\textit{Caso $k = l$.} La suma es el desarrollo de $\det(A)$ por la fila $l$, luego $(A \cdot adj(A))_{ll} = \det(A)$.

\textit{Caso $k \neq l$.} La suma es el desarrollo por la fila $l$ del determinante de la matriz $A^{(kl)}$ dada por
\[
(A^{(kl)})_{ij} = \begin{cases} a_{ij} & \text{si } i \neq l, \\ a_{kj} & \text{si } i = l, \end{cases}
\]
que tiene las filas $k$ y $l$ iguales, por lo que su determinante es $0$.

Así $A \cdot adj(A) = \det(A) \cdot I_n$, y si $\det(A) \neq 0$ se despeja $A^{-1} = \frac{1}{\det(A)} \cdot adj(A)$.
\end{proof}

\subsection{Regla de Cramer}

\begin{proposicion}{Regla de Cramer}{cramer}
Sea $A=(a_{ij}) \in \mathbb{K}^{n \times n}$ una matriz inversible, y sea $b \in \mathbb{K}^{n \times 1}$.
Entonces la (única) solución del sistema lineal $A \cdot x = b$ está dada por:
\[
x_i = \frac{\det(A_i)}{\det(A)} \quad \text{para } i=1, \dots, n
\]
donde $\det(A_i)$ representa el determinante de la matriz que se obtiene reemplazando la columna $i$-ésima de $A$ por el vector $b$.
\end{proposicion}

\begin{proof}
Multiplicando $A \cdot x = b$ por $adj(A)$ y usando que $adj(A) \cdot A = A \cdot adj(A) = \det(A) \cdot I_n$ (Proposición \ref{prop:adjunta}),
\[
\det(A) \cdot x = adj(A) \cdot b.
\]
Tomando la coordenada $i$ y la definición de adjunta (Definición \ref{def:adjunta}),
\[
\det(A) \cdot x_i = \sum_{j=1}^{n} (adj(A))_{ij} \cdot b_j = \sum_{j=1}^{n} (-1)^{i+j} b_j \det(A(j|i)),
\]
que es el desarrollo por la columna $i$ del determinante de la matriz $A_i$ obtenida al reemplazar la columna $i$ de $A$ por $b$. Por lo tanto $x_i = \det(A_i)/\det(A)$.
\end{proof}

\subsection{Otra Fórmula para el Determinante}

Dada $A=(a_{ij}) \in \mathbb{K}^{n \times n}$, al escribir $A$ por sus filas $A = (F_1, \dots, F_n)^t$ y usar que el determinante es multilineal alternado por filas, cada fila se descompone en la base canónica $e_1, \dots, e_n$. Desarrollando la multilinealidad se obtiene una suma sobre todos los índices de columnas posibles:
\[
\det(A) = \sum_{1 \le i_1, \dots, i_n \le n} a_{1 i_1} a_{2 i_2} \dots a_{n i_n} \det \begin{pmatrix} e_{i_1} \\ \vdots \\ e_{i_n} \end{pmatrix}
\]
Dado que el determinante se anula si la matriz tiene dos filas iguales, en la suma anterior solo subsisten los términos donde los índices $i_1, \dots, i_n$ son todos distintos, es decir, una reordenación de $\{1, \dots, n\}$.

Para formalizar esto se introduce la siguiente definición.

\begin{definicion}{Permutaciones}{permutaciones}
Sea $I_n = \{1, 2, \dots, n\} \subseteq \mathbb{N}$. Una \textbf{permutación} de $I_n$ es una función biyectiva $\sigma: I_n \to I_n$.
El conjunto de todas las permutaciones de $I_n$ se nota $S_n$.
\end{definicion}

Con esta notación, la suma se restringe a los elementos de $S_n$:
\[
\det(A) = \sum_{\sigma \in S_n} a_{1 \sigma(1)} \dots a_{n \sigma(n)} \det \begin{pmatrix} e_{\sigma(1)} \\ \vdots \\ e_{\sigma(n)} \end{pmatrix}
\]
El determinante de la matriz formada por los vectores canónicos permutados $\begin{pmatrix} e_{\sigma(1)} \\ \vdots \\ e_{\sigma(n)} \end{pmatrix}$ solo depende de la permutación $\sigma$ y siempre vale $1$ o $-1$, ya que esta matriz se obtiene intercambiando filas de la matriz identidad $I_n$.
Esto permite definir el signo de la permutación.

\begin{definicion}{Signo de una Permutación}{signo_perm}
Dada una permutación $\sigma \in S_n$, se define el \textbf{signo} de $\sigma$, denotado $sg(\sigma)$, como el determinante de la matriz de permutación correspondiente:
\[
sg(\sigma) = \det \begin{pmatrix} e_{\sigma(1)} \\ \vdots \\ e_{\sigma(n)} \end{pmatrix}
\]
\end{definicion}

Sustituyendo esto en la expresión anterior se obtiene la fórmula explícita del determinante (fórmula de Leibniz).

\begin{proposicion}{Fórmula de Permutaciones}{leibniz}
Sea $A = (a_{ij}) \in \mathbb{K}^{n \times n}$. Entonces:
\[
\det(A) = \sum_{\sigma \in S_n} sg(\sigma) a_{1 \sigma(1)} a_{2 \sigma(2)} \dots a_{n \sigma(n)}
\]
\end{proposicion}

\obs Para calcular el signo de una permutación $\sigma \in S_n$ basta contar cuántos intercambios de filas se deben realizar en la matriz identidad para conseguir la matriz permutada.
Si el número de intercambios es $r$, el signo de la permutación es $(-1)^r$.

\begin{definicion}{Determinante de un Endomorfismo}{det_endo}
Sea $f: V \to V$ un endomorfismo de un espacio vectorial de dimensión finita. El \textbf{determinante} de $f$, denotado $\det(f)$, se define como el determinante de la matriz asociada a $f$ en cualquier base dada. Este valor es independiente de la base elegida.
\end{definicion}

\obs La independencia de la base se justifica así: dos matrices $A$ y $B$ asociadas a $f$ en bases distintas son semejantes, $B = P^{-1} A P$ (véase el Teorema \ref{teo:cambio_base_semejanza}, más adelante), y por la propiedad multiplicativa (Proposición \ref{prop:producto}),
\[
\det(B) = \det(P^{-1}) \det(A) \det(P) = \det(A),
\]
pues $\det(P^{-1}) \det(P) = \det(I_n) = 1$.

\newpage

\section{Autovalores, Autovectores y Diagonalización}

\subsection{Definiciones}

\begin{definicion}{Matriz Diagonal}{matriz_diagonal}
Una matriz cuadrada $D \in \mathbb{K}^{n \times n}$ se dice \textbf{diagonal} si todos sus elementos fuera de la diagonal principal son nulos.
Formalmente, si sus entradas $d_{ij}$ satisfacen:
\[ d_{ij} = 0 \quad \text{para todo } i \neq j \]
Visualmente, tiene la siguiente estructura:
\[ D = \begin{pmatrix}
\lambda_1 & 0 & \dots & 0 \\
0 & \lambda_2 & \dots & 0 \\
\vdots & \vdots & \ddots & \vdots \\
0 & 0 & \dots & \lambda_n
\end{pmatrix} \]
\end{definicion}

\obs Las matrices diagonales son sencillas de operar:
\begin{itemize}
    \item su determinante es el producto de los elementos de la diagonal;
    \item elevarlas a una potencia $k$ equivale a elevar cada elemento de la diagonal a la $k$.
\end{itemize}

\begin{definicion}{Autovalor y Autovector}{autovalor_autovector}
Sea $V$ un $\mathbb{K}$-espacio vectorial y $f: V \to V$ un endomorfismo.
Un escalar $\alpha \in \mathbb{K}$ es un \textbf{autovalor} de $f$ si existe un vector $v \neq 0$ en $V$ tal que:
\[ f(v) = \alpha v \]
En tal caso, $v$ es un \textbf{autovector} asociado a $\alpha$.
\end{definicion}

\begin{definicion}{Autoespacio}{autoespacio}
Dado un autovalor $\alpha$, el conjunto de todos los autovectores asociados a $\alpha$, junto con el vector nulo, forma un subespacio de $V$, denominado \textbf{autoespacio}, que se denota usualmente por $S_\alpha$:
\[ S_\alpha = \{ v \in V \mid f(v) = \alpha v \} = \ker(f - \alpha \operatorname{id}_V), \]
donde $\operatorname{id}_V$ es la transformación identidad.
\end{definicion}

\obs De la definición se desprenden las siguientes propiedades para $S_\alpha$:
\begin{itemize}
    \item $S_\alpha$ es un subespacio vectorial de $V$ (es el núcleo de la transformación lineal $f - \alpha \operatorname{id}_V$).
    \item El subespacio es invariante bajo $f$: si $x \in S_\alpha$, entonces $f(x) = \alpha x \in S_\alpha$, por ser $S_\alpha$ subespacio.
\end{itemize}

\subsection{Ejemplos}

\begin{ejemplo}{Rotación en el Plano Real}{rotacion_plano}
Sea $V = \mathbb{R}^2$ y $f$ la transformación lineal que rota un vector 90 grados:
\[ f(x, y) = (-y, x) \]
Geométricamente, el vector transformado es ortogonal al original ($f(v) \perp v$). Para que exista un autovalor $\alpha$ haría falta que $f(v)$ fuera paralelo a $v$. Dado que ningún vector no nulo puede ser ortogonal y paralelo a sí mismo simultáneamente, esta transformación no posee autovalores en $\mathbb{R}$.
\end{ejemplo}

\subsection{Cálculo de Autovalores}

\begin{proposicion}{Equivalencias}{equivalencias_autovalor}
Sea $V$ un espacio vectorial de dimensión finita $n$. Son equivalentes:
\begin{enumerate}
    \item[i)] $\alpha$ es autovalor de $f$.
    \item[ii)] El endomorfismo $(f - \alpha \operatorname{id}_V)$ no es inyectivo.
    \item[iii)] $\det(f - \alpha \operatorname{id}_V) = 0$.
\end{enumerate}
\end{proposicion}

\begin{proof}
\textit{(i) $\Leftrightarrow$ (ii).} Se tiene $f(v) = \alpha v \iff (f - \alpha \operatorname{id}_V)(v) = 0$, y que exista $v \neq 0$ con esta propiedad equivale a que el núcleo de $f - \alpha \operatorname{id}_V$ no sea nulo, esto es, a que no sea inyectivo.

\textit{(ii) $\Leftrightarrow$ (iii).} En dimensión finita, un endomorfismo es inyectivo si y solo si es isomorfismo (Corolario \ref{cor:cor_endomorfismos}), y es isomorfismo si y solo si su determinante es no nulo (Proposición \ref{prop:inversibilidad}). Por el contrarrecíproco, $f - \alpha \operatorname{id}_V$ no es inyectivo $\iff \det(f - \alpha \operatorname{id}_V) = 0$.
\end{proof}

\obs El ítem (ii) da una forma de encontrar autovectores dado un autovalor.

\begin{definicion}{Polinomio Característico}{polinomio_caracteristico}
Se define el \textbf{polinomio característico} de $f$, denotado por $P(\alpha)$, como:
\[ P(\alpha) = \det(f - \alpha \operatorname{id}_V) \]
\end{definicion}

\obs El problema de encontrar autovalores se reduce a calcular las raíces de un polinomio.

\subsection{Semejanza de Matrices}

\begin{definicion}{Matrices Semejantes}{matrices_semejantes}
Sean $A, B \in \mathbb{K}^{n \times n}$ dos matrices cuadradas. Se dice que $A$ es \textbf{semejante} a $B$ si existe una matriz invertible $P \in \mathbb{K}^{n \times n}$ tal que:
\[ B = P^{-1} A P \]
\end{definicion}

\obs La semejanza es una relación de equivalencia.

\begin{teorema}{Cambio de Base y Semejanza}{cambio_base_semejanza}
Sean $V$ un espacio vectorial de dimensión finita y $\mathcal{B}, \mathcal{B}'$ dos bases de $V$.
Si $A = [f]_\mathcal{B}$ y $B = [f]_{\mathcal{B}'}$ son las matrices asociadas a un mismo endomorfismo $f: V \to V$ en dichas bases, entonces $A$ y $B$ son semejantes.
\end{teorema}

\begin{proof}
Sea $P$ la matriz de cambio de base de $\mathcal{B}'$ a $\mathcal{B}$ (sus columnas son las coordenadas de los vectores de $\mathcal{B}'$ en la base $\mathcal{B}$), que cumple $[v]_\mathcal{B} = P [v]_{\mathcal{B}'}$ para todo $v \in V$ y, por ser invertible, $[v]_{\mathcal{B}'} = P^{-1} [v]_\mathcal{B}$. Por definición de matriz asociada, $[f(v)]_\mathcal{B} = A [v]_\mathcal{B}$ y $[f(v)]_{\mathcal{B}'} = B [v]_{\mathcal{B}'}$. Entonces
\[
[f(v)]_{\mathcal{B}'} = P^{-1} [f(v)]_\mathcal{B} = P^{-1} A [v]_\mathcal{B} = (P^{-1} A P) [v]_{\mathcal{B}'},
\]
y como esto vale para todo $v$, comparando con $[f(v)]_{\mathcal{B}'} = B [v]_{\mathcal{B}'}$ se obtiene $B = P^{-1} A P$. Luego $A$ y $B$ son semejantes.
\end{proof}

\obs Dado que $A$ y $B$ representan al mismo endomorfismo $f$, comparten todas las propiedades intrínsecas: tienen el mismo determinante, la misma traza y el mismo polinomio característico.

\subsection{Diagonalización}

\begin{definicion}{Matriz Diagonalizable}{matriz_diagonalizable}
Una matriz $A \in \mathbb{K}^{n \times n}$ se dice \textbf{diagonalizable} si es semejante a una matriz diagonal, esto es, si existe una matriz invertible $P$ y una matriz diagonal $D$ tales que:
\[ D = P^{-1} A P \]
\end{definicion}

\begin{definicion}{Endomorfismo Diagonalizable}{endo_diagonalizable}
Sea $V$ un espacio vectorial de dimensión finita. Un endomorfismo $f: V \to V$ se dice \textbf{diagonalizable} si existe una base $\mathcal{B}$ de $V$ tal que la matriz asociada $[f]_\mathcal{B}$ es una matriz diagonal.
\end{definicion}

\obs Si $f$ es diagonalizable, existe un sistema de coordenadas (la base $\mathcal{B}$) en el cual la acción de la transformación lineal se reduce a un escalamiento en cada dirección de los ejes coordenados.

\begin{proposicion}{Equivalencia de Definiciones}{equiv_diag}
Un endomorfismo $f$ es diagonalizable si y solo si su matriz asociada $[f]_\mathcal{B}$ (en cualquier base $\mathcal{B}$) es una matriz diagonalizable.
\end{proposicion}

\begin{proof}
Se prueban las dos implicaciones.

\textit{($\Rightarrow$).} Si $f$ es diagonalizable, existe una base $\mathcal{B}'$ con $[f]_{\mathcal{B}'} = D$ diagonal. Para cualquier otra base $\mathcal{B}$, con $A = [f]_\mathcal{B}$, por el Teorema \ref{teo:cambio_base_semejanza} $A$ y $D$ son semejantes, esto es, $D = P^{-1} A P$ para alguna $P$ invertible; luego $A$ es diagonalizable.

\textit{($\Leftarrow$).} Si $A = [f]_\mathcal{B}$ es diagonalizable, existen $D$ diagonal y $P$ invertible con $D = P^{-1} A P$. Las columnas de $P$, linealmente independientes, son las coordenadas de una base $\mathcal{B}'$ en función de $\mathcal{B}$, y por la fórmula de cambio de base $[f]_{\mathcal{B}'} = P^{-1} A P = D$; luego $f$ es diagonalizable.
\end{proof}

\begin{teorema}{Condición de Diagonalización}{condicion_diag}
Sea $f: V \to V$ un endomorfismo de un espacio vectorial de dimensión finita.
$f$ es diagonalizable si y solo si existe una base de $V$ formada por autovectores de $f$.
\end{teorema}

\begin{proof}
Si $\mathcal{B} = \{v_1, \dots, v_n\}$, la columna $j$ de $[f]_\mathcal{B}$ contiene las coordenadas de $f(v_j)$ en esa base.

\textit{($\Rightarrow$).} Si $f$ es diagonalizable, existe una base $\mathcal{B} = \{v_1, \dots, v_n\}$ con $[f]_\mathcal{B} = D$ diagonal, de modo que la columna $j$ tiene $\alpha_j$ en la posición $j$ y ceros en el resto; esto es, $f(v_j) = \alpha_j v_j$, de modo que cada $v_j$ es autovector y $\mathcal{B}$ es una base de autovectores.

\textit{($\Leftarrow$).} Si $\mathcal{B} = \{v_1, \dots, v_n\}$ es una base con $f(v_j) = \alpha_j v_j$, las coordenadas de $f(v_j)$ en $\mathcal{B}$ son $(0, \dots, \alpha_j, \dots, 0)^T$, y la matriz asociada es diagonal con $D_{jj} = \alpha_j$; luego $f$ es diagonalizable.
\end{proof}

\subsection{Operadores en Espacios con Producto Interno}

En las siguientes definiciones se supone que $(V, \langle \cdot, \cdot \rangle)$ es un espacio vectorial con producto interno de dimensión finita.

\begin{definicion}{Operador Autoadjunto}{operador_autoadjunto}
Sea $f: V \to V$ un endomorfismo. Se dice que $f$ es un \textbf{operador autoadjunto} si para todo par de vectores $u, v \in V$ se cumple la igualdad:
\[ \langle f(u), v \rangle = \langle u, f(v) \rangle. \]
\end{definicion}

\begin{definicion}{Operador Ortogonal}{operador_ortogonal}
Sea $f: V \to V$ un endomorfismo. Se dice que $f$ es un \textbf{operador ortogonal} si preserva el producto interno, esto es, para todo $u, v \in V$:
\[ \langle f(u), f(v) \rangle = \langle u, v \rangle. \]
\end{definicion}

\obs \textit{Intuición Geométrica:}

\begin{enumerate}
    \item[i)] Transformaciones Ortogonales (Preservan distancias):
    La condición $\langle f(v), f(w) \rangle = \langle v, w \rangle$ implica que se preservan las longitudes y los ángulos. Tomando $w = v$:
    \[ \|f(v)\|^2 = \langle f(v), f(v) \rangle = \langle v, v \rangle = \|v\|^2 \]
    Esto significa que la transformación no estira ni encoge los vectores; funciona como un movimiento rígido (como una rotación o una reflexión) que mantiene intacta la geometría del objeto.

    \item[ii)] Transformaciones Autoadjuntas (Simetría):
    La condición \[ \langle f(u), v \rangle = \langle u, f(v) \rangle \] refleja una simetría en la interacción entre los vectores.
    Geométricamente, estos operadores escalan el espacio a lo largo de ejes perpendiculares (sus autovectores), estirando o comprimiendo.
\end{enumerate}

\subsection{Representación Matricial de Operadores Autoadjuntos}

\begin{proposicion}{Matriz Hermítica y Operador Autoadjunto}{matriz_hermitica_iff}
Sea $(V, \langle \cdot, \cdot \rangle)$ un espacio vectorial con producto interno de dimensión finita y $\mathcal{B} = \{e_1, \dots, e_n\}$ una base ortonormal.
Sea $f: V \to V$ un endomorfismo y $A = [f]_\mathcal{B}$ su matriz asociada.
Entonces:
\[ f \text{ es autoadjunto} \iff A \text{ es hermítica, es decir, } A = A^*, \]
donde $A^*$ denota la transpuesta conjugada de $A$ (en el caso real, $A = A^T$: matriz simétrica).
\end{proposicion}

\begin{proof}
Como $\mathcal{B}$ es ortonormal, $a_{ij} = \langle f(e_j), e_i \rangle$.

\textit{($\Rightarrow$).} Si $f$ es autoadjunto, por la simetría conjugada del producto interno,
\[ a_{ij} = \langle f(e_j), e_i \rangle = \langle e_j, f(e_i) \rangle = \overline{\langle f(e_i), e_j \rangle} = \overline{a_{ji}}, \]
luego $A = A^*$.

\textit{($\Leftarrow$).} Si $A$ es hermítica ($a_{ij} = \overline{a_{ji}}$), sean $u = \sum_i x_i e_i$ y $v = \sum_j y_j e_j$. Por linealidad,
\[ \langle f(u), v \rangle = \sum_{i,j} x_i \overline{y_j} \langle f(e_i), e_j \rangle = \sum_{i,j} x_i \overline{y_j}\, a_{ji}, \]
\[ \langle u, f(v) \rangle = \sum_{i,j} x_i \overline{y_j} \langle e_i, f(e_j) \rangle = \sum_{i,j} x_i \overline{y_j}\, \overline{a_{ij}}. \]
Como $a_{ji} = \overline{a_{ij}}$, ambas sumas coinciden, de modo que $\langle f(u), v \rangle = \langle u, f(v) \rangle$.
\end{proof}

\begin{proposicion}{Matriz Unitaria y Operador Ortogonal}{matriz_ortogonal_iff}
Sea $(V, \langle \cdot, \cdot \rangle)$ un espacio vectorial con producto interno de dimensión finita y $\mathcal{B} = \{e_1, \dots, e_n\}$ una base ortonormal.
Sea $f: V \to V$ un endomorfismo y $A = [f]_\mathcal{B}$ su matriz asociada.
Entonces:
\[ f \text{ es ortogonal} \iff A \text{ es unitaria (} A^{-1} = A^* \text{)} \]
\end{proposicion}

\begin{proof}
Sean $C_1, \dots, C_n$ las columnas de $A$; por definición de matriz asociada, $C_j$ son las coordenadas de $f(e_j)$ en $\mathcal{B}$. Como $\mathcal{B}$ es ortonormal, $\langle C_j, C_i \rangle_{\mathbb{K}^n} = \langle f(e_j), f(e_i) \rangle_V$, y como la fila $i$ de $A^*$ es la conjugada traspuesta de $C_i$,
\[ (A^* A)_{ij} = \langle C_j, C_i \rangle_{\mathbb{K}^n} = \langle f(e_j), f(e_i) \rangle_V. \]

\textit{($\Rightarrow$).} Si $f$ es ortogonal, $\langle f(e_j), f(e_i) \rangle = \langle e_j, e_i \rangle = \delta_{ij}$, de modo que $(A^* A)_{ij} = \delta_{ij}$, es decir, $A^* A = I$ y $A^{-1} = A^*$.

\textit{($\Leftarrow$).} Si $A^* A = I$, entonces $\langle f(e_j), f(e_i) \rangle = (A^* A)_{ij} = \delta_{ij}$, de modo que $\{f(e_1), \dots, f(e_n)\}$ es una base ortonormal. Para $u = \sum_k x_k e_k$ y $v = \sum_k y_k e_k$,
\[ \langle f(u), f(v) \rangle = \sum_{k,l} x_k \overline{y_l} \langle f(e_k), f(e_l) \rangle = \sum_k x_k \overline{y_k} = \langle u, v \rangle, \]
de modo que $f$ es ortogonal.
\end{proof}

\subsection{Propiedades Espectrales de Operadores Autoadjuntos}

\begin{teorema}{Autovalores Reales}{autovalores_reales}
Sea $(V, \langle \cdot, \cdot \rangle)$ un espacio vectorial con producto interno y $f$ un endomorfismo autoadjunto.
Entonces, todos los autovalores de $f$ son números reales.
\end{teorema}

\begin{proof}
Sea $\lambda$ un autovalor de $f$ y $v \neq 0$ un autovector asociado, $f(v) = \lambda v$. Por la linealidad en la primera variable, $\langle f(v), v \rangle = \lambda \|v\|^2$; y como $f$ es autoadjunto, $\langle f(v), v \rangle = \langle v, f(v) \rangle = \overline{\lambda} \|v\|^2$. Igualando, $\lambda \|v\|^2 = \overline{\lambda} \|v\|^2$, y como $\|v\| > 0$, se tiene $\lambda = \overline{\lambda}$, esto es, $\lambda \in \mathbb{R}$.
\end{proof}

\begin{lema}\label{lema_invariancia}
Sea $f$ un endomorfismo autoadjunto en $V$ y sea $S \subseteq V$ un subespacio invariante bajo $f$ y no nulo. Entonces:
\begin{enumerate}
    \item[i)] $S^\perp$ es también invariante bajo $f$.
    \item[ii)] $S$ contiene al menos un autovector de $f$.
\end{enumerate}
\end{lema}

\begin{proof}
\textit{(i).} Sea $y \in S^\perp$; se prueba que $f(y) \in S^\perp$, esto es, que $f(y)$ es ortogonal a todo $x \in S$. Como $f$ es autoadjunto, $\langle x, f(y) \rangle = \langle f(x), y \rangle$, y como $S$ es invariante, $f(x) \in S$, de modo que $\langle f(x), y \rangle = 0$ por ser $y \in S^\perp$. Como esto vale para todo $x \in S$, $f(y) \perp S$, es decir, $f(y) \in S^\perp$.

\textit{(ii).} Considérese la restricción $f|_S: S \to S$, bien definida por ser $S$ invariante, y autoadjunta por serlo $f$. Por el Teorema Fundamental del Álgebra, el polinomio característico de $f|_S$ tiene una raíz $\lambda$ (considerando la matriz asociada como matriz compleja), que por el Teorema \ref{teo:autovalores_reales} es real. Entonces existe $v \in S$, $v \neq 0$, con $f|_S(v) = \lambda v$, y como $f|_S(v) = f(v)$, $v$ es un autovector de $f$ en $S$.
\end{proof}

\begin{proposicion}{Ortogonalidad de Autoespacios}{ortogonalidad_autoespacios}
Sea $f$ un endomorfismo autoadjunto. Si $\alpha$ y $\beta$ son dos autovalores distintos de $f$ ($\alpha \neq \beta$), entonces sus autoespacios asociados son ortogonales entre sí.
\end{proposicion}

\begin{proof}
Sean $u \in S_\alpha$ y $v \in S_\beta$, con $f(u) = \alpha u$ y $f(v) = \beta v$. Se calcula $\langle f(u), v \rangle$ de dos formas: por un lado, $\langle f(u), v \rangle = \alpha \langle u, v \rangle$; por otro, como $f$ es autoadjunto y sus autovalores son reales ($\overline{\beta} = \beta$), $\langle f(u), v \rangle = \langle u, f(v) \rangle = \beta \langle u, v \rangle$. Igualando, $(\alpha - \beta) \langle u, v \rangle = 0$, y como $\alpha \neq \beta$, $\langle u, v \rangle = 0$. Luego $S_\alpha \perp S_\beta$.
\end{proof}

\obs Este resultado es útil en la práctica: para diagonalizar ortogonalmente una matriz simétrica basta hallar una base de cada autoespacio por separado, pues los autoespacios de autovalores distintos ya son ortogonales entre sí; solo resta ortonormalizar (con Gram-Schmidt) dentro de cada autoespacio.

\obs La propiedad anterior ($S_\alpha \perp S_\beta$) no vale en general.
    
    Para un endomorfismo cualquiera, si $\alpha \neq \beta$, lo único que puede asegurarse es que los autovectores son linealmente independientes (forman una base, si hay suficientes), pero pueden no formar un ángulo de 90 grados.
    
    La condición de ser autoadjunto es la que fuerza a estos ejes de estiramiento a ser perpendiculares entre sí.

\subsection{Teorema Espectral}

\begin{teorema}{Descomposición Espectral}{descomposicion_espectral}
Sea $f$ un endomorfismo autoadjunto en un espacio $V$ de dimensión finita.
Entonces $V$ es suma directa ortogonal de los autoespacios de $f$:
\[ V = S_{\lambda_1} \oplus^\perp S_{\lambda_2} \oplus^\perp \dots \oplus^\perp S_{\lambda_k} \]
\end{teorema}

\begin{proof}
La estrategia es extraer autoespacios ortogonales de forma sucesiva. Por el Teorema Fundamental del Álgebra y por ser $f$ autoadjunto, existe un autovalor real $\lambda_1$; sea $W_1 = S_{\lambda_1} = \ker(f - \lambda_1 I)$, invariante bajo $f$. Por el Lema \ref{lema_invariancia}, $W_1^\perp$ también es invariante. Si $W_1^\perp = \{0\}$, entonces $V = W_1$ y la descomposición consta de un solo autoespacio.

Si $W_1^\perp \neq \{0\}$, la restricción $f|_{W_1^\perp}$ es autoadjunta, de modo que tiene un autovalor $\lambda_2$ con autovector en $W_1^\perp$. Necesariamente $\lambda_2 \neq \lambda_1$: si fuera $\lambda_2 = \lambda_1$, el autovector estaría en $S_{\lambda_1} \cap S_{\lambda_1}^\perp = \{0\}$. Sea $W_2 = S_{\lambda_2}$; como $\lambda_1 \neq \lambda_2$, $W_1 \perp W_2$. La suma $W_1 \oplus W_2$ es invariante y, por el Lema \ref{lema_invariancia}, su complemento ortogonal también, de modo que el proceso se repite sobre él.

Como $V$ tiene dimensión finita, el proceso termina en un número finito de pasos y se obtiene
\[ V = S_{\lambda_1} \oplus^\perp S_{\lambda_2} \oplus^\perp \dots \oplus^\perp S_{\lambda_k}. \]
\end{proof}

\obs El teorema espectral, en su forma general, se enuncia en espacios de Hilbert.

\subsection{Corolario: Diagonalización Ortogonal}

\begin{corolario}{Diagonalización por Matriz Ortogonal}{diag_ortogonal}
Sea $A \in \mathbb{R}^{n \times n}$ una matriz simétrica.
Entonces existe una matriz ortogonal $P$ y una matriz diagonal $D$ tales que:
\[ P^T A P = D \]
\end{corolario}

\begin{proof}
Por el Teorema \ref{teo:descomposicion_espectral}, $V = S_{\lambda_1} \oplus^\perp \dots \oplus^\perp S_{\lambda_k}$. Tómese una base ortonormal $\mathcal{B}_i$ de cada $S_{\lambda_i}$ y sea
\[ \mathcal{B} = \mathcal{B}_1 \cup \dots \cup \mathcal{B}_k = \{ v_1, \dots, v_n \}, \]
que es una base ortonormal de $V$. La matriz $P$ con columnas $v_1, \dots, v_n$ es ortogonal, luego $P^{-1} = P^T$. El elemento $(i,j)$ de $M = P^T A P$ es
\[ m_{ij} = \langle v_i, A v_j \rangle, \]
interpretando el producto matricial como producto interno de filas y columnas. Como $v_j$ es autovector, $A v_j = \lambda_j v_j$, de modo que
\[ m_{ij} = \lambda_j \langle v_i, v_j \rangle = \lambda_j \delta_{ij}. \]
Por lo tanto $M$ es diagonal con $m_{ii} = \lambda_i$, esto es,
\[ P^T A P = \begin{pmatrix} \lambda_1 & 0 & \dots & 0 \\ 0 & \lambda_2 & \dots & 0 \\ \vdots & \vdots & \ddots & \vdots \\ 0 & 0 & \dots & \lambda_n \end{pmatrix} = D. \]
\end{proof}

\obs El análogo complejo también vale: si $A \in \mathbb{C}^{n \times n}$ es hermítica, existe una matriz unitaria $U$ tal que $U^* A U$ es diagonal (con entradas reales); la demostración es la misma, reemplazando transpuestas por transpuestas conjugadas.

\subsection{Teorema de Cayley-Hamilton}

\begin{teorema}{Cayley-Hamilton}{cayley_hamilton}
Sea $A \in \mathbb{K}^{n \times n}$ una matriz diagonalizable.
Sea $P(\lambda) = \det(A - \lambda I)$ su polinomio característico.
Entonces, la matriz satisface su propia ecuación característica:
\[ P(A) = 0 \]
Donde $0$ representa la matriz nula de tamaño $n \times n$.
\end{teorema}

\begin{proof}
Como $A$ es diagonalizable, $A = C D C^{-1}$ con $D$ diagonal cuyas entradas son los autovalores $\lambda_1, \dots, \lambda_n$. Si $P(t) = a_n t^n + \dots + a_0$, usando $(C D C^{-1})^k = C D^k C^{-1}$,
\[ P(A) = a_n C D^n C^{-1} + \dots + a_0 C I C^{-1} = C\, P(D)\, C^{-1}. \]
Como $D$ es diagonal, $P(D)$ es diagonal con entradas $P(\lambda_i)$, y como los $\lambda_i$ son las raíces de $P$, $P(\lambda_i) = 0$ para todo $i$; luego $P(D) = 0$ y
\[ P(A) = C \cdot 0 \cdot C^{-1} = 0. \]
\end{proof}

\obs El Teorema de Cayley-Hamilton vale para cualquier matriz cuadrada, sea diagonalizable o no. La demostración en el caso general requiere herramientas más avanzadas, como la forma canónica de Jordan o argumentos de continuidad, que no se desarrollan aquí.

\newpage

\section{Aplicación: Optimización de Formas Cuadráticas}

\subsection{Formas Cuadráticas}

Todo polinomio homogéneo de grado 2 en $m$ variables ($x_1, \dots, x_m$) puede escribirse de forma compacta utilizando notación matricial.
Sea $f(x) = \sum_{i,j=1}^m a_{ij} x_i x_j$. Se define la matriz simétrica $A$ de coeficientes y el vector $x = (x_1, \dots, x_m)^T$ tales que:
\[ f(x) = x^T A x = \langle x, Ax \rangle \]
A esta función $f: \mathbb{R}^m \to \mathbb{R}$ se la denomina \textbf{forma cuadrática} asociada a $A$.

\subsection{El Problema de Optimización}

Se plantea el siguiente problema de extremos ligados:
\begin{quote}
    Hallar los valores máximo y mínimo de la función $f(x) = x^T A x$ sujeto a la restricción de que $x$ pertenezca a la esfera unitaria $S^{m-1}$.
\end{quote}
Formalmente:
\[ \text{Optimizar } f(x) = \langle x, Ax \rangle \quad \text{sujeto a } \quad g(x) = \|x\|^2 = 1 \]

\subsection{Resolución por Multiplicadores de Lagrange}

Para encontrar los puntos críticos se utiliza el método de los multiplicadores de Lagrange: se buscan $x$ y un escalar $\lambda$ tales que los gradientes de la función objetivo y de la restricción sean paralelos:
\[ \nabla f(x) = \lambda \nabla g(x) \]

Los gradientes son:
\begin{enumerate}
    \item $\nabla f(x) = \nabla (x^T A x) = 2Ax$ (aquí se usa que $A$ es simétrica),
    \item $\nabla g(x) = \nabla (x^T x) = 2x$.
\end{enumerate}

Sustituyendo en la ecuación de Lagrange:
\[ 2Ax = \lambda (2x) \implies Ax = \lambda x \]

\obs La condición necesaria para que un punto $x$ sea un extremo sobre la esfera es que $x$ sea un autovector de la matriz $A$, y el multiplicador de Lagrange $\lambda$ sea su autovalor asociado.

\subsection{Análisis de los Valores Extremos}

Se evalúa la función objetivo $f$ en estos puntos críticos. Sea $x$ un autovector unitario ($x^T x = 1$) con autovalor asociado $\lambda$.
\[ f(x) = x^T A x = x^T (\lambda x) = \lambda (x^T x) = \lambda \cdot 1 = \lambda \]

Esto conduce al siguiente teorema:

\begin{teorema}{Extremos sobre la Esfera (Teorema del Cociente de Rayleigh)}{extremos_esfera}
Sea $A \in \mathbb{R}^{m \times m}$ una matriz simétrica con autovalores reales ordenados $\lambda_1 \geq \lambda_2 \geq \dots \geq \lambda_m$.
Entonces, sobre la esfera unitaria $\|x\|=1$:
\begin{enumerate}
    \item El valor máximo absoluto de la forma cuadrática es el mayor autovalor, $\lambda_{\max} = \lambda_1$, y se alcanza en el autovector asociado $v_1$.
    \item El valor mínimo absoluto es el menor autovalor, $\lambda_{\min} = \lambda_m$, y se alcanza en el autovector asociado $v_m$.
\end{enumerate}
\end{teorema}

\begin{proof}
La esfera unitaria es compacta y $f$ es continua, de modo que el máximo y el mínimo se alcanzan. Por el análisis anterior (multiplicadores de Lagrange), todo extremo se alcanza en un autovector unitario $x$ de $A$, y allí $f(x) = \lambda$, su autovalor asociado. El máximo es entonces el mayor de los autovalores, $\lambda_1$, alcanzado en $v_1$, y el mínimo es el menor, $\lambda_m$, alcanzado en $v_m$.
\end{proof}

\begin{corolario}{Ortogonalidad de los Extremos}{ortogonalidad_extremos}
Si el autovalor máximo y el mínimo son distintos, entonces los puntos donde se alcanzan el máximo y el mínimo son ortogonales entre sí.
\end{corolario}

\begin{proof}
El máximo se alcanza en $v_1$ (autovector de $\lambda_1$) y el mínimo en $v_m$ (autovector de $\lambda_m$). Como $A$ es simétrica, por la Proposición \ref{prop:ortogonalidad_autoespacios} los autovectores asociados a autovalores distintos son ortogonales, de modo que $v_1 \perp v_m$.
\end{proof}

\newpage

\section{Valores Singulares: Generalización de la Diagonalización}

\subsection{Motivación}

En la sección anterior se trató la diagonalización ($A = P D P^{-1}$), que tiene dos limitaciones:
\begin{enumerate}
    \item Solo se aplica a matrices cuadradas ($A \in \mathbb{K}^{n \times n}$).
    \item No toda matriz cuadrada es diagonalizable (depende de tener una base completa de autovectores).
\end{enumerate}

Si $A \in \mathbb{R}^{m \times n}$ es rectangular y representa una transformación lineal entre espacios de distinta dimensión ($f: \mathbb{R}^n \to \mathbb{R}^m$), la diagonalización clásica no se aplica, porque los vectores de entrada y de salida viven en espacios diferentes. Los valores singulares resuelven esto: en lugar de una única base que diagonalice la transformación, se buscan dos bases ortonormales, una para el dominio y otra para el codominio.

\subsection{Definiciones y Teorema Principal}

Para encontrar estas bases se usan las propiedades de los operadores autoadjuntos ya demostradas. Dada cualquier matriz $A \in \mathbb{R}^{m \times n}$, la matriz $A^T A \in \mathbb{R}^{n \times n}$ es siempre simétrica (y por ende autoadjunta) y semidefinida positiva: $v^T (A^T A) v = \|Av\|^2 \ge 0$ para todo $v$; en particular, sus autovalores son no negativos (si $A^T A v = \lambda v$ con $v \neq 0$, entonces $\lambda \|v\|^2 = \|Av\|^2 \ge 0$).

\begin{definicion}{Valores Singulares}{valores_singulares}
Sea $A \in \mathbb{R}^{m \times n}$. Considérese la matriz simétrica $A^T A$. Por el Teorema Espectral, sus autovalores $\lambda_1, \lambda_2, \dots, \lambda_n$ son reales y no negativos ($\lambda_i \ge 0$).

Se definen los \textbf{valores singulares} de $A$, denotados convencionalmente por $\sigma_i$, como la raíz cuadrada positiva de los autovalores de $A^T A$:
\[ \sigma_i = \sqrt{\lambda_i} \]
Usualmente, se ordenan de mayor a menor: $\sigma_1 \ge \sigma_2 \ge \dots \ge \sigma_r > 0$, donde $r$ es el número de valores singulares distintos de cero.
\end{definicion}

\begin{teorema}{Descomposición en Valores Singulares (SVD)}{svd}
Toda matriz $A \in \mathbb{R}^{m \times n}$ puede factorizarse de la forma:
\[ A = U \Sigma V^T \]
Donde:
\begin{itemize}
    \item $U \in \mathbb{R}^{m \times m}$ es una matriz ortogonal ($U^T U = I$). Sus columnas se llaman \textbf{vectores singulares izquierdos}.
    \item $V \in \mathbb{R}^{n \times n}$ es una matriz ortogonal ($V^T V = I$). Sus columnas se llaman \textbf{vectores singulares derechos}.
    \item $\Sigma \in \mathbb{R}^{m \times n}$ es una matriz diagonal rectangular, cuyos elementos en la diagonal principal son los valores singulares $\sigma_i$ ordenados de forma decreciente, y el resto de los elementos son cero.
\end{itemize}
\end{teorema}

\begin{proof}
La demostración es constructiva: se construyen $V$, $\Sigma$ y $U$.

Considérese $A^T A \in \mathbb{R}^{n \times n}$, simétrica y semidefinida positiva. Por el Corolario \ref{cor:diag_ortogonal} existe una base ortonormal $\{v_1, \dots, v_n\}$ de $\mathbb{R}^n$ formada por autovectores de $A^T A$, con autovalores $\lambda_i \ge 0$. Ordenados de mayor a menor y con $\sigma_i = \sqrt{\lambda_i}$, sea $r$ el rango de $A$, que coincide con la cantidad de autovalores no nulos: $\ker(A^T A) = \ker(A)$ (si $A^T A v = 0$, entonces $\|Av\|^2 = v^T A^T A v = 0$, y la otra inclusión es directa), de modo que
\[ \sigma_1 \ge \dots \ge \sigma_r > 0 \quad \text{y} \quad \sigma_{r+1} = \dots = \sigma_n = 0. \]
La matriz $V = (v_1 \mid \dots \mid v_n)$ es ortogonal por tener columnas ortonormales.

Para $i = 1, \dots, r$ se define $u_i = \frac{1}{\sigma_i} A v_i \in \mathbb{R}^m$. Este conjunto es ortonormal: usando $A^T A v_j = \sigma_j^2 v_j$,
\[ \langle u_i, u_j \rangle = \frac{1}{\sigma_i \sigma_j} v_i^T A^T A v_j = \frac{\sigma_j^2}{\sigma_i \sigma_j}\, v_i^T v_j = \delta_{ij}. \]
Extendiendo $\{u_1, \dots, u_r\}$ a una base de $\mathbb{R}^m$ y ortonormalizando con Gram-Schmidt, se lo completa a una base ortonormal $\{u_1, \dots, u_m\}$ de $\mathbb{R}^m$, y $U = (u_1 \mid \dots \mid u_m)$ es ortogonal. Sea $\Sigma \in \mathbb{R}^{m \times n}$ la matriz con los $\sigma_i$ en la diagonal principal y ceros en el resto.

Como $V$ es ortogonal, $A = U \Sigma V^T$ equivale a $A V = U \Sigma$, que se verifica columna a columna: la columna $i$ de $AV$ es $A v_i$ y la de $U \Sigma$ es $\sigma_i u_i$.
\begin{itemize}
    \item Para $i \le r$, de $u_i = \frac{1}{\sigma_i} A v_i$ se obtiene $A v_i = \sigma_i u_i$.
    \item Para $i > r$, $\sigma_i = 0$ y $\lambda_i = 0$, de modo que $\|A v_i\|^2 = v_i^T A^T A v_i = \lambda_i \|v_i\|^2 = 0$, esto es, $A v_i = 0 = \sigma_i u_i$.
\end{itemize}
Luego $A V = U \Sigma$ y, multiplicando por $V^T$ a derecha, $A = U \Sigma V^T$.
\end{proof}

\obs \textit{Intuición Geométrica:}

El teorema SVD muestra que cualquier transformación lineal matricial $A$ puede descomponerse en tres pasos geométricos secuenciales:
\begin{enumerate}
    \item $V^T$: Una rotación (o reflexión) en el espacio de entrada $\mathbb{R}^n$.
    \item $\Sigma$: Un escalamiento a lo largo de los ejes coordenados (que además cambia la dimensión de $n$ a $m$). Los valores singulares $\sigma_i$ indican exactamente cuánto se «estira» el espacio en cada una de esas direcciones ortogonales.
    \item $U$: Una nueva rotación (o reflexión) en el espacio de salida $\mathbb{R}^m$.
\end{enumerate}
Esto generaliza la diagonalización, donde la rotación de ida y vuelta se hacía con la misma matriz ($P$ y $P^{-1}$), restringida a un único espacio.

\subsection{SVD como Generalización de Autovalores}

\begin{proposicion}{Relación entre SVD y Autovectores}{relacion_svd_autovectores}
Sea $A = U \Sigma V^T$ la Descomposición en Valores Singulares de $A$. Entonces:
\begin{enumerate}
    \item[i)] Las columnas de $V$ son los autovectores ortonormales de $A^T A$.
    \item[ii)] Las columnas de $U$ son los autovectores ortonormales de $A A^T$.
\end{enumerate}
\end{proposicion}

\begin{proof}
\textit{(i).} Sustituyendo la SVD,
\[ A^T A = (U \Sigma V^T)^T (U \Sigma V^T) = V \Sigma^T U^T U \Sigma V^T = V (\Sigma^T \Sigma) V^T, \]
pues $U^T U = I_m$. La matriz $\Sigma^T \Sigma \in \mathbb{R}^{n \times n}$ es diagonal con entradas $\sigma_i^2 = \lambda_i$ (Definición \ref{def:valores_singulares}). Multiplicando por $V$ a derecha y usando $V^T V = I_n$, $(A^T A) V = V (\Sigma^T \Sigma)$, y columna a columna $(A^T A) v_j = \sigma_j^2 v_j = \lambda_j v_j$; luego cada $v_j$ es autovector de $A^T A$.

\textit{(ii).} Análogamente,
\[ A A^T = (U \Sigma V^T)(U \Sigma V^T)^T = U \Sigma V^T V \Sigma^T U^T = U (\Sigma \Sigma^T) U^T, \]
con $\Sigma \Sigma^T$ diagonal de entradas $\sigma_i^2$. El mismo razonamiento muestra que las columnas de $U$ son autovectores de $A A^T$.
\end{proof}

\newpage

\section{Análisis de Componentes Principales}

\subsection{El Problema de la Dimensionalidad}

Considérese un conjunto de datos grande: por ejemplo, una base con $m = 1000$ usuarios, cada uno con $n = 50$ variables distintas (edad, altura, peso, ingresos y tiempo de conexión, entre otras).

Procesar una matriz de $1000 \times 50$ es costoso, y 50 dimensiones no admiten visualización directa. Además, muchas de estas variables suelen estar correlacionadas (por ejemplo, la cantidad de clics y el tiempo en la página tienden a crecer juntos), lo que indica información redundante.

El \textbf{Análisis de Componentes Principales} (PCA) busca reducir esas 50 variables a unas pocas variables artificiales que resuman casi toda la información original. Esto comprime los datos, acelera los algoritmos y permite graficarlos.

\subsection{La Intuición Geométrica}

Considérese un caso más simple, con solo dos variables, altura y peso de varias personas. En un plano 2D, los datos forman una nube de puntos alargada en diagonal, pues a mayor altura suele corresponder mayor peso.

Para resumir la posición de cada persona con un solo número (reducir de 2D a 1D), proyectar sobre el eje $X$ (ignorando el peso) aplasta los puntos y pierde gran parte de la información sobre el tamaño general; lo mismo ocurre al proyectar sobre el eje $Y$.

La idea del PCA es rotar los ejes: trazar una nueva recta que atraviese la nube de puntos a lo largo de su máxima extensión.
\begin{itemize}
    \item A este nuevo eje se lo llama \textbf{primera componente principal} (PC1); sobre él los datos están lo más dispersos posible. En estadística, la dispersión se mide por la varianza, que se interpreta como información.
    \item El segundo eje (PC2) se traza perpendicular al primero y captura la varianza residual.
\end{itemize}

Por ejemplo, si PC1 (interpretable como el tamaño general de la persona) contiene el 95\% de la varianza, puede descartarse PC2 y pasar de 2 a 1 dimensión perdiendo solo un 5\% de la información.

\subsection{Formalización: Varianza y Covarianza}

Para traducir esta intuición a matemáticas, se supone que la matriz de datos $X \in \mathbb{R}^{m \times n}$ está centrada (se ha restado la media a cada columna, de modo que el origen de coordenadas está en el centro de la nube de puntos).

Se busca una dirección, dada por un vector unitario $w \in \mathbb{R}^n$ ($\|w\| = 1$), sobre la cual proyectar los datos de modo que la varianza sea máxima.

\begin{definicion}{Matriz de Covarianzas}{matriz_covarianza}
Dada la matriz de datos centrados $X \in \mathbb{R}^{m \times n}$, se define la \textbf{matriz de covarianzas} muestral $C \in \mathbb{R}^{n \times n}$ como:
\[ C = \frac{1}{m-1} X^T X \]
El elemento $C_{ij}$ indica cómo varían juntas la variable $i$ y la variable $j$. La diagonal principal $C_{ii}$ contiene la varianza individual de cada variable.
\end{definicion}

\obs Es la generalización natural a dimensiones superiores del concepto de varianza de una variable aleatoria escalar.

\obs La matriz $C$ es, por construcción, simétrica y semidefinida positiva, lo que permite usar la teoría de operadores autoadjuntos ya desarrollada.

La proyección de los $m$ datos sobre el vector $w$ es el producto $Xw$. La varianza de estos datos proyectados es una forma cuadrática:
\[ \text{Var}(Xw) = \frac{1}{m-1} (Xw)^T (Xw) = \frac{1}{m-1} w^T X^T X w = w^T C w \]

\subsection{PCA como Optimización de una Forma Cuadrática}

El objetivo de encontrar el eje de mayor dispersión es ahora un problema matemático preciso: maximizar la forma cuadrática $f(w) = w^T C w$ sujeto a que el vector sea unitario ($\|w\|^2 = 1$).

Este es exactamente el problema de optimización sobre la esfera unitaria que se resolvió en la sección anterior.

\begin{proposicion}{Componentes Principales}{pca_autovectores}
Las direcciones ortogonales que maximizan sucesivamente la varianza de los datos (las componentes principales) son exactamente los autovectores de la matriz de covarianzas $C$. La cantidad de información (varianza) capturada por cada componente es exactamente su autovalor asociado.
\end{proposicion}

\begin{proof}
Por el Teorema del Cociente de Rayleigh \ref{teo:extremos_esfera}, el valor máximo de la forma cuadrática $w^T C w$ sobre la esfera unitaria se alcanza en el autovector $v_1$ asociado al mayor autovalor $\lambda_1$ de $C$, que es la primera componente principal. La segunda componente se obtiene maximizando la varianza en las direcciones perpendiculares a $v_1$; por el Corolario \ref{cor:ortogonalidad_extremos}, ese máximo es el segundo mayor autovalor $\lambda_2$, alcanzado en el autovector ortogonal $v_2$. El proceso se repite hasta agotar las $n$ dimensiones.
\end{proof}

\subsection{PCA mediante SVD}

La construcción directa del PCA pasa por formar la matriz $C = \frac{1}{m-1} X^T X$ y hallar sus autovectores. En la práctica computacional moderna, sin embargo, esto rara vez se hace así.

Calcular $X^T X$ cuando $X$ es grande es lento y puede amplificar los errores de redondeo. La Descomposición en Valores Singulares ofrece una alternativa numéricamente estable.

\begin{teorema}{Equivalencia entre SVD y PCA}{equivalencia_svd_pca}
Sea $X \in \mathbb{R}^{m \times n}$ una matriz de datos centrada y sea $X = U \Sigma V^T$ su Descomposición en Valores Singulares. Entonces:
\begin{enumerate}
    \item[i)] Las componentes principales de $X$ son exactamente las columnas de la matriz ortogonal $V$ (los vectores singulares derechos).
    \item[ii)] Las varianzas de los datos sobre cada componente (los autovalores $\lambda_i$ de $C$) están dadas por:
    \[ \lambda_i = \frac{\sigma_i^2}{m-1} \]
    donde $\sigma_i$ son los valores singulares de $X$ (la diagonal de $\Sigma$).
\end{enumerate}
\end{teorema}

\begin{proof}
Sustituyendo la SVD $X = U \Sigma V^T$ en la matriz de covarianzas y usando $U^T U = I_m$,
\[ C = \frac{1}{m-1} (U \Sigma V^T)^T (U \Sigma V^T) = V \left( \frac{\Sigma^T \Sigma}{m-1} \right) V^T. \]
Como $V$ es ortogonal, $V^T = V^{-1}$, de modo que
% Corrección: decía Σ² dentro de la fórmula, pero Σ es rectangular y no puede elevarse
% al cuadrado; la matriz diagonal correcta es Σ^T Σ, como en la línea anterior.
\[ C = V \left( \frac{\Sigma^T \Sigma}{m-1} \right) V^{-1} \]
es la diagonalización de $C$. Por lo tanto, las columnas de $V$ son los autovectores de $C$ (las componentes principales) y los autovalores son $\lambda_i = \frac{\sigma_i^2}{m-1}$.
\end{proof}

\obs Para aplicar PCA basta pasar la matriz $X$ a un algoritmo de SVD: las columnas de $V$ dan la rotación de los ejes y los valores singulares de $\Sigma$ indican cuánta información retiene cada eje, lo que permite descartar los menos relevantes y comprimir los datos.

\subsection{Propiedades de la Transformación}

El Análisis de Componentes Principales reduce dimensiones y, además, transforma la naturaleza estadística de los datos. Al proyectar la matriz $X$ sobre la base de autovectores $V$ se obtiene un nuevo conjunto de datos transformados $Y = XV$.

Esta transformación lineal tiene dos propiedades destacadas:

\begin{proposicion}{Decorrelación de los Datos}{decorrelacion}
Las nuevas variables obtenidas mediante PCA (las componentes principales) están completamente decorrelacionadas entre sí. Es decir, la matriz de covarianzas de los datos transformados $Y$ es una matriz diagonal.
\end{proposicion}

\begin{proof}
La matriz de covarianzas de $Y = XV$ es
\[ C_Y = \frac{1}{m-1} Y^T Y = \frac{1}{m-1} (XV)^T (XV) = V^T \left( \frac{1}{m-1} X^T X \right) V = V^T C_X V. \]
Como las columnas de $V$ son los autovectores de $C_X$, la matriz $V$ diagonaliza $C_X$, de modo que $C_Y = V^T C_X V = D$ es diagonal con los autovalores $\lambda_i$. Al ser diagonal, las covarianzas cruzadas son nulas y los datos quedan decorrelacionados.
\end{proof}

\obs La base de autovectores provista por PCA (o SVD) es óptima para la reconstrucción de la información: para aproximar un conjunto de datos de dimensión $n$ mediante un subespacio de menor dimensión $k$ (con $k < n$), el teorema de Eckart-Young-Mirsky garantiza que proyectar sobre las primeras $k$ componentes principales minimiza el error cuadrático medio de reconstrucción. Ninguna otra base ortogonal retiene más varianza (información) en tan pocas dimensiones.

\obs Dado que esta herramienta matemática es universal, recibe distintos nombres según la disciplina que la aplique:
\begin{itemize}
    \item En Estadística y ciencia de datos, se la conoce como Análisis de Componentes Principales (PCA).
    \item En Procesamiento de Señales e Imágenes, se la conoce como Transformada de Karhunen-Loève (KLT). Se utiliza extensamente para la compresión de imágenes: los píxeles adyacentes de una imagen están altamente correlacionados. La KLT decorrelaciona estos píxeles, concentrando la mayor parte de la energía visual de la imagen en unos pocos coeficientes singulares, permitiendo descartar el resto sin pérdida de calidad perceptible.
\end{itemize}

\newpage

\section{Matrices en la Computadora}

\subsection{Representación en Memoria}

La memoria de una computadora es un arreglo lineal: las celdas se identifican con direcciones consecutivas $0, 1, 2, \dots$ Para guardar una matriz, que es un objeto bidimensional, hay que elegir el orden en que sus entradas se vuelcan sobre esa recta.

\begin{definicion}{Órdenes por Filas y por Columnas}{row_column_major}
Sea $A \in \mathbb{K}^{m \times n}$.
\begin{itemize}
    \item En el orden \textbf{row-major} (por filas), las entradas se almacenan fila por fila: la entrada $a_{ij}$ (con índices desde $1$) ocupa la posición
    \[
    (i-1) \, n + (j-1)
    \]
    dentro del bloque de memoria (contando desde $0$).
    \item En el orden \textbf{column-major} (por columnas), se almacenan columna por columna: $a_{ij}$ ocupa la posición $(j-1) \, m + (i-1)$.
\end{itemize}
\end{definicion}

\obs La elección es una convención de cada lenguaje: C y sus derivados usan row-major; Fortran, MATLAB y R usan column-major. La distinción importa en la práctica: recorrer la matriz en el mismo orden en que está almacenada produce accesos a direcciones consecutivas, que el hardware atiende mucho más rápido que los accesos dispersos (localidad de referencia). Un mismo algoritmo puede cambiar su tiempo real de ejecución en un factor grande solo por el orden de sus bucles, aunque su cantidad de operaciones aritméticas sea idéntica.

Una \textbf{partición por bloques} de $A \in \mathbb{K}^{m \times n}$ se obtiene cortando sus filas en tramos consecutivos de tamaños $m_1, \dots, m_q$ (con $m_1 + \dots + m_q = m$) y sus columnas en tramos consecutivos de tamaños $n_1, \dots, n_r$ (con $n_1 + \dots + n_r = n$): la matriz queda dividida en $q \cdot r$ \textbf{bloques} $A_{IJ} \in \mathbb{K}^{m_I \times n_J}$, donde $A_{IJ}$ es la submatriz formada por las filas del tramo $I$ y las columnas del tramo $J$. Hay dos maneras de materializar una partición:
\begin{enumerate}
    \item \textit{Por copia:} se reserva memoria para cada bloque y se copian sus entradas. Cuesta tiempo proporcional a la cantidad de entradas copiadas, $\Theta(mn)$.
    \item \textit{Por cálculo de índices:} un bloque queda descrito por tres datos —la dirección de su esquina superior izquierda, sus dimensiones y el ancho de la matriz original (el \emph{stride}, que separa en memoria dos filas consecutivas del bloque)—, sin tocar ninguna entrada. Particionar cuesta $\Theta(1)$, y toda modificación del bloque actualiza la matriz original, pues ocupan la misma memoria.
\end{enumerate}
En lo que sigue se supone que las particiones se hacen por cálculo de índices, en tiempo $\Theta(1)$.

La utilidad de los bloques proviene de que el producto de matrices los respeta:

\begin{proposicion}{Producto por Bloques}{producto_bloques}
Sean $A \in \mathbb{K}^{m \times n}$ y $B \in \mathbb{K}^{n \times p}$, particionadas de modo \emph{conformante}: la partición de las columnas de $A$ y la de las filas de $B$ son la misma, $n = n_1 + \dots + n_r$. Sea $C = AB$, particionada según las filas de $A$ ($m = m_1 + \dots + m_q$) y las columnas de $B$ ($p = p_1 + \dots + p_s$). Entonces, para todos $I \in \{1, \dots, q\}$ y $J \in \{1, \dots, s\}$,
\[
C_{IJ} = \sum_{K=1}^{r} A_{IK} \, B_{KJ},
\]
donde cada producto está bien definido: $A_{IK} \in \mathbb{K}^{m_I \times n_K}$ y $B_{KJ} \in \mathbb{K}^{n_K \times p_J}$. Es decir, los bloques se multiplican con la misma regla «fila por columna» que las entradas.
\end{proposicion}

\begin{proof}
Sean $i$ una fila del tramo $I$ y $j$ una columna del tramo $J$, y sea $t_K = n_1 + \dots + n_K$ (con $t_0 = 0$), de modo que el tramo $K$ de la partición común consta de los índices $t_{K-1} + 1, \dots, t_K$. Separando la suma de la definición del producto según esos tramos,
\[
c_{ij} = \sum_{k=1}^{n} a_{ik} \, b_{kj} = \sum_{K=1}^{r} \; \sum_{k = t_{K-1}+1}^{t_K} a_{ik} \, b_{kj}.
\]
Para cada $K$, la suma interna recorre las columnas de $A$ del tramo $K$ contra las filas de $B$ del mismo tramo: es la regla «fila por columna» aplicada a la fila de $A_{IK}$ correspondiente a $i$ y a la columna de $B_{KJ}$ correspondiente a $j$, es decir, la entrada de $A_{IK} B_{KJ}$ en la posición relativa que $(i, j)$ ocupa dentro del bloque $(I, J)$. Sumando sobre $K$ se obtiene la entrada correspondiente de $\sum_{K} A_{IK} B_{KJ}$, como se afirmó.
\end{proof}

El caso que usan los algoritmos de divide y vencerás es el cuadrado con todos los tramos iguales: $A, B \in \mathbb{K}^{n \times n}$ con $n$ par (potencia de $2$ si se particiona recursivamente) y $q = r = s = 2$ con tramos de tamaño $n/2$,
\[
A = \begin{pmatrix} A_{11} & A_{12} \\ A_{21} & A_{22} \end{pmatrix},
\qquad
\begin{aligned}
C_{11} &= A_{11} B_{11} + A_{12} B_{21}, & C_{12} &= A_{11} B_{12} + A_{12} B_{22}, \\
C_{21} &= A_{21} B_{11} + A_{22} B_{21}, & C_{22} &= A_{21} B_{12} + A_{22} B_{22}.
\end{aligned}
\]

\subsection{El Costo de la Multiplicación}

Calcular $C = AB$ directamente desde la definición requiere, para cada una de las $n^2$ entradas $c_{ij}$, una suma de $n$ productos: en total $\Theta(n^3)$ operaciones aritméticas, organizadas en tres bucles anidados sobre $i$, $j$ y $k$.

La Proposición \ref{prop:producto_bloques} sugiere un algoritmo de divide y vencerás: calcular los cuatro bloques de $C$ mediante $8$ multiplicaciones de matrices de $n/2 \times n/2$, más las sumas de bloques. Con particiones por cálculo de índices, su tiempo cumple la recurrencia
\[
T(n) = 8 \, T(n/2) + \Theta(1),
\]
cuya solución es $\Theta(n^3)$ (la recurrencia se resuelve con el método maestro; véanse los apuntes de \emph{Fundamentos} de Algorítmica): dividir en bloques, por sí solo, no mejora el costo asintótico, porque la cantidad de subproblemas ($8$) compensa exactamente la reducción de tamaño.

El \textbf{algoritmo de Strassen} obtiene, mediante una combinación de sumas y restas de bloques, los cuatro bloques de $C$ con solo $7$ multiplicaciones de tamaño $n/2$, a cambio de una cantidad constante de operaciones de bloques de costo $\Theta(n^2)$. Aplicado recursivamente, su tiempo cumple
\[
T(n) = 7 \, T(n/2) + \Theta(n^2),
\]
cuya solución es $\Theta(n^{\lg 7}) = O(n^{2.81})$: asintóticamente mejor que $\Theta(n^3)$. Las siete fórmulas y su verificación figuran en los apuntes de \emph{Fundamentos} de Algorítmica; para estos apuntes bastan su existencia y su costo. En la práctica solo compensa para matrices grandes, por las constantes ocultas y por un uso menos regular de la memoria (véase la observación sobre localidad de la subsección anterior).

\obs El exponente óptimo de la multiplicación de matrices, denotado $\omega$, es un problema abierto: se sabe que $2 \le \omega$ y, mediante refinamientos sucesivos de la idea de Strassen, que $\omega \le 2.372$ (resultados que se citan sin demostración). Se conjetura que $\omega = 2$.

\subsection{Sistemas Lineales y Factorización LU}

El problema computacional central del álgebra lineal es resolver $Ax = b$ con $A \in \mathbb{K}^{n \times n}$ invertible. La Regla de Cramer (Proposición \ref{prop:cramer}) lo resuelve en teoría, pero exige calcular $n+1$ determinantes; el camino eficiente es factorizar $A$ en matrices triangulares, porque los sistemas triangulares se resuelven leyendo las ecuaciones en orden: la primera involucra una sola incógnita, la segunda dos —una ya conocida—, y así sucesivamente.

Una matriz triangular inferior es \textbf{unitriangular} si su diagonal consta de unos.

\begin{proposicion}{Sustitución hacia Adelante y hacia Atrás}{sustituciones}
Sean $L \in \mathbb{K}^{n \times n}$ unitriangular inferior y $U \in \mathbb{K}^{n \times n}$ triangular superior con diagonal no nula.
\begin{enumerate}
    \item El sistema $Ly = b$ tiene solución única, calculable en orden $i = 1, \dots, n$ (\textbf{sustitución hacia adelante}):
    \[
    y_i = b_i - \sum_{j=1}^{i-1} l_{ij} \, y_j.
    \]
    \item El sistema $Ux = y$ tiene solución única, calculable en orden $i = n, \dots, 1$ (\textbf{sustitución hacia atrás}):
    \[
    x_i = \frac{1}{u_{ii}} \Big( y_i - \sum_{j=i+1}^{n} u_{ij} \, x_j \Big).
    \]
\end{enumerate}
Cada sustitución realiza $\Theta(n^2)$ operaciones.
\end{proposicion}

\begin{proof}
La ecuación $i$ de $Ly = b$ es $\sum_{j \le i} l_{ij} y_j = b_i$ con $l_{ii} = 1$; despejando $y_i$ queda la fórmula del enunciado, que solo usa los $y_j$ con $j < i$, ya calculados. La ecuación $i$ de $Ux = y$ es $\sum_{j \ge i} u_{ij} x_j = y_i$ con $u_{ii} \neq 0$; despejando $x_i$ queda la fórmula, que solo usa los $x_j$ con $j > i$. En ambos casos cada valor queda determinado de manera única, y la ecuación $i$ cuesta $O(n)$ operaciones: $\Theta(n^2)$ en total.
\end{proof}

\begin{definicion}{Factorización LU}{def_lu}
Una \textbf{factorización LU} de $A \in \mathbb{K}^{n \times n}$ es una escritura $A = LU$ con $L$ unitriangular inferior y $U$ triangular superior.
\end{definicion}

La factorización es la formulación matricial de la eliminación gaussiana, y conviene construirla examinando \emph{un} paso de eliminación. Supóngase $a_{11} \neq 0$ y particiónese
\[
A = \begin{pmatrix} a_{11} & w^T \\ v & A' \end{pmatrix},
\]
con $v \in \mathbb{K}^{n-1}$ el resto de la primera columna, $w^T$ el resto de la primera fila y $A' \in \mathbb{K}^{(n-1) \times (n-1)}$ el resto de la matriz. Eliminar la primera columna es restar, a cada fila $i \ge 2$, la primera fila multiplicada por el \textbf{multiplicador} $v_i / a_{11}$; con eso la entrada $(i,1)$ queda en $v_i - \frac{v_i}{a_{11}} \, a_{11} = 0$, y al resto de la fila $i$ se le resta $\frac{v_i}{a_{11}} \, w^T$. Apilando esas correcciones para todas las filas, lo que se le resta al bloque $A'$ es el producto exterior $\frac{v \, w^T}{a_{11}}$ (la matriz cuya fila $i$ es $\frac{v_i}{a_{11}} w^T$). El resultado de la eliminación es entonces
\[
\begin{pmatrix} a_{11} & w^T \\ 0 & S \end{pmatrix},
\qquad S = A' - \frac{v \, w^T}{a_{11}} \in \mathbb{K}^{(n-1) \times (n-1)},
\]
donde $S$ —la matriz del sistema que queda sobre las incógnitas $2, \dots, n$ una vez eliminada la primera— se llama \textbf{complemento de Schur} de $A$ respecto del \textbf{pivote} $a_{11}$.

El paso de eliminación se deshace multiplicando por la matriz unitriangular que \emph{vuelve a sumar} lo restado:
\[
\begin{pmatrix} 1 & 0 \\ v / a_{11} & I_{n-1} \end{pmatrix}
\begin{pmatrix} a_{11} & w^T \\ 0 & S \end{pmatrix}
= \begin{pmatrix} a_{11} & w^T \\[2pt] \frac{v}{a_{11}} \, a_{11} + 0 & \frac{v}{a_{11}} \, w^T + S \end{pmatrix}
= \begin{pmatrix} a_{11} & w^T \\ v & A' \end{pmatrix} = A,
\]
donde el bloque inferior derecho recupera $A'$ porque $\frac{v w^T}{a_{11}} + S = A'$ por la definición de $S$. Esto ya es una factorización de $A$ con primer factor unitriangular inferior; el segundo factor todavía no es triangular superior, pero el único bloque que lo impide es $S$, de tamaño $n - 1$: se le aplica el mismo proceso. Si recursivamente $S = L' U'$, con $L'$ unitriangular inferior y $U'$ triangular superior, entonces
\[
\begin{pmatrix} a_{11} & w^T \\ 0 & S \end{pmatrix}
= \begin{pmatrix} 1 & 0 \\ 0 & L' \end{pmatrix}
\begin{pmatrix} a_{11} & w^T \\ 0 & U' \end{pmatrix}
\]
(el producto por bloques da $(a_{11} \mid w^T)$ arriba y $(0 \mid L'U') = (0 \mid S)$ abajo), y los dos factores unitriangulares se combinan en uno:
\[
A = \begin{pmatrix} 1 & 0 \\ v/a_{11} & I_{n-1} \end{pmatrix}
\begin{pmatrix} 1 & 0 \\ 0 & L' \end{pmatrix}
\begin{pmatrix} a_{11} & w^T \\ 0 & U' \end{pmatrix}
= \underbrace{\begin{pmatrix} 1 & 0 \\ v/a_{11} & L' \end{pmatrix}}_{L}
\underbrace{\begin{pmatrix} a_{11} & w^T \\ 0 & U' \end{pmatrix}}_{U}.
\]
La estructura del resultado se lee directamente: $U$ acumula, fila por fila, los pivotes con sus filas tal como quedan tras cada eliminación, y $L$ guarda, columna por columna, los multiplicadores usados para producirlas ($v / a_{11}$ en la primera columna; dentro de $L'$, los de los pasos siguientes). La eliminación gaussiana y la factorización LU son, así, el mismo cálculo. El costo está dominado por el cálculo de $S$, que cuesta $\Theta((n-1)^2)$ operaciones; el tiempo total cumple $T(n) = T(n-1) + \Theta(n^2) = \Theta(n^3)$.

\begin{teorema}{Unicidad de la Factorización LU}{unicidad_lu}
Si $A$ es invertible, su factorización LU, cuando existe, es única.
\end{teorema}

\begin{proof}
Sean $A = L_1 U_1 = L_2 U_2$ dos factorizaciones. Como $\det(A) = \det(L_1) \det(U_1) \neq 0$ y $\det(L_1) = 1$ (Proposiciones \ref{prop:producto}, \ref{prop:transpuesta} y \ref{prop:triangular}), las matrices $U_1$ y $U_2$ son invertibles, y las unitriangulares $L_1, L_2$ lo son siempre (determinante $1$). Multiplicando la igualdad $L_1 U_1 = L_2 U_2$ a izquierda por $L_2^{-1}$ y a derecha por $U_1^{-1}$,
\[
L_2^{-1} L_1 = U_2 \, U_1^{-1}.
\]
El lado izquierdo es unitriangular inferior y el derecho es triangular superior (el producto de dos triangulares del mismo tipo es triangular de ese tipo, y la inversa de una matriz unitriangular o triangular invertible es del mismo tipo; ambas afirmaciones se verifican de forma directa por inducción). Una matriz que es a la vez unitriangular inferior y triangular superior es la identidad, de modo que $L_2^{-1} L_1 = I_n = U_2 U_1^{-1}$, es decir, $L_1 = L_2$ y $U_1 = U_2$.
\end{proof}

No toda matriz invertible admite factorización LU: el pivote puede anularse. Por ejemplo,
\[
A = \begin{pmatrix} 0 & 1 \\ 1 & 1 \end{pmatrix}
\]
es invertible, pero si fuera $A = LU$, la entrada $(1,1)$ del producto sería $u_{11} = 0$, y entonces la entrada $(2,1)$ sería $l_{21} u_{11} = 0 \neq 1$. La solución es permutar filas antes de eliminar; por precisión numérica, en la práctica se elige como pivote la entrada de mayor valor absoluto de la columna, incluso cuando el pivote natural no es nulo (dividir por pivotes pequeños amplifica los errores de redondeo).

\begin{definicion}{Factorización LUP}{def_lup}
Una \textbf{factorización LUP} de $A \in \mathbb{K}^{n \times n}$ es una terna $(L, U, P)$ con $L$ unitriangular inferior, $U$ triangular superior y $P$ una matriz de permutación (la identidad con sus filas reordenadas), tal que
\[
P A = L U.
\]
\end{definicion}

La existencia se prueba con la misma eliminación de antes, asegurada por intercambios: si el pivote natural es nulo, se sube a la primera posición —permutando dos filas— alguna entrada no nula de la primera columna, que existe porque $A$ es invertible, y se elimina normalmente. El único punto delicado es el registro de los multiplicadores: los niveles más profundos de la recursión permutan filas del complemento de Schur $S$, y los multiplicadores ya guardados, que acompañan a esas filas, deben permutarse de la misma manera; en la demostración que sigue, ese es exactamente el papel del bloque $P'v/a$ dentro de $L$.

\begin{teorema}{Existencia de la Factorización LUP}{existencia_lup}
Toda matriz invertible $A \in \mathbb{K}^{n \times n}$ admite una factorización LUP.
\end{teorema}

\begin{proof}
Por inducción en $n$. Para $n = 1$, $A = (a)$ con $a \neq 0$, y basta tomar $P = L = (1)$, $U = (a)$.

Sea $n > 1$. Como $A$ es invertible, su primera columna no es nula (si lo fuera, las columnas no serían linealmente independientes, contra la Proposición \ref{prop:prop_matrices_invertibles}); sea $a_{k1} \neq 0$ y sea $P_1$ la permutación que intercambia las filas $1$ y $k$. Entonces
\[
P_1 A = \begin{pmatrix} a & w^T \\ v & A' \end{pmatrix} \quad \text{con } a = a_{k1} \neq 0,
\qquad
P_1 A = \begin{pmatrix} 1 & 0 \\ v/a & I_{n-1} \end{pmatrix} \begin{pmatrix} a & w^T \\ 0 & S \end{pmatrix},
\]
con $S = A' - v w^T / a$ el complemento de Schur. Tomando determinantes en la segunda igualdad y desarrollando el segundo factor por su primera columna (Corolario \ref{cor:calculo_det}), $\det(P_1 A) = a \cdot \det(S)$; como $\det(P_1 A) = \pm \det(A) \neq 0$, la matriz $S \in \mathbb{K}^{(n-1) \times (n-1)}$ es invertible. Por hipótesis inductiva existe una factorización $P' S = L' U'$. Defínanse
\[
P = \begin{pmatrix} 1 & 0 \\ 0 & P' \end{pmatrix} P_1, \qquad
L = \begin{pmatrix} 1 & 0 \\ P'v/a & L' \end{pmatrix}, \qquad
U = \begin{pmatrix} a & w^T \\ 0 & U' \end{pmatrix}.
\]
La matriz $P$ es de permutación (producto de permutaciones), $L$ es unitriangular inferior y $U$ es triangular superior. Multiplicando por bloques,
\[
P A = \begin{pmatrix} 1 & 0 \\ 0 & P' \end{pmatrix} \begin{pmatrix} a & w^T \\ v & A' \end{pmatrix}
= \begin{pmatrix} a & w^T \\ P'v & P'A' \end{pmatrix},
\qquad
L U = \begin{pmatrix} a & w^T \\ P'v & \frac{P'v \, w^T}{a} + L'U' \end{pmatrix},
\]
y ambas coinciden: por la hipótesis inductiva, $L'U' = P'S = P' \big( A' - v w^T/a \big)$, de modo que $\frac{P'v \, w^T}{a} + L'U' = P'A'$.
\end{proof}

Con la factorización $PA = LU$, resolver $Ax = b$ se reduce a dos sustituciones. Como $P$ es invertible (solo reordena las ecuaciones), $Ax = b$ equivale a $PAx = Pb$, esto es, a $L(Ux) = Pb$; llamando $y = Ux$, se resuelve primero $Ly = Pb$ por sustitución hacia adelante y luego $Ux = y$ por sustitución hacia atrás, todo en tiempo $\Theta(n^2)$ (Proposición \ref{prop:sustituciones}). Factorizar cuesta $\Theta(n^3)$, pero se hace una sola vez: resolver después el sistema para muchos vectores $b$ distintos cuesta solo $\Theta(n^2)$ por vector.

\begin{ejemplo}{Resolución de un Sistema por LU}{ej_lu}
Sean
\[
A = \begin{pmatrix} 1 & 2 & 0 \\ 3 & 4 & 4 \\ 5 & 6 & 3 \end{pmatrix}, \qquad b = \begin{pmatrix} 3 \\ 7 \\ 8 \end{pmatrix}.
\]
\textit{Factorización.} El primer pivote es $a_{11} = 1$, con $v = (3, 5)^T$ y $w^T = (2, 0)$:
\[
S = \begin{pmatrix} 4 & 4 \\ 6 & 3 \end{pmatrix} - \begin{pmatrix} 3 \\ 5 \end{pmatrix} \begin{pmatrix} 2 & 0 \end{pmatrix}
= \begin{pmatrix} 4 - 6 & 4 \\ 6 - 10 & 3 \end{pmatrix} = \begin{pmatrix} -2 & 4 \\ -4 & 3 \end{pmatrix}.
\]
El segundo pivote es $-2$, con multiplicador $\frac{-4}{-2} = 2$ y complemento de Schur $3 - 2 \cdot 4 = -5$. Reuniendo los multiplicadores en $L$ y los pivotes con sus filas en $U$:
\[
L = \begin{pmatrix} 1 & 0 & 0 \\ 3 & 1 & 0 \\ 5 & 2 & 1 \end{pmatrix}, \qquad
U = \begin{pmatrix} 1 & 2 & 0 \\ 0 & -2 & 4 \\ 0 & 0 & -5 \end{pmatrix}
\]
(la primera columna de $L$ guarda los multiplicadores $3$ y $5$ del primer pivote, y la segunda, el multiplicador $2$ del segundo; no hizo falta permutar, $P = I_3$, y se comprueba directamente que $LU = A$).

\textit{Sustitución hacia adelante} ($Ly = b$): $y_1 = 3$; $y_2 = 7 - 3 \cdot 3 = -2$; $y_3 = 8 - 5 \cdot 3 - 2 \cdot (-2) = -3$.

\textit{Sustitución hacia atrás} ($Ux = y$): $x_3 = \frac{-3}{-5} = \frac{3}{5}$; $x_2 = \frac{1}{-2} \big( {-2} - 4 \cdot \frac{3}{5} \big) = \frac{11}{5}$; $x_1 = 3 - 2 \cdot \frac{11}{5} = -\frac{7}{5}$.

\textit{Verificación:} $A x = \big( {-\frac{7}{5}} + \frac{22}{5}, \ {-\frac{21}{5}} + \frac{44}{5} + \frac{12}{5}, \ {-\frac{35}{5}} + \frac{66}{5} + \frac{9}{5} \big)^T = (3, 7, 8)^T = b$.
\end{ejemplo}

\obs Dos subproductos de la factorización LUP:
\begin{itemize}
    \item \textit{Determinante:} de $\det(P) \det(A) = \det(L) \det(U)$ se despeja $\det(A) = (-1)^s \prod_{i=1}^n u_{ii}$, donde $s$ es la cantidad de intercambios de filas que realiza $P$. Es la forma práctica de calcular determinantes, en $\Theta(n^3)$, en lugar de los $n!$ términos de la fórmula de permutaciones (Proposición \ref{prop:leibniz}).
    \item \textit{Inversa:} las columnas de $A^{-1}$ son las soluciones de los $n$ sistemas $A x_i = e_i$; con la factorización ya calculada, cuestan $\Theta(n^2)$ cada uno: $\Theta(n^3)$ en total. En la práctica casi nunca se necesita $A^{-1}$ explícita: resolver el sistema es menos costoso y numéricamente más estable.
\end{itemize}

\obs La factorización puede calcularse \emph{in situ}: como $L$ es unitriangular (su diagonal de unos no necesita guardarse) y $U$ es triangular superior, ambas caben exactamente en el arreglo de $A$, sobrescribiendo los multiplicadores debajo de la diagonal y $U$ desde la diagonal hacia arriba. Es un ejemplo de cómo la representación en memoria (primera subsección) condiciona el diseño del algoritmo.

\subsection{Factorización QR y Mínimos Cuadrados}

La factorización QR es la versión matricial del proceso de Gram-Schmidt (Teorema \ref{teo:gram_schmidt}). Aplicado a las columnas $a_1, \dots, a_n$ de $A$, Gram-Schmidt produce vectores ortonormales $q_1, \dots, q_n$ que conservan los generados:
\[
\operatorname{span}\{a_1, \dots, a_k\} = \operatorname{span}\{q_1, \dots, q_k\} \quad \text{para todo } k.
\]
Esa conservación es el punto central de lo que sigue: dice que cada $a_k$ es combinación lineal de los primeros $k$ vectores ortonormales —no necesita los posteriores—, de modo que la matriz de coeficientes que expresa las columnas de $A$ en términos de $q_1, \dots, q_n$ es \emph{triangular superior}. Esa observación, escrita como producto de matrices, es precisamente la factorización QR.

\begin{definicion}{Factorización QR}{def_qr}
Sea $A \in \mathbb{R}^{m \times n}$ con columnas linealmente independientes ($m \ge n$). Una \textbf{factorización QR} de $A$ es una escritura $A = QR$ con $Q \in \mathbb{R}^{m \times n}$ de columnas ortonormales y $R \in \mathbb{R}^{n \times n}$ triangular superior con diagonal positiva.
\end{definicion}

\obs Toda matriz $Q$ de columnas ortonormales satisface dos identidades que se usarán repetidamente. Primera: $Q^T Q = I_n$ \emph{aunque $Q$ no sea cuadrada}, pues la entrada $(i, j)$ de $Q^T Q$ es el producto de la fila $i$ de $Q^T$ (es decir, $q_i$) contra la columna $j$ de $Q$: $\langle q_i, q_j \rangle = \delta_{ij}$. Segunda: cuando $m > n$, $Q Q^T \in \mathbb{R}^{m \times m}$ \emph{no} es la identidad, sino la matriz de la proyección ortogonal sobre el espacio columna de $Q$: para $b \in \mathbb{R}^m$,
\[
Q Q^T b = \sum_{i=1}^{n} \langle b, q_i \rangle \, q_i = P_{\operatorname{Col}(Q)}(b),
\]
por la fórmula de la proyección en base ortonormal (Corolario \ref{cor:caso_ortonormal}).

\begin{teorema}{Existencia y Unicidad de la QR}{existencia_qr}
Toda matriz $A \in \mathbb{R}^{m \times n}$ con columnas linealmente independientes admite una única factorización QR.
\end{teorema}

\begin{proof}
\textit{Existencia.} Sean $a_1, \dots, a_n$ las columnas de $A$, linealmente independientes por hipótesis, y aplíquese Gram-Schmidt, normalizando:
\[
v_k = a_k - \sum_{j=1}^{k-1} \frac{\langle a_k, v_j \rangle}{\|v_j\|^2} \, v_j, \qquad q_k = \frac{v_k}{\|v_k\|}
\]
(los $v_k$ son no nulos por la independencia lineal, como en el Teorema \ref{teo:gram_schmidt}, de modo que la normalización está bien definida). Cada término de la suma es la proyección de $a_k$ sobre la dirección de $v_j$, que escrita con el vector normalizado toma la forma
\[
\frac{\langle a_k, v_j \rangle}{\|v_j\|^2} \, v_j
= \Big\langle a_k, \tfrac{v_j}{\|v_j\|} \Big\rangle \, \tfrac{v_j}{\|v_j\|}
= \langle a_k, q_j \rangle \, q_j.
\]
Despejando $a_k$ de la definición de $v_k$ y usando $v_k = \|v_k\| \, q_k$,
\[
a_k = \underbrace{\langle a_k, q_1 \rangle}_{=: \, r_{1k}} \, q_1 + \dots + \underbrace{\langle a_k, q_{k-1} \rangle}_{=: \, r_{k-1,k}} \, q_{k-1} + \underbrace{\|v_k\|}_{=: \, r_{kk}} \, q_k \, :
\]
cada columna de $A$ es combinación lineal de $q_1, \dots, q_k$ \emph{solamente}. Defínase entonces $R \in \mathbb{R}^{n \times n}$ con las entradas $r_{jk}$ indicadas y $r_{jk} = 0$ para $j > k$: es triangular superior, con diagonal $r_{kk} = \|v_k\| > 0$. La igualdad de arriba dice que $a_k = Q \, r_k$, donde $r_k$ es la columna $k$ de $R$ —el producto $Q r_k$ es la combinación lineal de las columnas de $Q$ con los coeficientes de $r_k$ (Proposición \ref{prop:prop_img_cols}), y los coeficientes desde el $k+1$ en adelante son nulos—. Apilando las $n$ igualdades columna a columna, $A = QR$.

\textit{Unicidad.} La estructura triangular fuerza las columnas de $Q$ y de $R$ una por una. Sea $A = QR$ una factorización \emph{cualquiera} en las condiciones del enunciado; se prueba por inducción en $k$ que la columna $q_k$ y los coeficientes $r_{1k}, \dots, r_{kk}$ quedan completamente determinados por $A$, con lo cual dos factorizaciones cualesquiera coinciden. El punto de partida es la lectura columna a columna de $A = QR$, la misma de la existencia:
\[
a_k = r_{1k} \, q_1 + \dots + r_{k-1,k} \, q_{k-1} + r_{kk} \, q_k
\]
(los coeficientes con índice mayor que $k$ son nulos porque $R$ es triangular superior).

\textit{Caso $k = 1$:} la igualdad es $a_1 = r_{11} \, q_1$. Tomando normas, $\|a_1\| = r_{11} \|q_1\| = r_{11}$ (pues $\|q_1\| = 1$ y $r_{11} > 0$): forzosamente $r_{11} = \|a_1\|$ y $q_1 = a_1 / \|a_1\|$.

\textit{Paso inductivo:} supóngase que $q_1, \dots, q_{k-1}$ y las primeras $k-1$ columnas de $R$ ya quedaron determinadas. Tomando en la igualdad de arriba el producto interno contra $q_i$ con $i < k$, la ortonormalidad de las columnas de $Q$ anula todos los términos salvo uno:
\[
\langle a_k, q_i \rangle = \sum_{j=1}^{k} r_{jk} \, \langle q_j, q_i \rangle = r_{ik},
\]
de modo que los coeficientes $r_{ik}$ con $i < k$ están determinados ($q_i$ lo está, por hipótesis inductiva). Conocidos esos coeficientes, en la igualdad solo queda una incógnita, y se la despeja:
\[
r_{kk} \, q_k = a_k - \sum_{j=1}^{k-1} r_{jk} \, q_j =: z_k,
\]
donde $z_k$ es un vector ya determinado, y no nulo, pues $\|z_k\| = r_{kk} \|q_k\| = r_{kk} > 0$. Tomando normas, $r_{kk} = \|z_k\|$, y dividiendo, $q_k = z_k / \|z_k\|$: toda la columna $k$, de $Q$ y de $R$, queda determinada. La inducción concluye.
\end{proof}

\obs En la demostración de la unicidad, el vector despejado $z_k$ es exactamente el $v_k$ de Gram-Schmidt, y las fórmulas obtenidas ($r_{ik} = \langle a_k, q_i \rangle$, $r_{kk} = \|z_k\|$) coinciden con las de la existencia: toda factorización QR reproduce necesariamente el proceso de Gram-Schmidt.

\obs Las columnas de $Q$ generan lo mismo que las de $A$: cada $a_k$ es combinación de $q_1, \dots, q_k$ (existencia) y, despejando con $R^{-1}$, cada $q_k$ es combinación de $a_1, \dots, a_k$. Por lo tanto $\operatorname{Col}(Q) = \operatorname{Col}(A)$, y la observación anterior dice que $Q Q^T$ es la proyección ortogonal sobre $\operatorname{Col}(A)$: la factorización QR reúne, en forma matricial, una base ortonormal del espacio columna y la matriz de coeficientes $R$ que expresa las columnas originales en esa base.

\begin{ejemplo}{Una Factorización QR}{ej_qr}
Sea
\[
A = \begin{pmatrix} 1 & 1 \\ 1 & 0 \\ 0 & 1 \end{pmatrix},
\]
con columnas $a_1 = (1,1,0)^T$ y $a_2 = (1,0,1)^T$. Gram-Schmidt da $v_1 = a_1$, con $\|v_1\| = \sqrt{2}$, y
\[
v_2 = a_2 - \frac{\langle a_2, v_1 \rangle}{\|v_1\|^2} v_1 = (1,0,1)^T - \frac{1}{2}(1,1,0)^T = \Big( \frac{1}{2}, -\frac{1}{2}, 1 \Big)^T,
\qquad \|v_2\| = \frac{\sqrt{6}}{2}.
\]
Con $r_{12} = \langle a_2, q_1 \rangle = \frac{1}{\sqrt{2}}$, la factorización es
\[
Q = \begin{pmatrix} \frac{1}{\sqrt{2}} & \frac{1}{\sqrt{6}} \\ \frac{1}{\sqrt{2}} & -\frac{1}{\sqrt{6}} \\ 0 & \frac{2}{\sqrt{6}} \end{pmatrix}, \qquad
R = \begin{pmatrix} \sqrt{2} & \frac{1}{\sqrt{2}} \\ 0 & \frac{\sqrt{6}}{2} \end{pmatrix},
\]
y se comprueba columna a columna que $QR = A$: la primera es $\sqrt{2} \, q_1 = a_1$ y la segunda es $\frac{1}{\sqrt{2}} q_1 + \frac{\sqrt{6}}{2} q_2 = \big( \frac{1}{2}, \frac{1}{2}, 0 \big)^T + \big( \frac{1}{2}, -\frac{1}{2}, 1 \big)^T = a_2$.
\end{ejemplo}

Cuando hay más ecuaciones que incógnitas ($m > n$) —por ejemplo, al ajustar un modelo con pocos parámetros a muchos datos con ruido— el sistema $Ax = b$ es en general incompatible: los vectores de la forma $Ax$ recorren solo el subespacio $\operatorname{Col}(A)$, de dimensión a lo sumo $n$, y el $b$ medido no cae en él. Lo mejor que puede hacerse es minimizar la norma del \textbf{residuo} $b - Ax$, y la teoría de la proyección ortogonal ya resolvió ese problema: el punto de $\operatorname{Col}(A)$ más cercano a $b$ es único y es la proyección $P_{\operatorname{Col}(A)}(b)$ (Teoremas \ref{teo:mejor_aprox} y \ref{teo:descomposicion_ortogonal}), caracterizada porque el residuo queda \emph{perpendicular} al subespacio. Imponer esa perpendicularidad contra cada columna de $A$ produce un sistema de $n$ ecuaciones lineales: las \emph{ecuaciones normales}.

\begin{teorema}{Ecuaciones Normales}{ecuaciones_normales}
Sean $A \in \mathbb{R}^{m \times n}$ y $b \in \mathbb{R}^m$. Un vector $\hat{x} \in \mathbb{R}^n$ minimiza $\|Ax - b\|$ si y solo si
\[
A^T A \, \hat{x} = A^T b.
\]
Si las columnas de $A$ son linealmente independientes, el minimizador es único.
\end{teorema}

\begin{proof}
Minimizar $\|Ax - b\|$ sobre $x \in \mathbb{R}^n$ es minimizar $\|y - b\|$ sobre $y \in \operatorname{Col}(A)$, y por el Teorema \ref{teo:mejor_aprox} el mínimo se alcanza exactamente en $y = P_{\operatorname{Col}(A)}(b)$. Por la caracterización de la proyección (Teorema \ref{teo:descomposicion_ortogonal}), $A\hat{x} = P_{\operatorname{Col}(A)}(b)$ si y solo si el residuo $b - A\hat{x}$ pertenece a $\operatorname{Col}(A)^\perp$; como las columnas $a_1, \dots, a_n$ generan $\operatorname{Col}(A)$, eso equivale a las $n$ condiciones $\langle a_i, \, b - A\hat{x} \rangle = 0$. La fila $i$ de $A^T$ es $a_i^T$, de modo que las $n$ condiciones, apiladas, son exactamente
\[
A^T (b - A\hat{x}) = 0 \iff A^T A \, \hat{x} = A^T b.
\]

Si las columnas de $A$ son linealmente independientes, $\ker(A^T A) = \ker(A) = \{0\}$ (como en la demostración del Teorema \ref{teo:svd}), de modo que $A^T A$ es invertible y $\hat{x}$ es único.
\end{proof}

\obs Con la factorización $A = QR$ el problema se resuelve sin formar $A^T A$, mediante una derivación geométrica. Como $\operatorname{Col}(A) = \operatorname{Col}(Q)$, la proyección sobre $\operatorname{Col}(A)$ es $Q Q^T$ (observaciones anteriores), y la condición del minimizador se reescribe
\[
A \hat{x} = P_{\operatorname{Col}(A)}(b)
\iff Q R \, \hat{x} = Q Q^T b
\iff R \, \hat{x} = Q^T b,
\]
donde la última equivalencia se obtiene multiplicando a izquierda por $Q^T$ (con $Q^T Q = I_n$) en una dirección, y por $Q$ en la otra. El sistema $R \hat{x} = Q^T b$ es triangular y se resuelve por sustitución hacia atrás. Los dos factores tienen interpretación directa: $Q^T b$ contiene las coordenadas de la proyección de $b$ en la base ortonormal $q_1, \dots, q_n$, y resolver con $R$ traduce esas coordenadas a coeficientes sobre las columnas originales de $A$. El mismo sistema se obtiene desde las ecuaciones normales: $A^T A = R^T Q^T Q R = R^T R$ y $A^T b = R^T Q^T b$, y cancelando $R^T$, invertible, queda $R \hat{x} = Q^T b$. Evitar el paso por $A^T A$ importa numéricamente, porque ese producto amplifica los errores de redondeo: es la misma razón por la que el PCA se calcula mediante la SVD y no diagonalizando la matriz de covarianzas.

\begin{ejemplo}{Mínimos Cuadrados}{ej_minimos_cuadrados}
Con la matriz $A$ del Ejemplo \ref{ej:ej_qr} y $b = (1, 1, 1)^T$, el sistema $Ax = b$ no tiene solución. Las ecuaciones normales dan
\[
A^T A = \begin{pmatrix} 2 & 1 \\ 1 & 2 \end{pmatrix}, \qquad A^T b = \begin{pmatrix} 2 \\ 2 \end{pmatrix}
\qquad \Longrightarrow \qquad \hat{x} = \Big( \frac{2}{3}, \frac{2}{3} \Big)^T.
\]
Por la vía QR: $Q^T b = \big( \sqrt{2}, \frac{2}{\sqrt{6}} \big)^T$, y la sustitución hacia atrás en $R \hat{x} = Q^T b$ da $\hat{x}_2 = \frac{2/\sqrt{6}}{\sqrt{6}/2} = \frac{2}{3}$ y $\hat{x}_1 = \frac{1}{\sqrt{2}} \big( \sqrt{2} - \frac{1}{\sqrt{2}} \cdot \frac{2}{3} \big) = \frac{2}{3}$: el mismo resultado. El residuo es
\[
b - A\hat{x} = (1,1,1)^T - \Big( \frac{4}{3}, \frac{2}{3}, \frac{2}{3} \Big)^T = \Big( {-\frac{1}{3}}, \frac{1}{3}, \frac{1}{3} \Big)^T,
\]
ortogonal a ambas columnas de $A$ ($-\frac{1}{3} + \frac{1}{3} = 0$ en los dos casos), como exige la proyección.
\end{ejemplo}

\obs El caso de uso clásico es el \emph{ajuste polinomial}: dados $m$ puntos $(t_i, b_i)$ y un grado $d$ con $d + 1 \le m$, encontrar el polinomio $p(t) = x_0 + x_1 t + \dots + x_d t^d$ que minimiza $\sum_i (p(t_i) - b_i)^2$ es exactamente el problema de mínimos cuadrados con la matriz de Vandermonde $a_{ij} = t_i^{\, j-1}$, de tamaño $m \times (d+1)$.

\subsection{El Método de la Potencia}

Las secciones de diagonalización, SVD y PCA suponen que los autovalores y autovectores pueden calcularse, pero no dicen cómo. La vía directa —hallar las raíces del polinomio característico— no es practicable numéricamente: para $n \ge 5$ no existe una fórmula cerrada para las raíces (teorema de Abel-Ruffini, que se cita sin demostración) y, peor, las raíces de un polinomio son extremadamente sensibles a perturbaciones de sus coeficientes. Los métodos reales son iterativos, y el más simple —núcleo conceptual de los demás— es el \textbf{método de la potencia}: multiplicar repetidamente por $A$ y normalizar. Dado un vector inicial $x^{(0)}$ de norma $1$, se itera
\[
x^{(t+1)} = \frac{A \, x^{(t)}}{\| A \, x^{(t)} \|},
\]
y el autovalor se estima con el \textbf{cociente de Rayleigh}
\[
\lambda^{(t)} = \langle x^{(t)}, A \, x^{(t)} \rangle,
\]
el valor de la forma cuadrática de $A$ sobre la esfera unitaria (Teorema \ref{teo:extremos_esfera}).

\begin{teorema}{Convergencia del Método de la Potencia}{potencia}
Sea $A \in \mathbb{R}^{n \times n}$ simétrica con autovalores $\lambda_1 > |\lambda_2| \ge \dots \ge |\lambda_n|$ y una base ortonormal de autovectores $v_1, \dots, v_n$ (Corolario \ref{cor:diag_ortogonal}). Si $c_1 = \langle x^{(0)}, v_1 \rangle \neq 0$, entonces el ángulo $\theta_t$ entre $x^{(t)}$ y $v_1$ cumple
\[
\tan \theta_t \le \Big( \frac{|\lambda_2|}{\lambda_1} \Big)^{\! t} \tan \theta_0,
\]
de modo que $x^{(t)}$ converge a $v_1$ o a $-v_1$ (según el signo de $c_1$) con razón geométrica $|\lambda_2| / \lambda_1$, y $\lambda^{(t)} \to \lambda_1$ con razón $(|\lambda_2| / \lambda_1)^2$.
\end{teorema}

\begin{proof}
Como cada normalización divide por un escalar positivo, $x^{(t)}$ es el vector $y^{(t)} = A^t x^{(0)}$ normalizado. Escríbase $x^{(0)} = \sum_i c_i v_i$ en la base ortonormal de autovectores; como $A^t v_i = \lambda_i^t v_i$,
\[
y^{(t)} = \sum_{i=1}^{n} c_i \, \lambda_i^t \, v_i.
\]
La componente de $y^{(t)}$ sobre $v_1$ es $c_1 \lambda_1^t \neq 0$, y la componente ortogonal a $v_1$ tiene norma $\big( \sum_{i \ge 2} c_i^2 \lambda_i^{2t} \big)^{1/2}$. La tangente del ángulo entre $y^{(t)}$ y $v_1$ es el cociente entre ambas:
\[
\tan \theta_t = \frac{\big( \sum_{i \ge 2} c_i^2 \, \lambda_i^{2t} \big)^{1/2}}{|c_1| \, \lambda_1^t}
\le \Big( \frac{|\lambda_2|}{\lambda_1} \Big)^{\! t} \, \frac{\big( \sum_{i \ge 2} c_i^2 \big)^{1/2}}{|c_1|}
= \Big( \frac{|\lambda_2|}{\lambda_1} \Big)^{\! t} \tan \theta_0,
\]
usando $|\lambda_i|^{2t} \le |\lambda_2|^{2t}$ para $i \ge 2$. Como $|\lambda_2| / \lambda_1 < 1$, el ángulo tiende a $0$; y la componente de $x^{(t)}$ sobre $v_1$, que es $c_1 \lambda_1^t / \|y^{(t)}\|$, conserva el signo de $c_1$, de modo que el límite es $\operatorname{sign}(c_1) \, v_1$.

Para el cociente de Rayleigh, escríbase $x^{(t)} = \sum_i w_i v_i$, con $\sum_i w_i^2 = 1$ y $w_1^2 = \cos^2 \theta_t$. Por ortonormalidad, $\lambda^{(t)} = \sum_i \lambda_i w_i^2$, y entonces
\[
\lambda_1 - \lambda^{(t)} = \sum_{i \ge 2} (\lambda_1 - \lambda_i) \, w_i^2 \le (\lambda_1 - \lambda_n) \sin^2 \theta_t,
\]
que decae con el \emph{cuadrado} de la razón, $(|\lambda_2| / \lambda_1)^{2t}$, pues $\sin \theta_t \le \tan \theta_t$.
\end{proof}

\begin{ejemplo}{El Método de la Potencia en una Matriz $2 \times 2$}{ej_potencia}
Sea $A = \begin{pmatrix} 2 & 1 \\ 1 & 2 \end{pmatrix}$, la matriz $A^T A$ del Ejemplo \ref{ej:ej_minimos_cuadrados}, con autovalores $\lambda_1 = 3$ y $\lambda_2 = 1$ y autovectores $(1, 1)^T$ y $(1, -1)^T$. Con $x^{(0)} = (1, 0)^T = \frac{1}{2} \big[ (1,1)^T + (1,-1)^T \big]$,
\[
A^t x^{(0)} = \frac{1}{2} \Big( 3^t \, (1, 1)^T + (1, -1)^T \Big) = \Big( \frac{3^t + 1}{2}, \ \frac{3^t - 1}{2} \Big)^{T}:
\]
sin normalizar, los iterados son $(1,0)$, $(2,1)$, $(5,4)$, $(14,13)$, $(41,40)$, \dots, cuya dirección tiende a la de $(1,1)$ con razón $|\lambda_2| / \lambda_1 = \frac{1}{3}$, como predice el teorema. Los cocientes de Rayleigh (por ejemplo, para $(2,1)$: $\frac{(2,1) \cdot (5,4)}{2^2 + 1^2} = \frac{14}{5}$) son
\[
\lambda^{(0)} = 2, \qquad \lambda^{(1)} = \frac{14}{5} = 2.8, \qquad \lambda^{(2)} = \frac{122}{41} \approx 2.9756, \qquad \lambda^{(3)} = \frac{1094}{365} \approx 2.9973,
\]
y los errores $1, \ 0.2, \ 0.024, \ 0.0027$ se reducen, en efecto, en un factor cercano a $(|\lambda_2| / \lambda_1)^2 = \frac{1}{9}$ por paso.
\end{ejemplo}

\obs Para las componentes siguientes de una matriz simétrica se usa la \emph{deflación}: hallado el par $(\lambda_1, v_1)$, se pasa a la matriz
\[
A' = A - \lambda_1 \, v_1 v_1^T.
\]
Desarrollando la diagonalización ortogonal $A = P D P^T$ columna a columna se obtiene $A = \sum_i \lambda_i \, v_i v_i^T$, de modo que $A'$ tiene los mismos autovectores que $A$ con el autovalor $\lambda_1$ reemplazado por $0$: el método de la potencia sobre $A'$ produce $(\lambda_2, v_2)$, y así sucesivamente. Aplicado a la matriz de covarianzas, este esquema calcula las primeras $k$ componentes principales del PCA (Sección 9) sin diagonalizar la matriz completa.

\obs Tres notas prácticas. Primera: un $x^{(0)}$ aleatorio cumple $c_1 \neq 0$ con probabilidad $1$. Segunda: cada iteración cuesta un producto matriz-vector, $\Theta(n^2)$, o proporcional a las entradas no nulas si la matriz es dispersa; el ejemplo célebre es PageRank, en esencia un método de la potencia sobre la matriz de enlaces de la web. Tercera: las bibliotecas numéricas (LAPACK) no usan este método para el problema completo, sino el algoritmo QR con desplazamientos, una iteración más elaborada que obtiene todos los autovalores a la vez y que se cita sin demostración; el método de la potencia sigue siendo la elección cuando solo se necesitan las primeras componentes de una matriz enorme y dispersa.

\newpage

\section{Algoritmos de Aprendizaje Automático}

\subsection{Clustering: el Problema k-means}

Dados $n$ puntos de $\mathbb{R}^d$ sin etiquetas, se los quiere agrupar en $k$ grupos de puntos cercanos entre sí. Como medida de cercanía se usa la \textbf{distancia cuadrática}
\[
\delta(x, y) = \|x - y\|^2 = \sum_{a=1}^{d} (x_a - y_a)^2,
\]
el cuadrado de la distancia euclídea. No es una métrica (no cumple la desigualdad triangular), pero es la elección convencional por su conveniencia algebraica, como muestran los teoremas de esta subsección.

\begin{definicion}{El Problema k-means}{kmeans}
Sea $S \subseteq \mathbb{R}^d$ un conjunto de $n$ puntos y sea $k$ un entero positivo. Un \textbf{k-clustering} de $S$ es una partición $\langle S^{(1)}, \dots, S^{(k)} \rangle$ de $S$ junto con una secuencia de \textbf{centros} $C = \langle c^{(1)}, \dots, c^{(k)} \rangle$ en $\mathbb{R}^d$. Su costo es la suma de las distancias cuadráticas de cada punto al centro de su grupo,
\[
f(S, C) = \sum_{\ell=1}^{k} \sum_{x \in S^{(\ell)}} \delta(x, c^{(\ell)}),
\]
y el \textbf{problema k-means} consiste en encontrar el k-clustering de costo mínimo.
\end{definicion}

Los dos teoremas siguientes resuelven por separado las dos mitades del problema: cuál es el mejor centro para un grupo dado, y cuál es el mejor agrupamiento para centros dados.

\begin{teorema}{El Centroide como Centro Óptimo}{centroide}
Sea $S^{(\ell)}$ un cluster no vacío. El único centro $c \in \mathbb{R}^d$ que minimiza $\sum_{x \in S^{(\ell)}} \delta(x, c)$ es su \textbf{centroide} (o media)
\[
c^{(\ell)} = \frac{1}{|S^{(\ell)}|} \sum_{x \in S^{(\ell)}} x.
\]
\end{teorema}

\begin{proof}
Sea $q = |S^{(\ell)}|$. La suma se separa por coordenadas:
\[
\sum_{x \in S^{(\ell)}} \delta(x, c) = \sum_{a=1}^{d} \sum_{x \in S^{(\ell)}} (x_a - c_a)^2
= \sum_{a=1}^{d} \Bigg( \sum_{x \in S^{(\ell)}} x_a^2 - 2 \Big( \sum_{x \in S^{(\ell)}} x_a \Big) c_a + q \, c_a^2 \Bigg).
\]
Para cada $a$, el término entre paréntesis es una cuadrática en $c_a$ con coeficiente principal $q > 0$; completando el cuadrado, con $\mu_a = \frac{1}{q} \sum_{x \in S^{(\ell)}} x_a$,
\[
\sum_{x \in S^{(\ell)}} x_a^2 - 2 q \mu_a c_a + q c_a^2 = q (c_a - \mu_a)^2 + \Big( \sum_{x \in S^{(\ell)}} x_a^2 - q \mu_a^2 \Big),
\]
que se minimiza, de manera única, en $c_a = \mu_a$. Como esto vale para cada coordenada por separado, el único minimizador es $c = (\mu_1, \dots, \mu_d)$, el centroide.
\end{proof}

\begin{teorema}{Asignación al Centro más Cercano}{asignacion_cercana}
Dados $S$ y una secuencia de centros $\langle c^{(1)}, \dots, c^{(k)} \rangle$, un k-clustering minimiza $f(S, C)$ si y solo si asigna cada punto $x \in S$ a un cluster $S^{(\ell)}$ cuyo centro minimiza $\delta(x, c^{(\ell)})$.
\end{teorema}

\begin{proof}
Cada punto $x \in S$ aporta a la suma $f(S, C)$ exactamente un término, $\delta(x, c^{(\ell)})$ para el $\ell$ de su cluster, y los aportes son independientes entre sí: la suma es mínima si y solo si cada aporte lo es, esto es, si cada punto se asigna a un centro más cercano (en distancia cuadrática o euclídea, que inducen el mismo orden).
\end{proof}

El \textbf{procedimiento de Lloyd} alterna las dos optimizaciones parciales que dan los teoremas:
\begin{enumerate}
    \item \textit{Inicializar:} elegir $k$ centros (por ejemplo, $k$ puntos de $S$ al azar).
    \item \textit{Asignar:} formar el clustering asignando cada punto de $S$ a un centro más cercano (Teorema \ref{teo:asignacion_cercana}); los empates se rompen de manera fija, por ejemplo por el índice menor.
    \item \textit{Recalcular:} reemplazar cada centro por el centroide de su cluster (Teorema \ref{teo:centroide}).
    \item Repetir los pasos 2 y 3 hasta que una iteración no cambie el clustering.
\end{enumerate}
Cada iteración cuesta $O(dkn)$ operaciones ($nk$ distancias de $d$ coordenadas, más los centroides); con $T$ iteraciones, el total es $O(T d k n)$.

\begin{proposicion}{Monotonía y Terminación de Lloyd}{lloyd_termina}
En el procedimiento de Lloyd, el costo $f(S, C)$ no aumenta nunca, y el procedimiento termina tras una cantidad finita de iteraciones.
\end{proposicion}

\begin{proof}
La monotonía es consecuencia directa de los dos teoremas: el paso de asignación no aumenta $f$, porque con los centros fijos la asignación al más cercano es óptima entre todas las asignaciones posibles (Teorema \ref{teo:asignacion_cercana}), y el paso de recálculo tampoco, porque con el clustering fijo el centroide es el centro óptimo de cada grupo (Teorema \ref{teo:centroide}).

Para la terminación se examinan los dos casos del paso de recálculo. Si modifica algún centro, entonces $f$ \emph{disminuye estrictamente}: el centroide es el \emph{único} minimizador (Teorema \ref{teo:centroide}), de modo que el centro reemplazado, distinto de él, tenía costo estrictamente mayor. Si no modifica ninguno, los centros quedaron fijos, la siguiente asignación —determinista por la regla fija de desempate— reproduce el clustering anterior, y el procedimiento se detiene.

En consecuencia, toda iteración que no es la última disminuye estrictamente $f$. Además, tras un recálculo el valor de $f$ queda determinado por la partición (los centros son sus centroides), y las particiones de $S$ en a lo sumo $k$ grupos son finitas: los valores que $f$ puede tomar forman un conjunto finito, y una sucesión estrictamente decreciente dentro de un conjunto finito es necesariamente finita. La cantidad de iteraciones es, por lo tanto, finita.
\end{proof}

\obs El procedimiento termina en un mínimo \emph{local}: ninguna reasignación ni recálculo lo mejora. No garantiza el óptimo global —el resultado depende de la inicialización, y en la práctica se lo corre varias veces con inicializaciones distintas—, y encontrar el óptimo global es computacionalmente difícil en general.

\obs El procedimiento de Lloyd ilustra el esquema común a buena parte del aprendizaje automático: definir un espacio de hipótesis parametrizado (aquí, las secuencias de $k$ centros), definir una función que mide qué tan mal ajusta cada hipótesis a los datos (aquí, $f(S, C)$) y minimizarla, al menos localmente, con un procedimiento de optimización. El descenso de gradiente de la subsección siguiente, y el Análisis de Componentes Principales de la Sección 9, responden al mismo esquema.

\begin{ejemplo}{Lloyd en la Recta}{ej_lloyd}
Sea $S = \{0, 2, 5, 6\} \subset \mathbb{R}$ (es decir, $d = 1$) con $k = 2$, y tómense como centros iniciales $c^{(1)} = 0$ y $c^{(2)} = 2$.

\textit{Iteración 1.} Asignación: $0 \to c^{(1)}$; $2, 5, 6 \to c^{(2)}$ (por ejemplo, $\delta(5, 0) = 25 > 9 = \delta(5, 2)$). Recálculo: $c^{(1)} = 0$, $c^{(2)} = \frac{2+5+6}{3} = \frac{13}{3}$. Costo:
\[
f = 0 + \Big( 2 - \tfrac{13}{3} \Big)^2 + \Big( 5 - \tfrac{13}{3} \Big)^2 + \Big( 6 - \tfrac{13}{3} \Big)^2 = \tfrac{49}{9} + \tfrac{4}{9} + \tfrac{25}{9} = \tfrac{26}{3}.
\]

\textit{Iteración 2.} Asignación con $c^{(1)} = 0$, $c^{(2)} = \frac{13}{3}$: el punto $2$ \emph{cambia de cluster}, pues $\delta(2, 0) = 4 < \frac{49}{9} = \delta \big( 2, \frac{13}{3} \big)$; quedan $S^{(1)} = \{0, 2\}$ y $S^{(2)} = \{5, 6\}$. Recálculo: $c^{(1)} = 1$, $c^{(2)} = \frac{11}{2}$. Costo: $f = 1 + 1 + \frac{1}{4} + \frac{1}{4} = \frac{5}{2}$.

\textit{Iteración 3.} La asignación con $c^{(1)} = 1$, $c^{(2)} = \frac{11}{2}$ no cambia nada ($\delta(2,1) = 1 < \frac{49}{4}$, $\delta(5, \frac{11}{2}) = \frac{1}{4} < 16$): el procedimiento termina con costo $\frac{5}{2}$, que en este caso es el óptimo global (las otras particiones en dos intervalos, $\{0\} \cup \{2,5,6\}$ y $\{0,2,5\} \cup \{6\}$, cuestan $\frac{26}{3}$ y $\frac{114}{9}$).
\end{ejemplo}

\obs La compresión de imágenes por \emph{cuantización vectorial} es una aplicación directa: cada píxel es un punto de $\mathbb{R}^3$ (sus intensidades RGB), y correr Lloyd con $k = 256$ produce una paleta de $256$ colores representativos; reemplazando cada píxel por su centro más cercano, la imagen pasa de $24$ a $8$ bits por píxel.

\subsection{Descenso de Gradiente}

La Proposición \ref{prop:prop_max_descenso} identificó la dirección en que una función diferenciable decrece más rápido: $-\nabla f(x)$. El \textbf{descenso de gradiente} la convierte en un algoritmo de minimización: desde un punto inicial $x^{(0)} \in \mathbb{R}^n$, cada iteración se desplaza en esa dirección, con una longitud regulada por un parámetro fijo $\eta > 0$ llamado \textbf{paso} (o \emph{tasa de aprendizaje}):
\[
x^{(t+1)} = x^{(t)} - \eta \, \nabla f(x^{(t)}), \qquad t = 0, 1, \dots, T - 1.
\]
Tras $T$ iteraciones, el algoritmo devuelve el \emph{promedio} $\bar{x} = \frac{1}{T} \big( x^{(0)} + \dots + x^{(T-1)} \big)$; se devuelve el promedio, y no el último iterado, porque es la variante para la cual la demostración de convergencia que sigue obtiene una garantía. Sin hipótesis sobre $f$, un método local de este tipo solo puede aspirar a un mínimo local; la hipótesis que convierte lo local en global es la convexidad:

\begin{definicion}{Función Convexa}{convexidad}
Una función $f: \mathbb{R}^n \to \mathbb{R}$ es \textbf{convexa} si para todos $x, y \in \mathbb{R}^n$ y todo $\lambda \in [0, 1]$,
\[
f(\lambda x + (1 - \lambda) y) \le \lambda f(x) + (1 - \lambda) f(y).
\]
\end{definicion}

\obs En una función convexa, todo mínimo local es global: si $x$ es mínimo local y existiera $y$ con $f(y) < f(x)$, los puntos $\lambda y + (1-\lambda) x$ con $\lambda \to 0^+$ estarían arbitrariamente cerca de $x$ con valor a lo sumo $\lambda f(y) + (1-\lambda) f(x) < f(x)$, contra la minimalidad local.

\begin{lema}\label{lem_tangente}
Sea $f: \mathbb{R}^n \to \mathbb{R}$ convexa y diferenciable. Para todos $x, y \in \mathbb{R}^n$,
\[
f(y) \ge f(x) + \langle \nabla f(x), y - x \rangle,
\]
es decir, la función queda por encima de sus planos tangentes.
\end{lema}

\begin{proof}
Por convexidad, para $\lambda \in (0, 1]$,
\[
f(x + \lambda (y - x)) = f((1-\lambda) x + \lambda y) \le (1 - \lambda) f(x) + \lambda f(y),
\]
de donde
\[
\frac{f(x + \lambda (y - x)) - f(x)}{\lambda} \le f(y) - f(x).
\]
Cuando $\lambda \to 0^+$, el lado izquierdo tiende a la derivada direccional de $f$ en $x$ según $y - x$, que por diferenciabilidad es $\langle \nabla f(x), y - x \rangle$ (Definición \ref{def:def_diferencial} y Teorema \ref{teo:teo_gradiente_riesz}). Reordenando queda la desigualdad.
\end{proof}

\begin{lema}\label{lem_jensen}
Sea $f: \mathbb{R}^n \to \mathbb{R}$ convexa. Para todos $x^{(0)}, \dots, x^{(T-1)} \in \mathbb{R}^n$,
\[
f \Bigg( \frac{x^{(0)} + \dots + x^{(T-1)}}{T} \Bigg) \le \frac{f(x^{(0)}) + \dots + f(x^{(T-1)})}{T}.
\]
\end{lema}

\begin{proof}
Por inducción en $T$; el caso $T = 1$ es una igualdad. Para el paso, el promedio de $T$ puntos se escribe como combinación convexa del último punto y el promedio de los $T - 1$ primeros:
\[
\frac{x^{(0)} + \dots + x^{(T-1)}}{T} = \frac{1}{T} \, x^{(T-1)} + \frac{T-1}{T} \cdot \frac{x^{(0)} + \dots + x^{(T-2)}}{T-1},
\]
y aplicando la Definición \ref{def:convexidad} con $\lambda = \frac{1}{T}$ y luego la hipótesis inductiva se obtiene la cota.
\end{proof}

\begin{teorema}{Convergencia del Descenso de Gradiente}{convergencia_gd}
Sea $f: \mathbb{R}^n \to \mathbb{R}$ convexa y diferenciable con minimizador $x^*$. Sean $R = \|x^{(0)} - x^*\|$ y $L$ una cota superior de $\|\nabla f(x^{(t)})\|$ sobre los iterados $t = 0, \dots, T-1$. Si el paso es $\eta = \frac{R}{L \sqrt{T}}$, entonces el promedio $\bar{x}$ que devuelve el descenso de gradiente cumple
\[
f(\bar{x}) - f(x^*) \le \frac{R L}{\sqrt{T}}.
\]
\end{teorema}

\begin{proof}
Se estudia cómo evoluciona la distancia al minimizador. Expandiendo la norma cuadrada de $x^{(t+1)} - x^* = (x^{(t)} - x^*) - \eta \nabla f(x^{(t)})$,
\[
\|x^{(t+1)} - x^*\|^2 = \|x^{(t)} - x^*\|^2 - 2 \eta \, \langle \nabla f(x^{(t)}), \, x^{(t)} - x^* \rangle + \eta^2 \, \|\nabla f(x^{(t)})\|^2.
\]
Despejando de esa igualdad el producto interno,
\[
\langle \nabla f(x^{(t)}), \, x^{(t)} - x^* \rangle
= \frac{\|x^{(t)} - x^*\|^2 - \|x^{(t+1)} - x^*\|^2}{2 \eta} + \frac{\eta}{2} \, \|\nabla f(x^{(t)})\|^2.
\]
Por otro lado, el Lema \ref{lem_tangente} con $x = x^{(t)}$ e $y = x^*$ da $f(x^*) \ge f(x^{(t)}) + \langle \nabla f(x^{(t)}), x^* - x^{(t)} \rangle$, es decir, $f(x^{(t)}) - f(x^*) \le \langle \nabla f(x^{(t)}), x^{(t)} - x^* \rangle$. Combinando ambas y acotando $\|\nabla f(x^{(t)})\| \le L$,
\[
f(x^{(t)}) - f(x^*) \le \frac{\|x^{(t)} - x^*\|^2 - \|x^{(t+1)} - x^*\|^2}{2 \eta} + \frac{\eta L^2}{2}.
\]
Al sumar sobre $t = 0, \dots, T-1$, las distancias de la primera fracción se cancelan en cadena (telescopía), queda solo la primera menos la última, y el término $-\|x^{(T)} - x^*\|^2 \le 0$ se descarta:
\[
\sum_{t=0}^{T-1} \big( f(x^{(t)}) - f(x^*) \big) \le \frac{\|x^{(0)} - x^*\|^2}{2 \eta} + \frac{T \eta L^2}{2} = \frac{R^2}{2 \eta} + \frac{T \eta L^2}{2}.
\]
Dividiendo por $T$ y aplicando el Lema \ref{lem_jensen} al lado izquierdo,
\[
f(\bar{x}) - f(x^*) \le \frac{1}{T} \sum_{t=0}^{T-1} \big( f(x^{(t)}) - f(x^*) \big) \le \frac{R^2}{2 \eta T} + \frac{\eta L^2}{2},
\]
y con $\eta = \frac{R}{L \sqrt{T}}$ los dos sumandos valen $\frac{RL}{2\sqrt{T}}$ cada uno.
\end{proof}

\obs Tres observaciones sobre la cota. Primera: para lograr error $\varepsilon$ bastan $T = (RL/\varepsilon)^2$ iteraciones, cada una con el costo de evaluar un gradiente: el error decae como $O(1/\sqrt{T})$. Segunda: el paso óptimo equilibra los dos sumandos de la demostración; un paso demasiado grande hace dominar el término $\eta L^2 / 2$ (el algoritmo sobrepasa el mínimo en cada iteración), y uno demasiado pequeño, el término $R^2 / (2\eta T)$ (el avance por iteración resulta insuficiente). Tercera: la demostración no necesita que $f$ decrezca en cada paso; con hipótesis adicionales (gradiente Lipschitz, convexidad fuerte) se obtienen tasas mejores, $O(1/T)$ y geométricas, que se citan sin demostración.

\begin{ejemplo}{Descenso de Gradiente y Mínimos Cuadrados}{ej_gd}
La función de error cuadrático $f(x) = \frac{1}{2} \|Ax - b\|^2$ es convexa: dados $x, y$ y $\lambda \in [0,1]$, como $A(\lambda x + (1-\lambda) y) - b = \lambda (Ax - b) + (1-\lambda)(Ay - b)$, la desigualdad triangular y la convexidad de $t \mapsto t^2$ en $[0, \infty)$ dan
\[
f(\lambda x + (1-\lambda) y) \le \tfrac{1}{2} \big( \lambda \|Ax - b\| + (1-\lambda) \|Ay - b\| \big)^2 \le \lambda f(x) + (1-\lambda) f(y).
\]
Su gradiente es $\nabla f(x) = A^T A x - A^T b$ (como en la sección de formas cuadráticas, usando que $A^T A$ es simétrica), de modo que $\nabla f(x) = 0$ son exactamente las ecuaciones normales (Teorema \ref{teo:ecuaciones_normales}): el descenso de gradiente converge a la solución de mínimos cuadrados, y es la alternativa práctica a QR cuando $A$ es demasiado grande para factorizarla.

Con la matriz $A$ y el vector $b$ del Ejemplo \ref{ej:ej_minimos_cuadrados} ($\hat{x} = (\frac{2}{3}, \frac{2}{3})^T$), partiendo de $x^{(0)} = (0,0)^T$ con paso $\eta = \frac{1}{4}$:
\begin{gather*}
\nabla f(x) = \begin{pmatrix} 2 & 1 \\ 1 & 2 \end{pmatrix} x - \begin{pmatrix} 2 \\ 2 \end{pmatrix};
\\
x^{(1)} = \tfrac{1}{4} \begin{pmatrix} 2 \\ 2 \end{pmatrix} = \begin{pmatrix} 1/2 \\ 1/2 \end{pmatrix}, \qquad
x^{(2)} = \begin{pmatrix} 5/8 \\ 5/8 \end{pmatrix}, \qquad
x^{(3)} = \begin{pmatrix} 21/32 \\ 21/32 \end{pmatrix}, \; \dots
\end{gather*}
(por ejemplo, $\nabla f(x^{(1)}) = (\frac{3}{2}, \frac{3}{2})^T - (2,2)^T = -(\frac{1}{2}, \frac{1}{2})^T$ y $x^{(2)} = x^{(1)} + \frac{1}{4} (\frac{1}{2}, \frac{1}{2})^T$). Los valores de $f$ decrecen hacia $f(\hat{x}) = \frac{1}{6}$:
\[
f(x^{(0)}) = \tfrac{3}{2}, \qquad f(x^{(1)}) = \tfrac{1}{4}, \qquad f(x^{(2)}) = \tfrac{11}{64}, \qquad f(x^{(3)}) = \tfrac{171}{1024} \approx 0.167.
\]
La convergencia aquí es mucho más rápida que la garantía general del Teorema \ref{teo:convergencia_gd}: los iterados viven en la recta de dirección $(1,1)^T$, autovector de $A^T A$ con autovalor $3$, y el error se contrae en el factor $|1 - \eta \cdot 3| = \frac{1}{4}$ por paso: $x^{(t)} = \frac{2}{3} - \frac{2}{3} \cdot 4^{-t}$ en cada coordenada. Es un anticipo del papel de los autovalores en la velocidad del descenso de gradiente.
\end{ejemplo}

\obs En el aprendizaje automático moderno, $f$ suele ser el error de un modelo sobre millones de datos y $x$ sus parámetros; evaluar el gradiente completo es costoso, y se lo reemplaza por el gradiente de una muestra aleatoria de datos (\emph{descenso de gradiente estocástico}, SGD). Además, las funciones de error de los modelos profundos no son convexas, de modo que las garantías de esta sección no aplican literalmente; aun así, el descenso de gradiente y sus variantes constituyen el procedimiento de entrenamiento estándar.

\end{document}
